
\Data\ID{1003021701}
\Updated{2022-08-25 10:08:24}
\Level{A}
\SubArea{Triangles}
\LANGen{
\Question{
Interior angles of a triangle \( ABC \) are in the ratio \( \alpha:\beta:\gamma=2:4:6 \). Calculate the measures of these angles.}
{\( \alpha=30^{\circ};\ \beta=60^{\circ};\ \gamma=90^{\circ} \)}{\( \alpha=20^{\circ};\ \beta=40^{\circ};\ \gamma=60^{\circ} \)}{\( \alpha=15^{\circ};\ \beta=30^{\circ};\ \gamma=135^{\circ} \)}{\( \alpha=90^{\circ};\ \beta=60^{\circ};\ \gamma=30^{\circ} \)}{}{}}
\LANGpl{
\Question{
Stosunek wewnętrznych kątów trójkąta  \( ABC \) wynosi \( \alpha:\beta:\gamma=2:4:6 \). Oblicz miary tych kątów.}
{\( \alpha=30^{\circ};\ \beta=60^{\circ};\ \gamma=90^{\circ} \)}{\( \alpha=20^{\circ};\ \beta=40^{\circ};\ \gamma=60^{\circ} \)}{\( \alpha=15^{\circ};\ \beta=30^{\circ};\ \gamma=135^{\circ} \)}{\( \alpha=90^{\circ};\ \beta=60^{\circ};\ \gamma=30^{\circ} \)}{}{}}
\LANGcs{
\Question{
V trojúhelníku \( ABC \) jsou vnitřní úhly v poměru \( \alpha:\beta:\gamma=2:4:6 \). Vypočítejte velikost těchto vnitřních úhlů.}
{\( \alpha=30^{\circ};\ \beta=60^{\circ};\ \gamma=90^{\circ} \)}{\( \alpha=20^{\circ};\ \beta=40^{\circ};\ \gamma=60^{\circ} \)}{\( \alpha=15^{\circ};\ \beta=30^{\circ};\ \gamma=135^{\circ} \)}{\( \alpha=90^{\circ};\ \beta=60^{\circ};\ \gamma=30^{\circ} \)}{}{}}
\LANGsk{
\Question{
V trojuholníku  \( ABC \) sú vnútorné uhly v pomere \( \alpha:\beta:\gamma=2:4:6 \). Vypočítajte veľkosti týchto vnútorných uhlov.}
{\( \alpha=30^{\circ};\ \beta=60^{\circ};\ \gamma=90^{\circ} \)}{\( \alpha=20^{\circ};\ \beta=40^{\circ};\ \gamma=60^{\circ} \)}{\( \alpha=15^{\circ};\ \beta=30^{\circ};\ \gamma=135^{\circ} \)}{\( \alpha=90^{\circ};\ \beta=60^{\circ};\ \gamma=30^{\circ} \)}{}{}}
\LANGes{
\Question{
Los ángulos interiores de un triángulo \( ABC \) están en proporción \( \alpha:\beta:\gamma=2:4:6 \). Calcula las medidas de los ángulos.}
{\( \alpha=30^{\circ};\ \beta=60^{\circ};\ \gamma=90^{\circ} \)}{\( \alpha=20^{\circ};\ \beta=40^{\circ};\ \gamma=60^{\circ} \)}{\( \alpha=15^{\circ};\ \beta=30^{\circ};\ \gamma=135^{\circ} \)}{\( \alpha=90^{\circ};\ \beta=60^{\circ};\ \gamma=30^{\circ} \)}{}{}}
\End

\Data\ID{1003021703}
\Updated{2022-08-25 10:08:24}
\Level{A}
\SubArea{Triangles}
\LANGen{
\Question{
The measure of an  exterior angle of an isosceles triangle is \( 84^{\circ} \). Calculate the measures of all interior angles of the triangle.}
{\( 96^{\circ};\ 42^{\circ};\ 42^{\circ} \)}{\( 84^{\circ};\ 48^{\circ};\ 48^{\circ} \)}{\( 12^{\circ};\ 84^{\circ};\ 84^{\circ} \)}{\( 96^{\circ};\ 96^{\circ};\ 12^{\circ} \)}{}{}}
\LANGpl{
\Question{
Miara zewnętrznego kąta trójkąta równoramiennego wynosi \( 84^{\circ} \). Oblicz miarę wszystkich jego kątów wewnętrznych.}
{\( 96^{\circ};\ 42^{\circ};\ 42^{\circ} \)}{\( 84^{\circ};\ 48^{\circ};\ 48^{\circ} \)}{\( 12^{\circ};\ 84^{\circ};\ 84^{\circ} \)}{\( 96^{\circ};\ 96^{\circ};\ 12^{\circ} \)}{}{}}
\LANGcs{
\Question{
Vnější úhel rovnoramenného trojúhelníku má velikost \( 84^{\circ} \). Vypočítejte velikost všech vnitřních úhlů trojúhelníku.}
{\( 96^{\circ};\ 42^{\circ};\ 42^{\circ} \)}{\( 84^{\circ};\ 48^{\circ};\ 48^{\circ} \)}{\( 12^{\circ};\ 84^{\circ};\ 84^{\circ} \)}{\( 96^{\circ};\ 96^{\circ};\ 12^{\circ} \)}{}{}}
\LANGsk{
\Question{
Vonkajší uhol rovnoramenného trojuholníka má veľkosť  \( 84^{\circ} \). Vypočítajte veľkosti všetkých vnútorných uhlov trojuholníka.}
{\( 96^{\circ};\ 42^{\circ};\ 42^{\circ} \)}{\( 84^{\circ};\ 48^{\circ};\ 48^{\circ} \)}{\( 12^{\circ};\ 84^{\circ};\ 84^{\circ} \)}{\( 96^{\circ};\ 96^{\circ};\ 12^{\circ} \)}{}{}}
\LANGes{
\Question{
La medida de un ángulo exterior de un triángulo isósceles es \( 84^{\circ} \). Calcula las medidas de todos los ángulos interiores del triángulo.}
{\( 96^{\circ};\ 42^{\circ};\ 42^{\circ} \)}{\( 84^{\circ};\ 48^{\circ};\ 48^{\circ} \)}{\( 12^{\circ};\ 84^{\circ};\ 84^{\circ} \)}{\( 96^{\circ};\ 96^{\circ};\ 12^{\circ} \)}{}{}}
\End

\Data\ID{1003021705}
\Updated{2022-08-25 10:08:24}
\Level{A}
\SubArea{Triangles}
\LANGen{
\Question{
Calculate the measures of interior angles \( \alpha \), \( \beta \) and \( \gamma \) of  a triangle if \( \gamma=2\beta \) and \( \beta=3\alpha \).}
{\( \alpha=18^{\circ};\ \beta=54^{\circ};\ \gamma=108^{\circ} \)}{\( \alpha=15^{\circ};\ \beta=45^{\circ};\ \gamma=90^{\circ} \)}{\( \alpha=12^{\circ};\ \beta=54^{\circ};\ \gamma=111^{\circ} \)}{\( \alpha=54^{\circ};\ \beta=18^{\circ};\ \gamma=108^{\circ} \)}{}{}}
\LANGpl{
\Question{
Oblicz miary kątów wewnętrznych  \( \alpha \), \( \beta \) i \( \gamma \) trójkąta, jeśli \( \gamma=2\beta \) i \( \beta=3\alpha \).}
{\( \alpha=18^{\circ};\ \beta=54^{\circ};\ \gamma=108^{\circ} \)}{\( \alpha=15^{\circ};\ \beta=45^{\circ};\ \gamma=90^{\circ} \)}{\( \alpha=12^{\circ};\ \beta=54^{\circ};\ \gamma=111^{\circ} \)}{\( \alpha=54^{\circ};\ \beta=18^{\circ};\ \gamma=108^{\circ} \)}{}{}}
\LANGcs{
\Question{
Vypočítejte velikosti vnitřních úhlů \( \alpha \), \( \beta \) a \( \gamma \) trojúhelníku,  jestliže platí \( \gamma=2\beta \) a \( \beta=3\alpha \).}
{\( \alpha=18^{\circ};\ \beta=54^{\circ};\ \gamma=108^{\circ} \)}{\( \alpha=15^{\circ};\ \beta=45^{\circ};\ \gamma=90^{\circ} \)}{\( \alpha=12^{\circ};\ \beta=54^{\circ};\ \gamma=111^{\circ} \)}{\( \alpha=54^{\circ};\ \beta=18^{\circ};\ \gamma=108^{\circ} \)}{}{}}
\LANGsk{
\Question{
Vypočítajte veľkosti vnútorných uhlov \( \alpha \), \( \beta \) a \( \gamma \) trojuholníka, ak platí \( \gamma=2\beta \) a \( \beta=3\alpha \).}
{\( \alpha=18^{\circ};\ \beta=54^{\circ};\ \gamma=108^{\circ} \)}{\( \alpha=15^{\circ};\ \beta=45^{\circ};\ \gamma=90^{\circ} \)}{\( \alpha=12^{\circ};\ \beta=54^{\circ};\ \gamma=111^{\circ} \)}{\( \alpha=54^{\circ};\ \beta=18^{\circ};\ \gamma=108^{\circ} \)}{}{}}
\LANGes{
\Question{
Calcula las medidas de los ángulos interiores \( \alpha \), \( \beta \) y \( \gamma \) de un triángulo si \( \gamma=2\beta \) y \( \beta=3\alpha \).}
{\( \alpha=18^{\circ};\ \beta=54^{\circ};\ \gamma=108^{\circ} \)}{\( \alpha=15^{\circ};\ \beta=45^{\circ};\ \gamma=90^{\circ} \)}{\( \alpha=12^{\circ};\ \beta=54^{\circ};\ \gamma=111^{\circ} \)}{\( \alpha=54^{\circ};\ \beta=18^{\circ};\ \gamma=108^{\circ} \)}{}{}}
\End

\Data\ID{1003021707}
\Updated{2022-08-25 10:08:24}
\Level{A}
\SubArea{Triangles}
\LANGen{
\Question{
Choose the false statement.}
{All altitudes in a right-angled triangle are perpendicular to one another.}{The centroid of a triangle divides each median in the ratio \( 2:1 \).}{The midline of a triangle is parallel to the third side and half as long.}{The medians of a triangle intersect in a single point, the centroid of a triangle.}{}{}}
\LANGpl{
\Question{
Wybierz fałszywe stwierdzenie.}
{Wszystkie wysokości w trójkącie prostokątnym są prostopadłe do siebie.}{Środek ciężkości trójkąta dzieli każdą środkową w stosunku \( 2:1 \).}{Linia środkowa trójkąta jest równoległa do trzeciego boku i równa jego połowie.}{Środkowe trójkąta przecinają się w jednym punkcie, środku ciężkości trójkąta.}{}{}}
\LANGcs{
\Question{
Vyberte nesprávné tvrzení.}
{Všechny výšky v pravoúhlém trojúhelníku jsou na sebe navzájem kolmé.}{Těžiště trojúhelníku dělí každou těžnici v poměru \( 2:1 \).}{Střední příčka v trojúhelníku má stejnou délku jako polovina strany, se kterou je rovnoběžná.}{Těžnice trojúhelníku se protínají v jednom bodě, který nazýváme těžiště trojúhelníku.}{}{}}
\LANGsk{
\Question{
Vyberte nesprávne tvrdenie.}
{Všetky výšky v pravouhlom trojuholníku sú na seba navzájom kolmé.}{Ťažisko trojuholníka delí každú ťažnicu v pomere  \( 2:1 \).}{Stredná priečka v trojuholníku má dĺžku rovnú polovici strany, s ktorou je rovnobežná.}{Ťažnice trojuholníka sa pretínajú v jednom bode, ktorý nazývame ťažiskom trojuholníka.}{}{}}
\LANGes{
\Question{
Elige la proposición falsa.}
{En un triángulo rectángulo, todas las alturas son perpendiculares entre sí.}{El baricentro de un triángulo divide a cada mediana a razón \( 2:1 \).}{La paralela media de un triángulo tiene la misma longitud que la mitad del lado con el que es paralela.}{Las tres medianas de un triángulo se cortan en un punto llamado baricentro.}{}{}}
\End

\Data\ID{1003021710}
\Updated{2022-08-25 10:08:24}
\Level{A}
\SubArea{Triangles}
\LANGen{
\Question{
In a right-angled triangle \( ABC \) with a right angle at the vertex \( A \) the measure of the angle \( ABC \) is \( 50^{\circ} \). Give the measure of  the angle \( BCA \).}
{\( 40^{\circ} \)}{\( 90^{\circ} \)}{\( 30^{\circ} \)}{\( 180^{\circ} \)}{}{}}
\LANGpl{
\Question{
W trójkącie prostokątnym  \( ABC \) z kątem prostym w wierzchołku \( A \) miara kąta \( ABC \) wynosi \( 50^{\circ} \). Oblicz miarę kąta \( BCA \).}
{\( 40^{\circ} \)}{\( 90^{\circ} \)}{\( 30^{\circ} \)}{\( 180^{\circ} \)}{}{}}
\LANGcs{
\Question{
V pravoúhlém trojúhelníku \( ABC \) s pravým úhlem u vrcholu \( A \) má úhel \( ABC \) velikost \( 50^{\circ} \). Vypočítejte velikost úhlu \( BCA \).}
{\( 40^{\circ} \)}{\( 90^{\circ} \)}{\( 30^{\circ} \)}{\( 180^{\circ} \)}{}{}}
\LANGsk{
\Question{
V pravouhlom trojuholníku  \( ABC \) s pravým uhlom pri vrchole \( A \) má uhol  \( ABC \) veľkosť \( 50^{\circ} \). Vypočítajte veľkosť uhla  \( BCA \).}
{\( 40^{\circ} \)}{\( 90^{\circ} \)}{\( 30^{\circ} \)}{\( 180^{\circ} \)}{}{}}
\LANGes{
\Question{
En un triángulo rectángulo \( ABC \), siendo \( A \) el vértice del ángulo recto, la medida del ángulo \( ABC \) es \( 50^{\circ} \). Calcula la medida del ángulo \( BCA \).}
{\( 40^{\circ} \)}{\( 90^{\circ} \)}{\( 30^{\circ} \)}{\( 180^{\circ} \)}{}{}}
\End

\Data\ID{1003021711}
\Updated{2022-08-25 10:08:24}
\Level{A}
\SubArea{Triangles}
\LANGen{
\Question{
What is the measure of all interior angles in an equilateral triangle?}
{\( 60^{\circ} \)}{\( 30^{\circ} \)}{\( 90^{\circ} \)}{\( 45^{\circ} \)}{}{}}
\LANGpl{
\Question{
Jaka jest miara wszystkich kątów wewnętrznych w trójkącie równobocznym?}
{\( 60^{\circ} \)}{\( 30^{\circ} \)}{\( 90^{\circ} \)}{\( 45^{\circ} \)}{}{}}
\LANGcs{
\Question{
Jakou velikost mají všechny vnitřní úhly v rovnostranném trojúhelníku?}
{\( 60^{\circ} \)}{\( 30^{\circ} \)}{\( 90^{\circ} \)}{\( 45^{\circ} \)}{}{}}
\LANGsk{
\Question{
Akú veľkosť majú všetky vnútorné uhly v rovnostrannom trojuholníku?}
{\( 60^{\circ} \)}{\( 30^{\circ} \)}{\( 90^{\circ} \)}{\( 45^{\circ} \)}{}{}}
\LANGes{
\Question{
¿Cuánto miden los ángulos interiores de un ángulo equilátero?}
{\( 60^{\circ} \)}{\( 30^{\circ} \)}{\( 90^{\circ} \)}{\( 45^{\circ} \)}{}{}}
\End

\Data\ID{1003076801}
\Updated{2020-07-15 03:07:12}
\Level{A}
\SubArea{Triangles}
\LANGen{
\Question{
\( ABC \) is a triangle with sides \( a \), \( b \), \( c \). Let \( a\leq b\leq c \). Two of its interior angles have measures of \( 70^{\circ} \) and \( 50^{\circ} \). Which of the following statements about the triangle \( ABC \) is true?}
{The third interior angle is opposite the side \( b \).}{The angle of the measure \( 70^{\circ} \) lies opposite the side \( a \).}{The angle of the measure \( 50^{\circ} \) lies opposite the side \( b \).}{The third interior angle is opposite the side \( c \).}{}{}}
\LANGpl{
\Question{
\( ABC \) jest trójkątem o bokach \( a \), \( b \), \( c \). Niech \( a\leq b\leq c \). Dwa z jego wewnętrznych kątów mają miarę \( 70^{\circ} \) i \( 50^{\circ} \). Które z poniższych stwierdzeń dotyczących trójkąta \( ABC \) jest prawdziwe?}
{Trzeci wewnętrzny kąt jest naprzeciwko boku  \( b \).}{Kąt o mierze \( 70^{\circ} \) znajduje się naprzeciwko boku \( a \).}{Kąt o mierze \( 50^{\circ} \) leży naprzeciw boku  \( b \).}{Trzeci kąt wewnętrzny znajduje się naprzeciw boku \( c \).}{}{}}
\LANGcs{
\Question{
V trojúhelníku \( ABC \) pro velikost stran \( a \), \( b \), \( c \) platí \( a\leq b\leq c \). Dva z jeho vnitřních úhlů mají velikost \( 70^{\circ} \) a \( 50^{\circ} \). Které z následujících tvrzení o trojúhelníku \( ABC \) je pravdivé?}
{Proti straně \( b \) leží třetí vnitřní úhel.}{Úhel o velikosti \( 70^{\circ} \) leží proti straně \( a \).}{Úhel o velikosti \( 50^{\circ} \) leží proti straně \( b \).}{Třetí vnitřní úhel leží proti straně \( c \).}{}{}}
\LANGsk{
\Question{
V trojuholníku  \( ABC \) pre veľkosti strán  \( a \), \( b \), \( c \) platí \( a\leq b\leq c \). Dva z jeho vnútorných uhlov majú veľkosť \( 70^{\circ} \) a \( 50^{\circ} \). Ktoré z nasledujúcich tvrdení o trojuholníku \( ABC \) je pravdivé?}
{Oproti strane  \( b \) leží tretí vnútorný uhol.}{Uhol s veľkosťou  \( 70^{\circ} \) leží oproti strane \( a \).}{Uhol s veľkosťou  \( 50^{\circ} \) leží oproti strane \( b \).}{Tretí vnútorný uhol leží oproti strane  \( c \).}{}{}}
\LANGes{
\Question{
Dado el triángulo \( ABC \) con los lados \( a \), \( b \), \( c \), suponiendo que \( a\leq b\leq c \). Dos de sus ángulos interiores miden \( 70^{\circ} \) y \( 50^{\circ} \). Identifica la proposición lógica sobre el triángulo \( ABC \).}
{El tercer ángulo interior es opuesto al lado \( b \).}{El ángulo de la medida \( 70^{\circ} \) es opuesto al lado \( a \).}{El ángulo de la medida \( 50^{\circ} \) es opuesto al lado \( b \).}{El tercer ángulo interior es opuesto al lado \( c \).}{}{}}
\End

\Data\ID{1003076802}
\Updated{2020-07-15 03:07:12}
\Level{A}
\SubArea{Triangles}
\LANGen{
\Question{
If the orthocenter of a triangle lies outside the triangle, then it is an/a:}
{obtuse triangle}{right-angled triangle}{acute triangle}{equilateral triangle}{}{}}
\LANGpl{
\Question{
Jeżeli punkt przecięcia wszystkich wysokości trójkąta znajduje się poza nim, jest to trójkąt:}
{rozwartokątny}{prostokątny}{ostrokątny}{równoboczny}{}{}}
\LANGcs{
\Question{
Jestliže průsečík výšek trojúhelníku leží mimo trojúhelník, tak jde o:}
{tupoúhlý trojúhelník}{pravoúhlý trojúhelník}{ostroúhlý trojúhelník}{rovnostranný trojúhelník}{}{}}
\LANGsk{
\Question{
Ak priesečník výšok trojuholníka leží mimo trojuholníka, tak ide o:}
{tupouhlý trojuholník}{pravouhlý trojuholník}{ostrouhlý trojuholník}{rovnostranný trojuholník}{}{}}
\LANGes{
\Question{
Si el ortocentro de las alturas de un triángulo está fuera del triángulo, se trata de un:}
{triángulo obtusángulo}{triángulo rectángulo}{triángulo acutángulo}{triángulo equilátero}{}{}}
\End

\Data\ID{1003076803}
\Updated{2021-04-19 06:04:07}
\Level{A}
\SubArea{Triangles}
\LANGen{
\Question{
Choose the false statement:}
{There exists a right-angled equilateral triangle.}{There exists an acute isosceles triangle.}{There exists an acute equilateral triangle.}{There exists an obtuse isosceles triangle.}{}{}}
\LANGpl{
\Question{
Wybierz fałszywe stwierdzenie.}
{Istnieje trójkąt prostokątny, który jest równoboczny.}{Istnieje trójkąt ostrokątny, który jest równoramienny.}{Istnieje trójkąt ostrokątny, który jest równoboczny.}{Istnieje trójkąt równoramienny, który jest rozwartokątny.}{}{}}
\LANGcs{
\Question{
Vyberte nepravdivé tvrzení:}
{Existuje pravoúhlý rovnostranný trojúhelník.}{Existuje ostroúhlý rovnoramenný trojúhelník.}{Existuje ostroúhlý rovnostranný trojúhelník.}{Existuje tupoúhlý rovnoramenný trojúhelník.}{}{}}
\LANGsk{
\Question{
Vyberte nepravdivé tvrdenie:}
{Existuje pravouhlý rovnostranný trojuholník.}{Existuje ostrouhlý rovnoramenný trojuholník.}{Existuje ostrouhlý rovnostranný trojuholník.}{Existuje tupouhlý rovnoramenný trojuholník.}{}{}}
\LANGes{
\Question{
Identifica la proposición falsa:}
{Hay triángulos equiláteros que son rectángulos.}{Hay triángulos isósceles que son acutángulos.}{Hay triángulos equiláteros que son acutángulos.}{Hay triángulos isósceles que son obtusángulos.}{}{}}
\End

\Data\ID{1003076804}
\Updated{2022-04-23 05:04:05}
\Level{A}
\SubArea{Triangles}
\LANGen{
\Question{
The measures of two interior angles of a triangle are \( 31^{\circ}20' \) and \( 25^{\circ} \). What is the measure of the third interior angle?}
{\( 123^{\circ}40' \)}{\( 56^{\circ}20' \)}{\( 5^{\circ}40' \)}{\( 124^{\circ}40' \)}{}{}}
\LANGpl{
\Question{
Miary dwóch kątów wewnętrznych trójkąta są równe \( 31^{\circ}20' \) i \( 25^{\circ} \). Jaka jest miara trzeciego kąta wewnętrznego tego trójkąta?}
{\( 123^{\circ}40' \)}{\( 56^{\circ}20' \)}{\( 5^{\circ}40' \)}{\( 124^{\circ}40' \)}{}{}}
\LANGcs{
\Question{
Velikosti dvou vnitřních úhlů trojúhelníku jsou \( 31^{\circ}20' \) a \( 25^{\circ} \). Jakou velikost má třetí vnitřní úhel?}
{\( 123^{\circ}40' \)}{\( 56^{\circ}20' \)}{\( 5^{\circ}40' \)}{\( 124^{\circ}40' \)}{}{}}
\LANGsk{
\Question{
Veľkosti dvoch vnútorných uhlov trojuholníka sú \( 31^{\circ}20' \) a \( 25^{\circ} \). Akú veľkosť má tretí vnútorný uhol?}
{\( 123^{\circ}40' \)}{\( 56^{\circ}20' \)}{\( 5^{\circ}40' \)}{\( 124^{\circ}40' \)}{}{}}
\LANGes{
\Question{
Las medidas de dos ángulos interiores de un triángulo son \( 31^{\circ}20' \) y \( 25^{\circ} \). ¿Cuánto mide el tercer ángulo interior?}
{\( 123^{\circ}40' \)}{\( 56^{\circ}20' \)}{\( 5^{\circ}40' \)}{\( 124^{\circ}40' \)}{}{}}
\End

\Data\ID{1003076805}
\Updated{2020-07-15 03:07:12}
\Level{A}
\SubArea{Triangles}
\LANGen{
\Question{
If two interior angles of a triangle \( ABC \) have measures of \( 60^{\circ} \), then the triangle \( ABC \) is an/a:}
{equilateral triangle}{isosceles triangle}{right-angled triangle}{obtuse triangle}{}{}}
\LANGpl{
\Question{
Jeżeli dwa wewnętrzne kąty trójkąta \( ABC \) mają miarę  \( 60^{\circ} \), wtedy trójkąt \( ABC \) jest:}
{trójkątem równobocznym}{trójkątem równoramiennym}{trójkątem prostokątnym}{trójkątem rozwartokątnym}{}{}}
\LANGcs{
\Question{
Jestliže mají dva vnitřní úhly v trojúhelníku \( ABC \) velikost \( 60^{\circ} \), tak je trojúhelník \( ABC \):}
{rovnostranný}{rovnoramenný}{pravoúhlý}{tupoúhlý}{}{}}
\LANGsk{
\Question{
Ak dva vnútorné uhly v trojuholníku  \( ABC \) majú veľkosť \( 60^{\circ} \), tak trojuholník \( ABC \) je:}
{rovnostranný}{rovnoramenný}{pravouhlý}{tupouhlý}{}{}}
\LANGes{
\Question{
Si dos ángulos interiores de un triángulo \( ABC \) miden \( 60^{\circ} \), se trata de un triángulo:}
{equilátero}{isósceles}{rectángulo}{obtusángulo}{}{}}
\End

\Data\ID{1003076806}
\Updated{2020-07-15 03:07:12}
\Level{A}
\SubArea{Triangles}
\LANGen{
\Question{
Choose the false statement:}
{In a triangle the side opposite the smallest interior angle is the longest side of a triangle.}{The sum of the  interior angles of a triangle is \( 180^{\circ} \).}{There is at most one obtuse interior angle in a triangle.}{In a triangle the sum of any two sides is greater than the third side.}{}{}}
\LANGpl{
\Question{
Wybierz fałszywe stwierdzenie.}
{W trójkącie strona przeciwległa do najmniejszego kąta wewnętrznego jest najdłuższym bokiem trójkąta.}{Suma miar kątów wewnętrznych trójkąta wynosi \( 180^{\circ} \).}{W trójkącie znajduje się co najwyżej jeden rozwarty kąt wewnętrzny.}{W trójkącie suma dowolnych dwóch boków jest większa niż trzeci bok.}{}{}}
\LANGcs{
\Question{
Vyberte nesprávné tvrzení:}
{Proti nejmenšímu vnitřnímu úhlu trojúhelníku leží nejdelší strana trojúhelníku.}{Součet velikostí všech vnitřních úhlů trojúhelníku je \( 180^{\circ} \).}{V trojúhelníku může být maximálně jeden vnitřní úhel tupý.}{Součet délek dvou libovolných stran trojúhelníku musí být větší než délka třetí strany.}{}{}}
\LANGsk{
\Question{
Vyberte nesprávne tvrdenie:}
{Oproti najmenšiemu vnútornému uhlu trojuholníka leží najdlhšia strana trojuholníka.}{Súčet veľkostí všetkých vnútorných uhlov trojuholníka je  \( 180^{\circ} \).}{V trojuholníku môže byť najviac jeden vnútorný uhol tupý.}{Súčet dĺžok dvoch ľubovoľných strán trojuholníka musí byť väčší ako dĺžka tretej strany.}{}{}}
\LANGes{
\Question{
Identifica la proposición falsa:}
{En un triángulo, el lado más largo es opuesto al ángulo interior más pequeño.}{La suma de los ángulos interiores de un triángulo es igual a \( 180^{\circ} \).}{Ningún triángulo puede tener más de un ángulo obtuso.}{En un triángulo, la suma de las logitudes de dos lados cualesquiera es mayor que el tercer lado.}{}{}}
\End

\Data\ID{1003076807}
\Updated{2020-07-15 03:07:12}
\Level{A}
\SubArea{Triangles}
\LANGen{
\Question{
The sum of all interior angles of an isosceles triangle with a right angle opposite the base is:}
{\( 180^{\circ} \)}{\( 90^{\circ} \)}{\( 60^{\circ} \)}{\( 30^{\circ} \)}{}{}}
\LANGpl{
\Question{
Suma miar wszystkich kątów wewnętrznych trójkąta równoramiennego z kątem prostym naprzeciw podstawy wynosi:}
{\( 180^{\circ} \)}{\( 90^{\circ} \)}{\( 60^{\circ} \)}{\( 30^{\circ} \)}{}{}}
\LANGcs{
\Question{
Součet velikostí všech vnitřních úhlů rovnoramenného trojúhelníku s pravým úhlem proti základně je:}
{\( 180^{\circ} \)}{\( 90^{\circ} \)}{\( 60^{\circ} \)}{\( 30^{\circ} \)}{}{}}
\LANGsk{
\Question{
Súčet veľkostí všetkých vnútorných uhlov rovnoramenného trojuholníka s pravým uhlom oproti základni je:}
{\( 180^{\circ} \)}{\( 90^{\circ} \)}{\( 60^{\circ} \)}{\( 30^{\circ} \)}{}{}}
\LANGes{
\Question{
La suma de las medidas de todos los ángulos interiores de un triángulo isósceles con un ángulo recto opuesto a la base es igual a:}
{\( 180^{\circ} \)}{\( 90^{\circ} \)}{\( 60^{\circ} \)}{\( 30^{\circ} \)}{}{}}
\End

\Data\ID{1003076810}
\Updated{2020-07-15 04:07:14}
\Level{A}
\SubArea{Triangles}
\LANGen{
\Question{
Interior angles of a triangle \( ABC \) are in the ratio \( 2:3:4 \). A circle is inscribed into the triangle \( ABC \). Points of tangency divide the circle into three arcs. What is the ratio of the lengths of these arcs?}
{\( 5:6:7 \)}{\( 4:5:6 \)}{\( 2:3:4 \)}{\( 3:4:5 \)}{}{}}
\LANGpl{
\Question{
Stosunek miar kątów wewnętrznych trójkąta \( ABC \) wynosi \( 2:3:4 \). W trójkąt \( ABC \) jest wpisany okrąg.  Punkty styczności dzielą okrąg na trzy łuki. Jaki jest stosunek długości tych łuków?}
{\( 5:6:7 \)}{\( 4:5:6 \)}{\( 2:3:4 \)}{\( 3:4:5 \)}{}{}}
\LANGcs{
\Question{
Vnitřní úhly trojúhelníku \( ABC \) jsou v poměru \( 2:3:4 \). Do tohoto trojúhelníku je vepsaná kružnice k. Body dotyku kružnice k se stranami trojúhelníku dělí kružnici na tři oblouky. V jakém poměru jsou délky těchto oblouků?}
{\( 5:6:7 \)}{\( 4:5:6 \)}{\( 2:3:4 \)}{\( 3:4:5 \)}{}{}}
\LANGsk{
\Question{
Vnútorné uhly trojuholníka \( ABC \) sú v pomere  \( 2:3:4 \). Do tohto trojuholníka je vpísaná kružnica k. Body dotyku kružnice k so stranami trojuholníka delia kružnicu na tri oblúky. V akom pomere sú dĺžky týchto oblúkov?}
{\( 5:6:7 \)}{\( 4:5:6 \)}{\( 2:3:4 \)}{\( 3:4:5 \)}{}{}}
\LANGes{
\Question{
Los ángulos interiores del triángulo \( ABC \) están en proporción \( 2:3:4 \). En el triángulo se inscribe una circunferencia \( k \). Los puntos de tangencia dividen a la circunferencia en tres arcos. ¿Cuál es la razón de las longitudes de estos arcos?}
{\( 5:6:7 \)}{\( 4:5:6 \)}{\( 2:3:4 \)}{\( 3:4:5 \)}{}{}}
\End

\Data\ID{1103021702}
\Updated{2021-04-19 06:04:44}
\Level{A}
\SubArea{Triangles}
\IMG{1103021702.pdf}
\LANGen{
\Question{
Given the triangle \( ABC \) (see the picture), where \( \alpha:\beta=5:7 \) and the angle \( \gamma \) is by \( 42^{\circ} \)  smaller than the angle \( \omega \), calculate the measure of \( \gamma \).}
{\( 108^{\circ} \)}{\( 42^{\circ} \)}{\( 30^{\circ} \)}{\( 60^{\circ} \)}{}{}}
\LANGpl{
\Question{
Dany jest trójkąt \( ABC \) (zobacz rysunek), gdzie \( \alpha:\beta=5:7 \), a kąt \( \gamma \) jest o  \( 42^{\circ} \)  mniejszy od kąta \( \omega \). Oblicz miarę kąta \( \gamma \).}
{\( 108^{\circ} \)}{\( 42^{\circ} \)}{\( 30^{\circ} \)}{\( 60^{\circ} \)}{}{}}
\LANGcs{
\Question{
Je dán trojúhelník \( ABC \) (viz obrázek), ve kterém \( \alpha:\beta=5:7 \) a úhel \( \gamma \)  je o  \( 42^{\circ} \)  menší než úhel \( \omega \). Vypočítejte velikost úhlu \( \gamma \).}
{\( 108^{\circ} \)}{\( 42^{\circ} \)}{\( 30^{\circ} \)}{\( 60^{\circ} \)}{}{}}
\LANGsk{
\Question{
Daný je trojuholník  \( ABC \) (pozri obrázok), v ktorom  \( \alpha:\beta=5:7 \) a uhol   \( \gamma \)  je o  \( 42^{\circ} \)  menší ako uhol  \( \omega \). Vypočítajte veľkosť uhla  \( \gamma \).}
{\( 108^{\circ} \)}{\( 42^{\circ} \)}{\( 30^{\circ} \)}{\( 60^{\circ} \)}{}{}}
\LANGes{
\Question{
Dado el triángulo \( ABC \), donde \( \alpha:\beta=5:7 \) y el ángulo \( \gamma \) es \( 42^{\circ} \)  menor que el ángulo \( \omega \), calcula la medida del ángulo \( \gamma \).}
{\( 108^{\circ} \)}{\( 42^{\circ} \)}{\( 30^{\circ} \)}{\( 60^{\circ} \)}{}{}}
\End

\Data\ID{1103021704}
\Updated{2020-07-15 03:07:12}
\Level{A}
\SubArea{Triangles}
\IMG{1103021704.pdf}
\LANGen{
\Question{
Decide according to the diagram, which of the line segments \( a \), \( b \), \( c \), \( d \), \( e \) is the longest.}
{\( d \)}{\( a \)}{\( b \)}{\( c \)}{}{}}
\LANGpl{
\Question{
Zdecyduj w odniesieniu do  rysunku, który z podanych odcinków \( a \), \( b \), \( c \), \( d \), \( e \) jest najdłuższy.}
{\( d \)}{\( a \)}{\( b \)}{\( c \)}{}{}}
\LANGcs{
\Question{
Podle údajů znázorněných na obrázku rozhodněte, která z úseček \( a \), \( b \), \( c \), \( d \), \( e \) je nejdelší.}
{\( d \)}{\( a \)}{\( b \)}{\( c \)}{}{}}
\LANGsk{
\Question{
Podľa údajov znázornených na obrázku rozhodnite, ktorá z úsečiek \( a \), \( b \), \( c \), \( d \), \( e \) je najdlhšia.}
{\( d \)}{\( a \)}{\( b \)}{\( c \)}{}{}}
\LANGes{
\Question{
A base de los datos en el diagrama decide, cuál de los segmentos \( a \), \( b \), \( c \), \( d \), \( e \) es el más largo.}
{\( d \)}{\( a \)}{\( b \)}{\( c \)}{}{}}
\End

\Data\ID{1103021706}
\Updated{2022-04-23 05:04:05}
\Level{A}
\SubArea{Triangles}
\IMG{1103021706_0.pdf}
\LANGen{
\Question{
In the triangle \( ABC \), \( \alpha=80^{\circ} \) and \( \beta=70^{\circ} \) (see the picture). Determine the measure of the angle between the altitude to the side \( AB \) and the altitude to the side \( BC \).}
{\( 70^{\circ} \)}{\( 120^{\circ} \)}{\( 30^{\circ} \)}{\( 60^{\circ} \)}{}{}}
\LANGpl{
\Question{
W trójkącie \( ABC \), \( \alpha=80^{\circ} \) i \( \beta=70^{\circ} \) (zobacz rysunek). Określ miarę kąta między wysokością opuszczoną na a bok \( AB \) i wysokością opuszczoną na bok  \( BC \).}
{\( 70^{\circ} \)}{\( 120^{\circ} \)}{\( 30^{\circ} \)}{\( 60^{\circ} \)}{}{}}
\LANGcs{
\Question{
V trojúhelníku \( ABC \) platí \( \alpha=80^{\circ} \) a \( \beta=70^{\circ} \) (viz obrázek). Jaký úhel svírá výška na stranu \( AB \) s výškou na stranu \( BC \)?}
{\( 70^{\circ} \)}{\( 120^{\circ} \)}{\( 30^{\circ} \)}{\( 60^{\circ} \)}{}{}}
\LANGsk{
\Question{
V trojuholníku  \( ABC \), platí \( \alpha=80^{\circ} \) a \( \beta=70^{\circ} \) (pozri obrázok). Aký uhol zviera výška na stranu \( AB \) s výškou na stranu  \( BC \)?}
{\( 70^{\circ} \)}{\( 120^{\circ} \)}{\( 30^{\circ} \)}{\( 60^{\circ} \)}{}{}}
\LANGes{
\Question{
Dado el triángulo \( ABC \) con \( \alpha=80^{\circ} \) y \( \beta=70^{\circ} \). ¿Qué ángulo forma la altura sobre el lado \( AB \) con la altura sobre el lado \( BC \)?}
{\( 70^{\circ} \)}{\( 120^{\circ} \)}{\( 30^{\circ} \)}{\( 60^{\circ} \)}{}{}}
\End

\Data\ID{1103021708}
\Updated{2022-04-23 05:04:05}
\Level{A}
\SubArea{Triangles}
\IMG{1103021708.pdf}
\LANGen{
\Question{
What is the measure of the angle \( \alpha \)? (See the picture.)}
{\( 7^{\circ} \)}{\( 76^{\circ} \)}{\( 83^{\circ} \)}{\( 16^{\circ} \)}{}{}}
\LANGpl{
\Question{
Jaka jest miara kąta \( \alpha \)? (Popatrz na rysunek)}
{\( 7^{\circ} \)}{\( 76^{\circ} \)}{\( 83^{\circ} \)}{\( 16^{\circ} \)}{}{}}
\LANGcs{
\Question{
Jakou velikost má úhel \( \alpha \)? (Viz obrázek)}
{\( 7^{\circ} \)}{\( 76^{\circ} \)}{\( 83^{\circ} \)}{\( 16^{\circ} \)}{}{}}
\LANGsk{
\Question{
Akú veľkosť má uhol  \( \alpha \)? (Pozri obrázok.)}
{\( 7^{\circ} \)}{\( 76^{\circ} \)}{\( 83^{\circ} \)}{\( 16^{\circ} \)}{}{}}
\LANGes{
\Question{
¿Cuál es la medida del ángulo \( \alpha \)?}
{\( 7^{\circ} \)}{\( 76^{\circ} \)}{\( 83^{\circ} \)}{\( 16^{\circ} \)}{}{}}
\End

\Data\ID{1103021709}
\Updated{2020-07-15 03:07:12}
\Level{A}
\SubArea{Triangles}
\IMG{1103021709.pdf}
\LANGen{
\Question{
The picture shows an isosceles triangle \( ABC \) with the base \( AB \). Calculate the measure of the angle \( \alpha \).}
{\( 31^{\circ} \)}{\( 62^{\circ} \)}{\( 118^{\circ} \)}{\( 59^{\circ} \)}{}{}}
\LANGpl{
\Question{
Rysunek przedstawia trójkąt równoramienny \( ABC \) o podstawie  \( AB \). Oblicz miarę kata \( \alpha \).}
{\( 31^{\circ} \)}{\( 62^{\circ} \)}{\( 118^{\circ} \)}{\( 59^{\circ} \)}{}{}}
\LANGcs{
\Question{
Na obrázku je rovnoramenný trojúhelník \( ABC \) se základnou \( AB \). Vypočítejte velikost úhlu \( \alpha \).}
{\( 31^{\circ} \)}{\( 62^{\circ} \)}{\( 118^{\circ} \)}{\( 59^{\circ} \)}{}{}}
\LANGsk{
\Question{
Na obrázku je rovnoramenný trojuholník  \( ABC \) so základňou  \( AB \). Vypočítajte veľkosť uhla  \( \alpha \).}
{\( 31^{\circ} \)}{\( 62^{\circ} \)}{\( 118^{\circ} \)}{\( 59^{\circ} \)}{}{}}
\LANGes{
\Question{
La imagen representa un triángulo isósceles \( ABC \) con la base \( AB \). Calcula la medida del ángulo \( \alpha \).}
{\( 31^{\circ} \)}{\( 62^{\circ} \)}{\( 118^{\circ} \)}{\( 59^{\circ} \)}{}{}}
\End

\Data\ID{1103076811}
\Updated{2020-07-15 03:07:12}
\Level{A}
\SubArea{Triangles}
\IMG{1103076811.pdf}
\LANGen{
\Question{
A circle is inscribed into an isosceles triangle. The base of the triangle is \( 4\,\mathrm{cm} \) long and the length of the altitude to the base is \( 10\,\mathrm{cm} \). Calculate the radius of the circle.}
{\( 1.64\,\mathrm{cm} \)}{\( 0.82\,\mathrm{cm} \)}{\( 0.20\,\mathrm{cm} \)}{\( 0.12\,\mathrm{cm} \)}{}{}}
\LANGpl{
\Question{
Okrąg jest wpisany w trójkąt równoramienny. Podstawa trójkąta ma długość \( 4\,\mathrm{cm} \), i długość wysokości do podstawy jest równa \( 10\,\mathrm{cm} \). Oblicz promień okręgu.}
{\( 1{,}64\,\mathrm{cm} \)}{\( 0{,}82\,\mathrm{cm} \)}{\( 0{,}20\,\mathrm{cm} \)}{\( 0{,}12\,\mathrm{cm} \)}{}{}}
\LANGcs{
\Question{
Do rovnoramenného trojúhelníku se základnou dlouhou \( 4\,\mathrm{cm} \) a výškou na základnu dlouhou \( 10\,\mathrm{cm} \) je vepsaná kružnice. Vypočítejte poloměr vepsané kružnice.}
{\( 1{,}64\,\mathrm{cm} \)}{\( 0{,}82\,\mathrm{cm} \)}{\( 0{,}20\,\mathrm{cm} \)}{\( 0{,}12\,\mathrm{cm} \)}{}{}}
\LANGsk{
\Question{
Do rovnoramenného trojuholníka so základňou dlhou  \( 4\,\mathrm{cm} \) a výškou na základňu dlhou \( 10\,\mathrm{cm} \) je vpísaná kružnica. Vypočítajte polomer vpísanej kružnice.}
{\( 1{,}64\,\mathrm{cm} \)}{\( 0{,}82\,\mathrm{cm} \)}{\( 0{,}20\,\mathrm{cm} \)}{\( 0{,}12\,\mathrm{cm} \)}{}{}}
\LANGes{
\Question{
Una circunferencia se inscribe en un triángulo isósceles. La base del triángulo mide \( 4\,\mathrm{cm} \) y la altura sobre la base mide \( 10\,\mathrm{cm} \). Calcula el radio de la circunferencia.}
{\( 1.64\,\mathrm{cm} \)}{\( 0.82\,\mathrm{cm} \)}{\( 0.20\,\mathrm{cm} \)}{\( 0.12\,\mathrm{cm} \)}{}{}}
\End

\Data\ID{1103076905}
\Updated{2021-04-19 06:04:35}
\Level{A}
\SubArea{Triangles}
\IMG{1103076905.pdf}
\LANGen{
\Question{
The triangle in the diagram is divided into two isosceles triangles \( AKC \) and \( KBC \) which have the same area. Calculate the size of the angle \( \beta \) if you know that the measure of \(\measuredangle AKC \) is \( 140^{\circ} \).}
{\( 70^{\circ} \)}{\( 60^{\circ} \)}{\( 50^{\circ} \)}{\( 40^{\circ} \)}{}{}}
\LANGpl{
\Question{
Trójkąt na rysunku jest podzielony na dwa trójkąty równoramienne \( AKC \) i \( KBC \) o tej samej powierzchni. Oszacuj miarę kąta \( \beta \), jeśli miara kąta \( AKC \) wynosi \( 140^{\circ} \).}
{\( 70^{\circ} \)}{\( 60^{\circ} \)}{\( 50^{\circ} \)}{\( 40^{\circ} \)}{}{}}
\LANGcs{
\Question{
Trojúhelník na obrázku je rozdělený na dva trojúhelníky \( AKC \) a \( KBC \), které jsou rovnoramenné a mají stejný obsah. Jakou velikost má úhel \( \beta \), jestliže  \(\measuredangle AKC \) má velikost \( 140^{\circ} \).}
{\( 70^{\circ} \)}{\( 60^{\circ} \)}{\( 50^{\circ} \)}{\( 40^{\circ} \)}{}{}}
\LANGsk{
\Question{
Trojuholník na obrázku je rozdelený na dva trojuholníky  \( AKC \) a \( KBC \), ktoré sú rovnoramenné  a majú rovnaký obsah. Akú veľkosť má uhol  \( \beta \), ak  \(\measuredangle AKC \) má veľkosť \( 140^{\circ} \).}
{\( 70^{\circ} \)}{\( 60^{\circ} \)}{\( 50^{\circ} \)}{\( 40^{\circ} \)}{}{}}
\LANGes{
\Question{
El triángulo de la imagen se divide en dos triángulos isósceles \( AKC \) y \( KBC \), cuyas áres son iguales. Calcula la medida del ángulo \( \beta \), si sabes que \(\measuredangle AKC \) mide \( 140^{\circ} \).}
{\( 70^{\circ} \)}{\( 60^{\circ} \)}{\( 50^{\circ} \)}{\( 40^{\circ} \)}{}{}}
\End

\Data\ID{2000003201}
\Updated{2021-12-01 04:12:22}
\Level{A}
\SubArea{Triangles}
\IMG{2000003201_7-figure0.pdf}
\LANGen{
\Question{
There is a triangle ABC in the picture. Choose the correct statement.}
{The triangle is isosceles.}{The triangle is equilateral.}{The triangle is right-angled.}{The triangle is obtuse.}{}{}}
\LANGpl{
\Question{
Na rysunku przedstawiony jest trójkąt ABC.  Wybierz poprawne stwierdzenie.}
{Trójkąt jest równoramienny.}{Trójkąt jest równoboczny.}{Trójkąt jest prostokątny.}{Trójkąt jest rozwarty.}{}{}}
\LANGcs{
\Question{
Na obrázku je trojúhelník ABC. Které tvrzení je pravdivé?}
{Trojúhelník je rovnoramenný.}{Trojúhelník je rovnostranný.}{Trojúhelník je pravoúhlý.}{Trojúhelník je tupoúhlý.}{}{}}
\LANGsk{
\Question{
Na obrázku je trojuholník ABC. Vyberte správne tvrdenie.}
{Trojuholník je rovnoramenný.}{Trojuholník  je rovnostranný.}{Trojuholník je pravouhlý.}{Trojuholník  je tupouhlý.}{}{}}
\LANGes{
\Question{
Dado el triángulo \( ABC \). ¿Cuál de las siguientes proposiciones es verdadera?}
{Es un triángulo isósceles.}{Es un triángulo equilátero.}{Es un triángulo rectángulo.}{Es un triángulo obtusángulo.}{}{}}
\End

\Data\ID{2000003202}
\Updated{2021-12-01 04:12:27}
\Level{A}
\SubArea{Triangles}
\IMG{2000003201_7-figure1.pdf}
\LANGen{
\Question{
An isosceles triangle \(ABC\) is given in the picture. What are the measures of its interior angles?}
{\( \alpha=27^{\circ};~\beta=27^{\circ};~\gamma=126^{\circ}\)}{\( \alpha=54^{\circ};~\beta=54^{\circ};~\gamma=72^{\circ}\)}{\( \alpha=63^{\circ};~\beta=63^{\circ};~\gamma=153^{\circ}\)}{\( \alpha=126^{\circ};~\beta=27^{\circ};~\gamma=27^{\circ}\)}{}{}}
\LANGpl{
\Question{
Na rysunku przedstawiono trójkąt równoramienny ABC. Jakie są miary jego kątów wewnętrznych?}
{\( \alpha=27^{\circ};~\beta=27^{\circ};~\gamma=126^{\circ}\)}{\( \alpha=54^{\circ};~\beta=54^{\circ};~\gamma=72^{\circ}\)}{\( \alpha=63^{\circ};~\beta=63^{\circ};~\gamma=153^{\circ}\)}{\( \alpha=126^{\circ};~\beta=27^{\circ};~\gamma=27^{\circ}\)}{}{}}
\LANGcs{
\Question{
Na obrázku je dán rovnoramenný trojúhelník \(ABC\). Jaká je velikost jeho vnitřních úhlů?}
{\( \alpha=27^{\circ};~\beta=27^{\circ};~\gamma=126^{\circ}\)}{\( \alpha=54^{\circ};~\beta=54^{\circ};~\gamma=72^{\circ}\)}{\( \alpha=63^{\circ};~\beta=63^{\circ};~\gamma=153^{\circ}\)}{\( \alpha=126^{\circ};~\beta=27^{\circ};~\gamma=27^{\circ}\)}{}{}}
\LANGsk{
\Question{
Na obrázku je rovnoramenný trojuholník  \(ABC\). Akú veľkosť majú jeho vnútorné uhly?}
{\( \alpha=27^{\circ};~\beta=27^{\circ};~\gamma=126^{\circ}\)}{\( \alpha=54^{\circ};~\beta=54^{\circ};~\gamma=72^{\circ}\)}{\( \alpha=63^{\circ};~\beta=63^{\circ};~\gamma=153^{\circ}\)}{\( \alpha=126^{\circ};~\beta=27^{\circ};~\gamma=27^{\circ}\)}{}{}}
\LANGes{
\Question{
Dado el triángulo isósceles \(ABC\). Calcula las medidas de los ángulos interiores.}
{\( \alpha=27^{\circ};~\beta=27^{\circ};~\gamma=126^{\circ}\)}{\( \alpha=54^{\circ};~\beta=54^{\circ};~\gamma=72^{\circ}\)}{\( \alpha=63^{\circ};~\beta=63^{\circ};~\gamma=153^{\circ}\)}{\( \alpha=126^{\circ};~\beta=27^{\circ};~\gamma=27^{\circ}\)}{}{}}
\End

\Data\ID{2000003203}
\Updated{2021-12-01 04:12:33}
\Level{A}
\SubArea{Triangles}
\IMG{2000003201_7-figure2.pdf}
\LANGen{
\Question{
A deltoid is composed of two isosceles triangles that have a common base. See the picture. Find the measures of the deltoids interior angles.}
{\( \alpha=36^{\circ};~\beta=134^{\circ};~\gamma=56^{\circ};~\delta=134^{\circ}\)}{\( \alpha=36^{\circ};~\beta=100^{\circ};~\gamma=56^{\circ};~\delta=100^{\circ}\)}{\( \alpha=56^{\circ};~\beta=134^{\circ};~\gamma=56^{\circ};~\delta=134^{\circ}\)}{\( \alpha=36^{\circ};~\beta=128^{\circ};~\gamma=56^{\circ};~\delta=128^{\circ}\)}{}{}}
\LANGpl{
\Question{
Deltoid składa się z dwóch trójkątów równoramiennych, które mają wspólną podstawę. Patrząc na rysunek, znajdź miary kątów wewnętrznych deltoidu.}
{\( \alpha=36^{\circ};~\beta=134^{\circ};~\gamma=56^{\circ};~\delta=134^{\circ}\)}{\( \alpha=36^{\circ};~\beta=100^{\circ};~\gamma=56^{\circ};~\delta=100^{\circ}\)}{\( \alpha=56^{\circ};~\beta=134^{\circ};~\gamma=56^{\circ};~\delta=134^{\circ}\)}{\( \alpha=36^{\circ};~\beta=128^{\circ};~\gamma=56^{\circ};~\delta=128^{\circ}\)}{}{}}
\LANGcs{
\Question{
Je dán deltoid, který tvoří dva rovnoramenné trojúhelníky se společnou základnou (viz obrázek). Určete velikost vnitřních úhlů deltoidu.}
{\( \alpha=36^{\circ};~\beta=134^{\circ};~\gamma=56^{\circ};~\delta=134^{\circ}\)}{\( \alpha=36^{\circ};~\beta=100^{\circ};~\gamma=56^{\circ};~\delta=100^{\circ}\)}{\( \alpha=56^{\circ};~\beta=134^{\circ};~\gamma=56^{\circ};~\delta=134^{\circ}\)}{\( \alpha=36^{\circ};~\beta=128^{\circ};~\gamma=56^{\circ};~\delta=128^{\circ}\)}{}{}}
\LANGsk{
\Question{
Deltoid na obrázku je zložený z dvoch rovnoramenných trojuholníkov, ktoré majú spoločnú základňu. Nájdite veľkosť jeho vnútorných uhlov.}
{\( \alpha=36^{\circ};~\beta=134^{\circ};~\gamma=56^{\circ};~\delta=134^{\circ}\)}{\( \alpha=36^{\circ};~\beta=100^{\circ};~\gamma=56^{\circ};~\delta=100^{\circ}\)}{\( \alpha=56^{\circ};~\beta=134^{\circ};~\gamma=56^{\circ};~\delta=134^{\circ}\)}{\( \alpha=36^{\circ};~\beta=128^{\circ};~\gamma=56^{\circ};~\delta=128^{\circ}\)}{}{}}
\LANGes{
\Question{
Dado el deltoide formado por dos triángulos isósceles con una base común. Calcula las medidas de los ángulos interiores del deltoide.}
{\( \alpha=36^{\circ};~\beta=134^{\circ};~\gamma=56^{\circ};~\delta=134^{\circ}\)}{\( \alpha=36^{\circ};~\beta=100^{\circ};~\gamma=56^{\circ};~\delta=100^{\circ}\)}{\( \alpha=56^{\circ};~\beta=134^{\circ};~\gamma=56^{\circ};~\delta=134^{\circ}\)}{\( \alpha=36^{\circ};~\beta=128^{\circ};~\gamma=56^{\circ};~\delta=128^{\circ}\)}{}{}}
\End

\Data\ID{2000003204}
\Updated{2021-12-13 07:12:09}
\Level{A}
\SubArea{Triangles}
\IMG{2000003201_7-figure3.pdf}
\LANGen{
\Question{
The picture shows a triangle \(ABC\) with a circumscribed circle \(k\) that is centered at \(S\). What is the measure of the angle \(\beta\)?}
{\( 58^{\circ}\)}{\( 32^{\circ}\)}{\( 148^{\circ}\)}{\( 64^{\circ}\)}{}{}}
\LANGpl{
\Question{
Rysunek przedstawia trójkąt \(ABC\)  z okręgiem opisanym \(k\), którego środek  \(S\) leży na boku \(AB\). Jaka jest miara kąta \(\beta\)?}
{\( 58^{\circ}\)}{\( 32^{\circ}\)}{\( 148^{\circ}\)}{\( 64^{\circ}\)}{}{}}
\LANGcs{
\Question{
Na obrázku je znázorněn trojúhelník \(ABC\) s opsanou kružnicí \(k\) se středem \(S\). Jaká je velikost úhlu \(\beta\)?}
{\( 58^{\circ}\)}{\( 32^{\circ}\)}{\( 148^{\circ}\)}{\( 64^{\circ}\)}{}{}}
\LANGsk{
\Question{
Trojuholníku \(ABC\) je opísaná kružnica \(k\) so stredom \(S\). (Pozri obrázok.) Akú veľkosť má uhol  \(\beta\)?}
{\( 58^{\circ}\)}{\( 32^{\circ}\)}{\( 148^{\circ}\)}{\( 64^{\circ}\)}{}{}}
\LANGes{
\Question{
Dado el triángulo \(ABC\) y la circunferencia circunscrita \(k\) cuyo centro es \(S\). Calcula la medida del ángulo \(\beta\).}
{\( 58^{\circ}\)}{\( 32^{\circ}\)}{\( 148^{\circ}\)}{\( 64^{\circ}\)}{}{}}
\End

\Data\ID{2000003205}
\Updated{2023-10-18 02:10:59}
\Level{A}
\SubArea{Triangles}
\IMG{2000003201_7-figure4.pdf}
\LANGen{
\Question{
There is an equilateral triangle \(ABC\) in the picture. Choose the correct statement.}
{The triangle KBC is acute-angled.}{The triangle KBC is right-angled.}{The triangle KBC is obtuse-angled.}{The triangle KBC is isosceles.}{}{}}
\LANGpl{
\Question{
Na rysunku znajduje się trójkąt równoboczny ABC. Wybierz poprawne stwierdzenie.}
{Trójkąt KBC jest ostry.}{Trójkąt KBC jest prostokątny.}{Trójkąt KBC jest rozwarty.}{Trójkąt KBC jest równoramienny.}{}{}}
\LANGcs{
\Question{
Na obrázku je znázorněn rovnostranný trojúhelník \(ABC\). Které tvrzení je pravdivé?}
{Trojúhelník KBC je ostroúhlý.}{Trojúhelník KBC je pravoúhlý.}{Trojúhelník KBC je tupoúhlý.}{Trojúhelník KBC je rovnoramenný.}{}{}}
\LANGsk{
\Question{
Na obrázku je rovnostranný trojuholník \(ABC\). Vyberte správne tvrdenie.}
{Trojuholník KBC je ostrouhlý.}{Trojuholník KBC je pravouhlý.}{Trojuholník KBC je tupouhlý.}{Trojuholník KBC je rovnoramenný.}{}{}}
\LANGes{
\Question{
Dado el triángulo equilátero \(ABC\). ¿Cuál de las proposiciones es verdadera?}
{El triángulo KBC es acutángulo.}{El triángulo KBC es rectángulo.}{El triángulo KBC es obtusángulo.}{El triángulo KBC es isósceles.}{}{}}
\End

\Data\ID{2000003206}
\Updated{2022-06-08 12:06:09}
\Level{A}
\SubArea{Triangles}
\IMG{2000003201_7-figure5.pdf}
\LANGen{
\Question{
Find the measure of the angle \(\alpha\), if the line \(a\) is parallel to the line \(b\). See the picture.}
{\( 53^{\circ}\)}{\( 55^{\circ}\)}{\( 125^{\circ}\)}{\( 72^{\circ}\)}{}{}}
\LANGpl{
\Question{
Patrząc na rysunek, znajdź miarę kąta \(\alpha\), jeśli prosta \(a\)  jest równoległa do prostej \(b\).}
{\( 53^{\circ}\)}{\( 55^{\circ}\)}{\( 125^{\circ}\)}{\( 72^{\circ}\)}{}{}}
\LANGcs{
\Question{
Určete velikost úhlu \(\alpha\), pokud platí, že přímka \(a\) je rovnoběžná s přímkou \(b\) (viz obrázek).}
{\( 53^{\circ}\)}{\( 55^{\circ}\)}{\( 125^{\circ}\)}{\( 72^{\circ}\)}{}{}}
\LANGsk{
\Question{
Nájdite veľkosť uhla  \(\alpha\), ak priamka  \(a\) je rovnobežná s priamkou  \(b\). Pozri obrázok.}
{\( 53^{\circ}\)}{\( 55^{\circ}\)}{\( 125^{\circ}\)}{\( 72^{\circ}\)}{}{}}
\LANGes{
\Question{
Calcula la medida del ángulo \(\alpha\), suponiendo que las rectas \(a\) y \(b\) son paralelas.}
{\( 53^{\circ}\)}{\( 55^{\circ}\)}{\( 125^{\circ}\)}{\( 72^{\circ}\)}{}{}}
\End

\Data\ID{2000003207}
\Updated{2021-12-01 04:12:30}
\Level{A}
\SubArea{Triangles}
\IMG{2000003201_7-figure6.pdf}
\LANGen{
\Question{
Two isosceles obtuse-angled triangles are symmetrical about the axis \(o\). See the picture. What is the measure of the angle \(\delta\)?}
{\( 280^{\circ}\)}{\( 160^{\circ}\)}{\( 320^{\circ}\)}{\( 340^{\circ}\)}{}{}}
\LANGpl{
\Question{
Popatrz na rysunek. Dwa trójkąty równoramienne rozwarte są symetryczne wokół osi \(o\). Jaka jest miara kąta \(\delta\)?}
{\( 280^{\circ}\)}{\( 160^{\circ}\)}{\( 320^{\circ}\)}{\( 340^{\circ}\)}{}{}}
\LANGcs{
\Question{
Dva rovnoramenné tupoúhlé trojúhelníky na obrázku jsou symetrické podle osy \(o\). Jaká je velikost úhu \(\delta\)?}
{\( 280^{\circ}\)}{\( 160^{\circ}\)}{\( 320^{\circ}\)}{\( 340^{\circ}\)}{}{}}
\LANGsk{
\Question{
Dva rovnoramenné, tupouhlé trojuholníky sú súmerné podľa osi  \(o\). Pozri obrázok. Akú veľkosť má uhol \(\delta\)?}
{\( 280^{\circ}\)}{\( 160^{\circ}\)}{\( 320^{\circ}\)}{\( 340^{\circ}\)}{}{}}
\LANGes{
\Question{
Dos triángulos isósceles de ángulos obtusos son simétricos respecto del eje \(o\). Calcula la medida del ángulo \(\delta\).}
{\( 280^{\circ}\)}{\( 160^{\circ}\)}{\( 320^{\circ}\)}{\( 340^{\circ}\)}{}{}}
\End

\Data\ID{2000003208}
\Updated{2021-12-01 04:12:38}
\Level{A}
\SubArea{Triangles}
\IMG{2000003201_7-figure7.pdf}
\LANGen{
\Question{
In the picture the angle \(\gamma\) is marked. What is the measure of \(\gamma\)?}
{\( 9^{\circ}\)}{\( 58^{\circ}\)}{\( 67^{\circ}\)}{\( 125^{\circ}\)}{}{}}
\LANGpl{
\Question{
Na rysunku zaznaczono kąt \(\gamma\). Jaka jest miara kąta \(\gamma\)?}
{\( 9^{\circ}\)}{\( 58^{\circ}\)}{\( 67^{\circ}\)}{\( 125^{\circ}\)}{}{}}
\LANGcs{
\Question{
Na obrázku je znázorněn úhel \(\gamma\). Jaká je jeho velikost?}
{\( 9^{\circ}\)}{\( 58^{\circ}\)}{\( 67^{\circ}\)}{\( 125^{\circ}\)}{}{}}
\LANGsk{
\Question{
Nájdite veľkosť uhla  \(\gamma\). Pozri obrázok.}
{\( 9^{\circ}\)}{\( 58^{\circ}\)}{\( 67^{\circ}\)}{\( 125^{\circ}\)}{}{}}
\LANGes{
\Question{
Calcula la medida del ángulo \(\gamma\).}
{\( 9^{\circ}\)}{\( 58^{\circ}\)}{\( 67^{\circ}\)}{\( 125^{\circ}\)}{}{}}
\End

\Data\ID{2010015201}
\Updated{2022-08-25 10:08:01}
\Level{A}
\SubArea{Triangles}
\LANGen{
\Question{
Interior angles of a triangle \( ABC \) are in the ratio \( \alpha:\beta:\gamma=3:5:7 \). Calculate the measures of these angles.}
{\( \alpha=36^{\circ};\ \beta=60^{\circ};\ \gamma=84^{\circ} \)}{\( \alpha=30^{\circ};\ \beta=50^{\circ};\ \gamma=70^{\circ} \)}{\( \alpha=16.5^{\circ};\ \beta=30^{\circ};\ \gamma=133.5^{\circ} \)}{\( \alpha=84^{\circ};\ \beta=60^{\circ};\ \gamma=36^{\circ} \)}{}{}}
\LANGpl{
\Question{
Kąty wewnętrzne trójkąta \( ABC \) są w stosunku \( \alpha:\beta:\gamma=3:5:7 \). Oblicz miary tych kątów.}
{\( \alpha=36^{\circ};\ \beta=60^{\circ};\ \gamma=84^{\circ} \)}{\( \alpha=30^{\circ};\ \beta=50^{\circ};\ \gamma=70^{\circ} \)}{\( \alpha=16{,}5^{\circ};\ \beta=30^{\circ};\ \gamma=133{,}5^{\circ} \)}{\( \alpha=84^{\circ};\ \beta=60^{\circ};\ \gamma=36^{\circ} \)}{}{}}
\LANGcs{
\Question{
Vnitřní úhly trojúhelníku \( ABC \) jsou v poměru \( \alpha:\beta:\gamma=3:5:7 \). Vypočtěte velikost těchto úhlů.}
{\( \alpha=36^{\circ};\ \beta=60^{\circ};\ \gamma=84^{\circ} \)}{\( \alpha=30^{\circ};\ \beta=50^{\circ};\ \gamma=70^{\circ} \)}{\( \alpha=16{,}5^{\circ};\ \beta=30^{\circ};\ \gamma=133{,}5^{\circ} \)}{\( \alpha=84^{\circ};\ \beta=60^{\circ};\ \gamma=36^{\circ} \)}{}{}}
\LANGsk{
\Question{
Vnútorné uhly trojuholníka  \( ABC \) sú v pomere \( \alpha:\beta:\gamma=3:5:7 \). Vypočítajte veľkosti týchto uhlov.}
{\( \alpha=36^{\circ};\ \beta=60^{\circ};\ \gamma=84^{\circ} \)}{\( \alpha=30^{\circ};\ \beta=50^{\circ};\ \gamma=70^{\circ} \)}{\( \alpha=16{,}5^{\circ};\ \beta=30^{\circ};\ \gamma=133{,}5^{\circ} \)}{\( \alpha=84^{\circ};\ \beta=60^{\circ};\ \gamma=36^{\circ} \)}{}{}}
\LANGes{
\Question{
Los ángulos interiores de un triángulo \( ABC \) están en proporción \( \alpha:\beta:\gamma=3:5:7 \). Calcula las medidas de los ángulos.}
{\( \alpha=36^{\circ};\ \beta=60^{\circ};\ \gamma=84^{\circ} \)}{\( \alpha=30^{\circ};\ \beta=50^{\circ};\ \gamma=70^{\circ} \)}{\( \alpha=16.5^{\circ};\ \beta=30^{\circ};\ \gamma=133.5^{\circ} \)}{\( \alpha=84^{\circ};\ \beta=60^{\circ};\ \gamma=36^{\circ} \)}{}{}}
\End

\Data\ID{2010015205}
\Updated{2022-08-25 10:08:04}
\Level{A}
\SubArea{Triangles}
\LANGen{
\Question{
The measures of two interior angles of a triangle are \( 61^{\circ}20' \) and \( 28^{\circ} \). What is the measure of the third interior angle?}
{\( 90^{\circ}40' \)}{\( 33^{\circ}20' \)}{\( 145^{\circ}20' \)}{\( 147^{\circ}40' \)}{}{}}
\LANGpl{
\Question{
Miary dwóch wewnętrznych kątów trójkąta to \( 61^{\circ}20' \) i \( 28^{\circ} \). Jaka jest miara trzeciego kąta wewnętrznego?}
{\( 90^{\circ}40' \)}{\( 33^{\circ}20' \)}{\( 145^{\circ}20' \)}{\( 147^{\circ}40' \)}{}{}}
\LANGcs{
\Question{
Velikost dvou vnitřních úhlů trojúhelníka je \( 61^{\circ}20' \) a \( 28^{\circ} \). Jaká je velikost třetího vnitřního úhlu?}
{\( 90^{\circ}40' \)}{\( 33^{\circ}20' \)}{\( 145^{\circ}20' \)}{\( 147^{\circ}40' \)}{}{}}
\LANGsk{
\Question{
Veľkosti dvoch vnútorných uhlov trojuholníka sú \( 61^{\circ}20' \) a \( 28^{\circ} \). Aká je veľkosť tretieho vnútorného uhla?}
{\( 90^{\circ}40' \)}{\( 33^{\circ}20' \)}{\( 145^{\circ}20' \)}{\( 147^{\circ}40' \)}{}{}}
\LANGes{
\Question{
Las medidas de dos ángulos interiores de un triángulo son \( 61^{\circ}20' \) y \( 28^{\circ} \). ¿Cuánto mide el tercer ángulo interior?}
{\( 90^{\circ}40' \)}{\( 33^{\circ}20' \)}{\( 145^{\circ}20' \)}{\( 147^{\circ}40' \)}{}{}}
\End

\Data\ID{2010015208}
\Updated{2022-08-25 10:08:04}
\Level{A}
\SubArea{Triangles}
\IMG{2010015201-figure2.pdf}
\LANGen{
\Question{
In the triangle \( ABC \), \( \alpha=80^{\circ} \) and \( \gamma=30^{\circ} \) (see the picture). Determine the measure of the angle between the altitude to the side \( AC \) and the altitude to the side \( AB \).}
{\( 80^{\circ} \)}{\(30^{\circ}\)}{\(70^{\circ}\)}{\(100^{\circ}\)}{}{}}
\LANGpl{
\Question{
W trójkącie \( ABC \), \( \alpha=80^{\circ} \) i \( \gamma=30^{\circ} \) (popatrz na rysunek). Wyznacz miarę kąta między wysokością do boku \( AC \) a wysokością do boku \( AB \).}
{\( 80^{\circ} \)}{\(30^{\circ}\)}{\(70^{\circ}\)}{\(100^{\circ}\)}{}{}}
\LANGcs{
\Question{
Je dán trojúhelník \( ABC \), kde \( \alpha=80^{\circ} \) a \( \gamma=30^{\circ} \) (viz obrázek). Jaký úhel svírá výška na stranu \( AC \) a výška na stranu \( AB \)?}
{\( 80^{\circ} \)}{\(30^{\circ}\)}{\(70^{\circ}\)}{\(100^{\circ}\)}{}{}}
\LANGsk{
\Question{
V trojuholníku \( ABC \), \( \alpha=80^{\circ} \) a \( \gamma=30^{\circ} \) (pozri obrázok). Aký uhol zviera výška na stranu  \( AC \) s výškou na stranu \( AB \).}
{\( 80^{\circ} \)}{\(30^{\circ}\)}{\(70^{\circ}\)}{\(100^{\circ}\)}{}{}}
\LANGes{
\Question{
Dado el triángulo \( ABC \), \( \alpha=80^{\circ} \) y \( \gamma=30^{\circ} \). ¿Qué ángulo forma la altura sobre el lado \( AC \) con la altura sobre el lado \( AB \)?}
{\( 80^{\circ} \)}{\(30^{\circ}\)}{\(70^{\circ}\)}{\(100^{\circ}\)}{}{}}
\End

\Data\ID{2010015209}
\Updated{2022-08-25 10:08:04}
\Level{A}
\SubArea{Triangles}
\IMG{2010015201-figure3.pdf}
\LANGen{
\Question{
What is the measure of the angle \( \alpha \)? (See the picture.)}
{\( 72^{\circ} \)}{\(44^{\circ}\)}{\(88^{\circ}\)}{\(108^{\circ}\)}{}{}}
\LANGpl{
\Question{
Jaka jest miara kąta \( \alpha \)? (Popatrz na rysunek.)}
{\( 72^{\circ} \)}{\(44^{\circ}\)}{\(88^{\circ}\)}{\(108^{\circ}\)}{}{}}
\LANGcs{
\Question{
Jaká je velikost úhlu \( \alpha \)? (Viz obrázek.)}
{\( 72^{\circ} \)}{\(44^{\circ}\)}{\(88^{\circ}\)}{\(108^{\circ}\)}{}{}}
\LANGsk{
\Question{
Akú veľkosť má uhol  \( \alpha \)? (Pozri obrázok.)}
{\( 72^{\circ} \)}{\(44^{\circ}\)}{\(88^{\circ}\)}{\(108^{\circ}\)}{}{}}
\LANGes{
\Question{
¿Cuál es la medida del ángulo \( \alpha \)?}
{\( 72^{\circ} \)}{\(44^{\circ}\)}{\(88^{\circ}\)}{\(108^{\circ}\)}{}{}}
\End

\Data\ID{9000121701}
\Updated{2020-07-15 03:07:37}
\Level{A}
\SubArea{Triangles}
\LANGen{
\Question{
Consider a triangle \(ABC\) with
sides of the lengths \(3\, \mathrm{cm}\),
\(4\, \mathrm{cm}\) and
\(4\, \mathrm{cm}\). Identify the type
of the triangle \(ABC\).}
{isosceles}{equilateral}{right}{obtuse}{}{}}
\LANGpl{
\Question{
Rozważ trójkąt  \(ABC\) o bokach długości \(3\, \mathrm{cm}\),
\(4\, \mathrm{cm}\) i
\(4\, \mathrm{cm}\). Jaki to rodzaj trójkąta?}
{równoramienny}{równoboczny}{prostokątny}{rozwartokątny}{}{}}
\LANGcs{
\Question{
Je dán trojúhelník \(ABC\),
jehož strany mají délky \(3\, \mathrm{cm}\),
\(4\, \mathrm{cm}\) a
\(4\, \mathrm{cm}\). Trojúhelník
\(ABC\) je:}
{rovnoramenný}{rovnostranný}{pravoúhlý}{tupoúhlý}{}{}}
\LANGsk{
\Question{
Je daný trojuholník \(ABC\),
ktorého strany majú dĺžky \(3\, \mathrm{cm}\),
\(4\, \mathrm{cm}\) a
\(4\, \mathrm{cm}\). Trojuholník
\(ABC\) je:}
{rovnoramenný}{rovnostranný}{pravouhlý}{tupouhlý}{}{}}
\LANGes{
\Question{
Dado un triángulo \(ABC\) cuyos lados miden \(3\, \mathrm{cm}\),
\(4\, \mathrm{cm}\) y
\(4\, \mathrm{cm}\). Se trata de un triángulo:}
{isósceles}{equilátero}{rectángulo}{obtusángulo}{}{}}
\End

\Data\ID{9000121702}
\Updated{2020-07-15 03:07:37}
\Level{A}
\SubArea{Triangles}
\LANGen{
\Question{
Consider a triangle \(ABC\) with
sides of the lengths \(3\, \mathrm{cm}\),
\(4\, \mathrm{cm}\) and
\(5\, \mathrm{cm}\). Identify the type
of the triangle \(ABC\).}
{right}{isosceles}{equilateral}{obtuse}{}{}}
\LANGpl{
\Question{
Rozważ trójkąt \(ABC\) o bokach długości \(3\, \mathrm{cm}\),
\(4\, \mathrm{cm}\) i
\(5\, \mathrm{cm}\). Określ rodzaj tego trójkąta.}
{prostokątny}{równoramienny}{równoboczny}{rozwartokątny}{}{}}
\LANGcs{
\Question{
Je dán trojúhelník \(ABC\),
jehož strany mají délky \(3\, \mathrm{cm}\),
\(4\, \mathrm{cm}\) a
\(5\, \mathrm{cm}\). Trojúhelník
\(ABC\) je:}
{pravoúhlý}{rovnoramenný}{rovnostranný}{tupoúhlý}{}{}}
\LANGsk{
\Question{
Je daný trojuholník \(ABC\),
ktorého strany majú dĺžky \(3\, \mathrm{cm}\),
\(4\, \mathrm{cm}\) a
\(5\, \mathrm{cm}\). Trojuholník
\(ABC\) je:}
{pravouhlý}{rovnoramenný}{rovnostranný}{tupouhlý}{}{}}
\LANGes{
\Question{
Dado un triángulo \(ABC\) cuyos lados miden \(3\, \mathrm{cm}\),
\(4\, \mathrm{cm}\) y
\(5\, \mathrm{cm}\). Se trata de un triángulo:}
{rectángulo}{isósceles}{equilátero}{obtusángulo}{}{}}
\End

\Data\ID{9000121703}
\Updated{2020-07-15 03:07:37}
\Level{A}
\SubArea{Triangles}
\LANGen{
\Question{
Consider a triangle \(ABC\) with
sides of the lengths \(4\, \mathrm{cm}\),
\(4\, \mathrm{cm}\) and
\(4\, \mathrm{cm}\). Identify the type
of the triangle \(ABC\).}
{equilateral}{isosceles}{right}{obtuse}{}{}}
\LANGpl{
\Question{
Rozważ trójkąt  \(ABC\) o bokach długości \(4\, \mathrm{cm}\),
\(4\, \mathrm{cm}\) i 
\(4\, \mathrm{cm}\). Jaki to rodzaj trójkąta?}
{równoboczny}{równoramienny}{prostokątny}{rozwartokątny}{}{}}
\LANGcs{
\Question{
Je dán trojúhelník \(ABC\),
jehož strany mají délky \(4\, \mathrm{cm}\),
\(4\, \mathrm{cm}\) a
\(4\, \mathrm{cm}\). Trojúhelník
\(ABC\) je:}
{rovnostranný}{rovnoramenný}{pravoúhlý}{tupoúhlý}{}{}}
\LANGsk{
\Question{
Je daný trojuholník \(ABC\),
ktorého strany majú dĺžky \(4\, \mathrm{cm}\),
\(4\, \mathrm{cm}\) a
\(4\, \mathrm{cm}\). Trojuholník
\(ABC\) je:}
{rovnostranný}{rovnoramenný}{pravouhlý}{tupouhlý}{}{}}
\LANGes{
\Question{
Dado un triángulo \(ABC\) cuyos lados miden \(4\, \mathrm{cm}\),
\(4\, \mathrm{cm}\) y
\(4\, \mathrm{cm}\). Se trata de un triángulo:}
{equilátero}{isósceles}{rectángulo}{obtusángulo}{}{}}
\End

\Data\ID{9000121704}
\Updated{2020-07-15 03:07:37}
\Level{A}
\SubArea{Triangles}
\LANGen{
\Question{
Find the measure of the angle \(CAB\) 
in an isosceles triangle \(ABC\) 
with the interior angle \( BCA\) 
of the measure \(120^{\circ }\).}
{\(30^{\circ }\)}{\(60^{\circ }\)}{\(15^{\circ }\)}{\(45^{\circ }\)}{}{}}
\LANGpl{
\Question{
Oblicz miarę kąta  \( CAB\) 
w trójkącie równoramiennym \(ABC\) 
o kącie wewnętrznym \(BCA\), którego miara wynosi
 \(120^{\circ }\).}
{\(30^{\circ }\)}{\(60^{\circ }\)}{\(15^{\circ }\)}{\(45^{\circ }\)}{}{}}
\LANGcs{
\Question{
Je dán rovnoramenný trojúhelník
\(ABC\), ve
kterém \(|\measuredangle BCA| = 120^{\circ }\).
Určete \(|\measuredangle CAB|\).}
{\(30^{\circ }\)}{\(60^{\circ }\)}{\(15^{\circ }\)}{\(45^{\circ }\)}{}{}}
\LANGsk{
\Question{
Je daný rovnoramenný trojuholník
\(ABC\), v
ktorom \(|\measuredangle BCA| = 120^{\circ }\).
Určte \(|\measuredangle CAB|\).}
{\(30^{\circ }\)}{\(60^{\circ }\)}{\(15^{\circ }\)}{\(45^{\circ }\)}{}{}}
\LANGes{
\Question{
Dado un triángulo isósceles
\(ABC\), en el que \(|\measuredangle BCA| = 120^{\circ }\).
Halla \(|\measuredangle CAB|\).}
{\(30^{\circ }\)}{\(60^{\circ }\)}{\(15^{\circ }\)}{\(45^{\circ }\)}{}{}}
\End

\Data\ID{9000121705}
\Updated{2022-07-21 05:07:58}
\Level{A}
\SubArea{Triangles}
\LANGen{
\Question{
Consider an isosceles triangle \(ABC\) with sides \(AC\) and \(BC\) of equal length. The measure of the angle \( BAC\) is \(40^{\circ }\). \(X\) is the point of intersection between the line $AB$ and the line through the vertex \(C\) perpendicular to it. Find the measure of the angle \( BCX\).}
{\(50^{\circ }\)}{\(80^{\circ }\)}{\(100^{\circ }\)}{\(40^{\circ }\)}{}{}}
\LANGpl{
\Question{
Rozważmy trójkąt równoramienny \(ABC\) o  bokach \(AC\) i \(BC\) równej długości. Miara kąta  \( BAC\) wynosi \(40^{\circ }\). Punkt \(X\) jest przecięciem linii \( AB \) i linii przez wierzchołek \(C\) prostopadłe do niej.  Znajdź miarę kąta \( BCX\).}
{\(50^{\circ }\)}{\(80^{\circ }\)}{\(100^{\circ }\)}{\(40^{\circ }\)}{}{}}
\LANGcs{
\Question{
Je dán rovnoramenný trojúhelník
\(ABC\), ve kterém
\(|\measuredangle BAC| = 40^{\circ }\). Bod
\(X\) je pata kolmice
vedené z bodu \(C\) 
na základnu \(c\).
Určete \(|\measuredangle BCX|\).}
{\(50^{\circ }\)}{\(80^{\circ }\)}{\(100^{\circ }\)}{\(40^{\circ }\)}{}{}}
\LANGsk{
\Question{
Je daný rovnoramenný trojuholník
\(ABC\), v ktorom
\(|\measuredangle BAC| = 40^{\circ }\). Bod
\(X\) je päta kolmice
vedené z bodu \(C\) 
na základňu \(c\).
Určte \(|\measuredangle BCX|\).}
{\(50^{\circ }\)}{\(80^{\circ }\)}{\(100^{\circ }\)}{\(40^{\circ }\)}{}{}}
\LANGes{
\Question{
Dado un triángulo isósceles \(ABC\) con los lados \(AC\) y \(BC\) de la misma longitud. La medida del ángulo \( BAC\) es \(40^{\circ }\). \(X\) es el punto de intersección entre la recta $AB$ y la recta que pasa por el vértice \(C\) y es perpendicular a la primera. Calcula la medida del ángulo \( BCX\).}
{\(50^{\circ }\)}{\(80^{\circ }\)}{\(100^{\circ }\)}{\(40^{\circ }\)}{}{}}
\End

\Data\ID{1003021803}
\Updated{2022-08-25 10:08:24}
\Level{B}
\SubArea{Triangles}
\IMG{2010015201-figure0.pdf}
\LANGen{
\Question{
The ladder is leaning against the wall of a house. Its length is \( 6 \) meters. How high does the ladder reach if the angle between it and the wall is \( 30^{\circ} \)? (See the picture.)}
{\( 3\sqrt3\,\mathrm{m} \)}{\( 3\,\mathrm{m} \)}{\( 6\,\mathrm{m} \)}{\( \frac{\sqrt3}2\,\mathrm{m} \)}{}{}}
\LANGpl{
\Question{
Drabina opiera się o ścianę budynku. Jej długość wynosi \( 6 \) metrów. Jak wysoko sięga drabina, jeśli kąt pomiędzy nią a ścianą ma miarę \( 30^{\circ} \)?}
{\( 3\sqrt3\,\mathrm{m} \)}{\( 3\,\mathrm{m} \)}{\( 6\,\mathrm{m} \)}{\( \frac{\sqrt3}2\,\mathrm{m} \)}{}{}}
\LANGcs{
\Question{
Žebřík se opírá o stěnu domu. Jeho délka je \( 6 \) metrů. Do jaké výšky stěny žebřík dosáhne, jestliže svírá se stěnou úhel \( 30^{\circ} \)?}
{\( 3\sqrt3\,\mathrm{m} \)}{\( 3\,\mathrm{m} \)}{\( 6\,\mathrm{m} \)}{\( \frac{\sqrt3}2\,\mathrm{m} \)}{}{}}
\LANGsk{
\Question{
Rebrík sa opiera o stenu domu. Jeho dĺžka je  \( 6 \) metrov. Do akej výšky steny rebrík dosiahne, ak so stenou zviera rebrík uhol \( 30^{\circ} \)?}
{\( 3\sqrt3\,\mathrm{m} \)}{\( 3\,\mathrm{m} \)}{\( 6\,\mathrm{m} \)}{\( \frac{\sqrt3}2\,\mathrm{m} \)}{}{}}
\LANGes{
\Question{
Una escalera se apoya en la pared de una casa. La escalera mide \( 6 \) metros. ¿A qué altura llega la escalera si forma un ángulo de \( 30^{\circ} \) con la pared?}
{\( 3\sqrt3\,\mathrm{m} \)}{\( 3\,\mathrm{m} \)}{\( 6\,\mathrm{m} \)}{\( \frac{\sqrt3}2\,\mathrm{m} \)}{}{}}
\End

\Data\ID{1003021805}
\Updated{2022-08-25 10:08:24}
\Level{B}
\SubArea{Triangles}
\LANGen{
\Question{
In a triangle with interior angles \( 30^{\circ} \), \( 60^{\circ} \) and \( 90^{\circ} \), the longest side is \( 10\,\mathrm{cm} \) long. Find the length of its shortest side.}
{\( 5\,\mathrm{cm} \)}{\( 5\sqrt3\,\mathrm{cm} \)}{\( 3\,\mathrm{cm} \)}{\( 8\,\mathrm{cm} \)}{}{}}
\LANGpl{
\Question{
W trójkącie kąty wewnętrzne równe są  \( 30^{\circ} \), \( 60^{\circ} \) i \( 90^{\circ} \), najdłuższy bok ma długość \( 10\,\mathrm{cm} \). Wyznacz długość najkrótszego boku.}
{\( 5\,\mathrm{cm} \)}{\( 5\sqrt3\,\mathrm{cm} \)}{\( 3\,\mathrm{cm} \)}{\( 8\,\mathrm{cm} \)}{}{}}
\LANGcs{
\Question{
Vnitřní úhly trojúhelníku mají velikost \( 30^{\circ} \), \( 60^{\circ} \) a \( 90^{\circ} \). Nejdelší strana trojúhelníku měří \( 10\,\mathrm{cm} \).  Vypočítejte délku nejkratší strany.}
{\( 5\,\mathrm{cm} \)}{\( 5\sqrt3\,\mathrm{cm} \)}{\( 3\,\mathrm{cm} \)}{\( 8\,\mathrm{cm} \)}{}{}}
\LANGsk{
\Question{
Vnútorné uhly trojuholníka majú veľkosti \( 30^{\circ} \), \( 60^{\circ} \) a \( 90^{\circ} \). Najdlhšia strana trojuholníka meria \( 10\,\mathrm{cm} \).  Vypočítajte dĺžku najkratšej strany.}
{\( 5\,\mathrm{cm} \)}{\( 5\sqrt3\,\mathrm{cm} \)}{\( 3\,\mathrm{cm} \)}{\( 8\,\mathrm{cm} \)}{}{}}
\LANGes{
\Question{
En un triángulo con ángulos interiores de \( 30^{\circ} \), \( 60^{\circ} \) y \( 90^{\circ} \), el lado más largo mide \( 10\,\mathrm{cm} \). Calcula la longitud de su lado más corto.}
{\( 5\,\mathrm{cm} \)}{\( 5\sqrt3\,\mathrm{cm} \)}{\( 3\,\mathrm{cm} \)}{\( 8\,\mathrm{cm} \)}{}{}}
\End

\Data\ID{1003021809}
\Updated{2022-08-25 10:08:24}
\Level{B}
\SubArea{Triangles}
\LANGen{
\Question{
\( ABC \) is a right-angled triangle with the right angle at vertex \( C \), side \( b=10\,\mathrm{cm} \) and the altitude to the hypotenuse \( v_c=5\,\mathrm{cm} \). Find the measure of the angle \( BAC \).}
{\( 30^{\circ} \)}{\( 45^{\circ} \)}{\( 60^{\circ} \)}{\( 90^{\circ} \)}{}{}}
\LANGpl{
\Question{
W trójkącie prostokątnym \( ABC \) z kątem prostym w wierzchołku \( C \) dane są bok \( b=10\,\mathrm{cm} \) a wysokość do przeciwprostokątnej \( v_c=5\,\mathrm{cm} \). Znajdź miarę kata \( BAC \).}
{\( 30^{\circ} \)}{\( 45^{\circ} \)}{\( 60^{\circ} \)}{\( 90^{\circ} \)}{}{}}
\LANGcs{
\Question{
V pravoúhlém trojúhelníku  \( ABC \) s pravým úhlem u vrcholu \( C \) je dána strana \( b=10\,\mathrm{cm} \) a výška na přeponu  \( v_c=5\,\mathrm{cm} \). Vypočítejte velikost úhlu \( BAC \).}
{\( 30^{\circ} \)}{\( 45^{\circ} \)}{\( 60^{\circ} \)}{\( 90^{\circ} \)}{}{}}
\LANGsk{
\Question{
V pravouhlom trojuholníku  \( ABC \) s pravým uhlom pri vrchole \( C \) je daná  strana  \( b=10\,\mathrm{cm} \) a výška na preponu   \( v_c=5\,\mathrm{cm} \). Vypočítajte veľkosť uhla \( BAC \).}
{\( 30^{\circ} \)}{\( 45^{\circ} \)}{\( 60^{\circ} \)}{\( 90^{\circ} \)}{}{}}
\LANGes{
\Question{
En un triángulo rectángulo \( ABC \), siendo \( C \) el vértice del ángulo recto, dados el lado \( b=10\,\mathrm{cm} \) y la altura sobre la hipotenusa \( v_c=5\,\mathrm{cm} \). Calcula la medida del ángulo \( BAC \).}
{\( 30^{\circ} \)}{\( 45^{\circ} \)}{\( 60^{\circ} \)}{\( 90^{\circ} \)}{}{}}
\End

\Data\ID{1003076808}
\Updated{2020-07-15 03:07:12}
\Level{B}
\SubArea{Triangles}
\LANGen{
\Question{
In a triangle \( ABC \) the measure of \( \measuredangle CAB \) is \( 45^{\circ} \) and the measure of \( \measuredangle CBA \) is \( 60^{\circ} \). The altitude to side \( AB \) is \( 1\,\mathrm{cm} \) long. Calculate the area of the triangle \( ABC \) in \(\mathrm{cm}^2 \).}
{\( \frac{\sqrt3+1}{2\sqrt3} \)}{\( \frac{\sqrt3+1}{\sqrt3} \)}{\( \frac{\sqrt3+1}{2} \)}{\( \frac{\sqrt3+1}{4} \)}{}{}}
\LANGpl{
\Question{
W trójkącie \( ABC \) miara kąta \( CAB \) wynosi  \( 45^{\circ} \), a miara kąta \( CBA \) wynosi  \( 60^{\circ} \). Wysokość do boku \( AB \) ma długość \( 1\,\mathrm{cm} \).   Oblicz pole tego trójkąta \( ABC \) w \(\mathrm{cm}^2 \).}
{\( \frac{\sqrt3+1}{2\sqrt3} \)}{\( \frac{\sqrt3+1}{\sqrt3} \)}{\( \frac{\sqrt3+1}{2} \)}{\( \frac{\sqrt3+1}{4} \)}{}{}}
\LANGcs{
\Question{
V trojúhelníku \( ABC \) má  \( \measuredangle CAB \) velikost \( 45^{\circ} \)  a  \( \measuredangle CBA \) má velikost \( 60^{\circ} \). Výška na stranu  \( AB \) má délku \( 1\,\mathrm{cm} \). Vypočítejte v \(\mathrm{cm}^2 \) obsah trojúhelníku \( ABC \).}
{\( \frac{\sqrt3+1}{2\sqrt3} \)}{\( \frac{\sqrt3+1}{\sqrt3} \)}{\( \frac{\sqrt3+1}{2} \)}{\( \frac{\sqrt3+1}{4} \)}{}{}}
\LANGsk{
\Question{
V trojuholníku  \( ABC \) má  \( \measuredangle CAB \) veľkosť \( 45^{\circ} \)  a  \( \measuredangle CBA \) má veľkosť \( 60^{\circ} \). Výška na stranu  \( AB \) má dĺžku \( 1\,\mathrm{cm} \). Vypočítajte v \(\mathrm{cm}^2 \) obsah trojuholníka  \( ABC \).}
{\( \frac{\sqrt3+1}{2\sqrt3} \)}{\( \frac{\sqrt3+1}{\sqrt3} \)}{\( \frac{\sqrt3+1}{2} \)}{\( \frac{\sqrt3+1}{4} \)}{}{}}
\LANGes{
\Question{
En un triángulo \( ABC \), la medida de \( \measuredangle CAB \) es \( 45^{\circ} \) y la medida de \( \measuredangle CBA \) es \( 60^{\circ} \). La altura sobre el lado \( AB \) mide \( 1\,\mathrm{cm} \). Calcula el área del triángulo \( ABC \) en \(\mathrm{cm}^2 \).}
{\( \frac{\sqrt3+1}{2\sqrt3} \)}{\( \frac{\sqrt3+1}{\sqrt3} \)}{\( \frac{\sqrt3+1}{2} \)}{\( \frac{\sqrt3+1}{4} \)}{}{}}
\End

\Data\ID{1003076906}
\Updated{2022-04-23 05:04:05}
\Level{B}
\SubArea{Triangles}
\LANGen{
\Question{
The lengths of the sides in a triangle are \( a \), \( b \), \( c \) and the opposite angles are \( \alpha \), \( \beta \), \( \gamma \). Give the measure of \( \alpha \) if \( a^2 = b^2 + c^2 +bc \).}
{\( 120^{\circ} \)}{\( 60^{\circ} \)}{\( 30^{\circ} \)}{\( 90^{\circ} \)}{}{}}
\LANGpl{
\Question{
Długości boków w trójkącie są  \( a \), \( b \), \( c \) i przeciwległe kąty są \( \alpha \), \( \beta \), \( \gamma \). Oblicz miarę kąta \( \alpha \) jeśli \( a^2 = b^2 + c^2 +bc \).}
{\( 120^{\circ} \)}{\( 60^{\circ} \)}{\( 30^{\circ} \)}{\( 90^{\circ} \)}{}{}}
\LANGcs{
\Question{
Délky stran v trojúhelníku jsou \( a \), \( b \), \( c \) a vnitřní úhly \( \alpha \), \( \beta \), \( \gamma \). Vypočítejte velikost úhlu \( \alpha \), jestliže \( a^2 = b^2 + c^2 +bc \).}
{\( 120^{\circ} \)}{\( 60^{\circ} \)}{\( 30^{\circ} \)}{\( 90^{\circ} \)}{}{}}
\LANGsk{
\Question{
Dĺžky strán v trojuholníku sú \( a \), \( b \), \( c \) a vnútorné uhly \( \alpha \), \( \beta \), \( \gamma \). Vypočítajte veľkosť  uhla  \( \alpha \) ak \( a^2 = b^2 + c^2 +bc \).}
{\( 120^{\circ} \)}{\( 60^{\circ} \)}{\( 30^{\circ} \)}{\( 90^{\circ} \)}{}{}}
\LANGes{
\Question{
En un triángulo, las longitudes de los lados son \( a \), \( b \), \( c \) y los ángulos opuestos son \( \alpha \), \( \beta \), \( \gamma \). Calcula la medida del ángulo \( \alpha \) si \( a^2 = b^2 + c^2 +bc \).}
{\( 120^{\circ} \)}{\( 60^{\circ} \)}{\( 30^{\circ} \)}{\( 90^{\circ} \)}{}{}}
\End

\Data\ID{1003076909}
\Updated{2020-07-15 03:07:21}
\Level{B}
\SubArea{Triangles}
\LANGen{
\Question{
\( ABC \) is a triangle. Given \( |AB|=3\,\mathrm{cm} \), the measure of \( \measuredangle CAB \) is \( 75^{\circ} \), and the measure of \(\measuredangle ABC \) is \( 45^{\circ} \), calculate the length of the side \( AC \).}
{\( \sqrt6\,\mathrm{cm} \)}{\( 3\sqrt2\,\mathrm{cm} \)}{\( 2\sqrt3\,\mathrm{cm} \)}{\( 3\frac{\sqrt3}{\sqrt2}\,\mathrm{cm} \)}{}{}}
\LANGpl{
\Question{
\( ABC \) jest trójkątem. Podano \( |AB|=3\,\mathrm{cm} \), miara kąta \( CAB \) jest \( 75^{\circ} \), miara kąta \( ABC \) jest równa \( 45^{\circ} \). Oblicz długość boku\( AC \).}
{\( \sqrt6\,\mathrm{cm} \)}{\( 3\sqrt2\,\mathrm{cm} \)}{\( 2\sqrt3\,\mathrm{cm} \)}{\( 3\frac{\sqrt3}{\sqrt2}\,\mathrm{cm} \)}{}{}}
\LANGcs{
\Question{
V trojúhelníku \( ABC \) je \( |AB|=3\,\mathrm{cm} \), velikost \( \measuredangle CAB \) je \( 75^{\circ} \) a velikost \(\measuredangle ABC \) je \( 45^{\circ} \). Vypočítejte délku strany \( AC \).}
{\( \sqrt6\,\mathrm{cm} \)}{\( 3\sqrt2\,\mathrm{cm} \)}{\( 2\sqrt3\,\mathrm{cm} \)}{\( 3\frac{\sqrt3}{\sqrt2}\,\mathrm{cm} \)}{}{}}
\LANGsk{
\Question{
V trojuholníku \( ABC \) je \( |AB|=3\,\mathrm{cm} \), veľkosť \( \measuredangle CAB \) je \( 75^{\circ} \) a veľkosť   \(\measuredangle ABC \) je \( 45^{\circ} \). Vypočítajte dĺžku strany  \( AC \).}
{\( \sqrt6\,\mathrm{cm} \)}{\( 3\sqrt2\,\mathrm{cm} \)}{\( 2\sqrt3\,\mathrm{cm} \)}{\( 3\frac{\sqrt3}{\sqrt2}\,\mathrm{cm} \)}{}{}}
\LANGes{
\Question{
Dado el triángulo \( ABC \) con \( |AB|=3\,\mathrm{cm} \), la medida de \( \measuredangle CAB \) es \( 75^{\circ} \), y la medida de \(\measuredangle ABC \) es \( 45^{\circ} \), calcula la longitud del lado \( AC \).}
{\( \sqrt6\,\mathrm{cm} \)}{\( 3\sqrt2\,\mathrm{cm} \)}{\( 2\sqrt3\,\mathrm{cm} \)}{\( 3\frac{\sqrt3}{\sqrt2}\,\mathrm{cm} \)}{}{}}
\End

\Data\ID{1003076910}
\Updated{2022-04-23 05:04:05}
\Level{B}
\SubArea{Triangles}
\LANGen{
\Question{
The lengths of the sides of a triangle are \( 4\,\mathrm{cm} \), \( 6\,\mathrm{cm} \) and \( 8\,\mathrm{cm} \). Calculate cosine of its smallest interior angle.}
{\( \frac78 \)}{\( 28.96^{\circ} \)}{\( -\frac14 \)}{\( -\frac78 \)}{}{}}
\LANGpl{
\Question{
Długości boków trójkąta wynoszą \( 4\,\mathrm{cm} \), \( 6\,\mathrm{cm} \) i \( 8\,\mathrm{cm} \). Oblicz cosinus jego najmniejszego kąta wewnętrznego.}
{\( \frac78 \)}{\( 28{,}96^{\circ} \)}{\( -\frac14 \)}{\( -\frac78 \)}{}{}}
\LANGcs{
\Question{
Délky stran trojúhelníku jsou \( 4\,\mathrm{cm} \), \( 6\,\mathrm{cm} \) a \( 8\,\mathrm{cm} \). Vypočítejte kosinus jeho nejmenšího vnitřního úhlu.}
{\( \frac78 \)}{\( 28{,}96^{\circ} \)}{\( -\frac14 \)}{\( -\frac78 \)}{}{}}
\LANGsk{
\Question{
Dĺžky strán trojuholníka sú \( 4\,\mathrm{cm} \), \( 6\,\mathrm{cm} \) a \( 8\,\mathrm{cm} \). Vypočítajte kosínus jeho najmenšieho vnútorného uhla.}
{\( \frac78 \)}{\( 28{,}96^{\circ} \)}{\( -\frac14 \)}{\( -\frac78 \)}{}{}}
\LANGes{
\Question{
Las longitudes de un triángulo son \( 4\,\mathrm{cm} \), \( 6\,\mathrm{cm} \) y \( 8\,\mathrm{cm} \). Calcula el valor del coseno de su ángulo interior más pequeño.}
{\( \frac78 \)}{\( 28.96^{\circ} \)}{\( -\frac14 \)}{\( -\frac78 \)}{}{}}
\End

\Data\ID{1003077001}
\Updated{2021-04-19 06:04:18}
\Level{B}
\SubArea{Triangles}
\LANGen{
\Question{
A triangle has sides whose lengths are \( 6\,\mathrm{cm} \), \( 7\,\mathrm{cm} \) and \( 9\,\mathrm{cm} \). Find cosine of its largest interior angle.}
{\( \frac1{21} \)}{\( 87.27^{\circ} \)}{\( 1.98 \)}{\( -\frac1{21} \)}{}{}}
\LANGpl{
\Question{
Dany jest trójkąt o bokach \( 6\,\mathrm{cm} \), \( 7\,\mathrm{cm} \) i \( 9\,\mathrm{cm} \). Oblicz cosinus jego największego kąta.}
{\( \frac1{21} \)}{\( 87{,}27^{\circ} \)}{\( 1{,}98 \)}{\( -\frac1{21} \)}{}{}}
\LANGcs{
\Question{
Je dán trojúhelník s délkami stran \( 6\,\mathrm{cm} \), \( 7\,\mathrm{cm} \) a \( 9\,\mathrm{cm} \). Vypočítejte kosinus jeho největšího vnitřního úhlu.}
{\( \frac1{21} \)}{\( 87{,}27^{\circ} \)}{\( 1{,}98 \)}{\( -\frac1{21} \)}{}{}}
\LANGsk{
\Question{
Daný je trojuholník s dĺžkami strán  \( 6\,\mathrm{cm} \), \( 7\,\mathrm{cm} \) a \( 9\,\mathrm{cm} \). Vypočítajte kosínus jeho najväčšieho vnútorného uhla.}
{\( \frac1{21} \)}{\( 87{,}27^{\circ} \)}{\( 1{,}98 \)}{\( -\frac1{21} \)}{}{}}
\LANGes{
\Question{
Dado el triángulo cuyas longitudes de los lados son \( 6\,\mathrm{cm} \), \( 7\,\mathrm{cm} \) y \( 9\,\mathrm{cm} \). Halla el valor del coseno de su ángulo interior más grande.}
{\( \frac1{21} \)}{\( 87.27^{\circ} \)}{\( 1.98 \)}{\( -\frac1{21} \)}{}{}}
\End

\Data\ID{1003077002}
\Updated{2020-07-15 03:07:21}
\Level{B}
\SubArea{Triangles}
\LANGen{
\Question{
Find the measure of the largest interior angle of a triangle whose sides are \( 3\,\mathrm{cm} \), \( 4\,\mathrm{cm} \) and \( 6\,\mathrm{cm} \) long.}
{\( 117.28^{\circ} \)}{\( 62.72^{\circ} \)}{\( 143.66^{\circ} \)}{\( 36.34^{\circ} \)}{}{}}
\LANGpl{
\Question{
Wyznacz miarę największego kąta wewnętrznego trójkąta o bokach długości \( 3\,\mathrm{cm} \), \( 4\,\mathrm{cm} \) i \( 6\,\mathrm{cm} \).}
{\( 117{,}28^{\circ} \)}{\( 62{,}72^{\circ} \)}{\( 143{,}66^{\circ} \)}{\( 36{,}34^{\circ} \)}{}{}}
\LANGcs{
\Question{
Vypočítejte velikost největšího vnitřního úhlu trojúhelníku, jehož strany mají délky \( 3\,\mathrm{cm} \), \( 4\,\mathrm{cm} \) a \( 6\,\mathrm{cm} \).}
{\( 117{,}28^{\circ} \)}{\( 62{,}72^{\circ} \)}{\( 143{,}66^{\circ} \)}{\( 36{,}34^{\circ} \)}{}{}}
\LANGsk{
\Question{
Vypočítajte veľkosť najväčšieho vnútorného uhla trojuholníka, ktorého strany majú dĺžky \( 3\,\mathrm{cm} \), \( 4\,\mathrm{cm} \) a \( 6\,\mathrm{cm} \).}
{\( 117{,}28^{\circ} \)}{\( 62{,}72^{\circ} \)}{\( 143{,}66^{\circ} \)}{\( 36{,}34^{\circ} \)}{}{}}
\LANGes{
\Question{
Calcula la medida del ángulo interior más grande de un triángulo cuyos lados miden \( 3\,\mathrm{cm} \), \( 4\,\mathrm{cm} \) y \( 6\,\mathrm{cm} \).}
{\( 117.28^{\circ} \)}{\( 62.72^{\circ} \)}{\( 143.66^{\circ} \)}{\( 36.34^{\circ} \)}{}{}}
\End

\Data\ID{1003077006}
\Updated{2020-07-15 03:07:21}
\Level{B}
\SubArea{Triangles}
\LANGen{
\Question{
Giving a right-angled triangle, the hypotenuse is \( 50\,\mathrm{cm} \) long, the perimeter of the triangle is \( 12\,\mathrm{dm} \) and its area is \( 600\,\mathrm{cm}^2 \). Find the measures of all interior angles of the triangle.}
{\( 90^{\circ};\ 36.87^{\circ};\ 53.13^{\circ} \)}{\( 90^{\circ};\ 30.96^{\circ};\ 59.04^{\circ} \)}{\( 90^{\circ};\ 38.65^{\circ};\ 51.35^{\circ} \)}{\( 90^{\circ};\ 33.13^{\circ};\ 56.87^{\circ} \)}{}{}}
\LANGpl{
\Question{
Dany jest trójkąt prostokątny o przeciwprostokątnej długości \( 50\,\mathrm{cm} \) i obwodzie  \( 12\,\mathrm{dm} \), a jego pole jest \( 600\,\mathrm{cm}^2 \). Wyznacz miarę wszystkich kątów wewnętrznych tego trójkąta.}
{\( 90^{\circ};\ 36{,}87^{\circ};\ 53{,}13^{\circ} \)}{\( 90^{\circ};\ 30{,}96^{\circ};\ 59{,}04^{\circ} \)}{\( 90^{\circ};\ 38{,}65^{\circ};\ 51{,}35^{\circ} \)}{\( 90^{\circ};\ 33{,}13^{\circ};\ 56{,}87^{\circ} \)}{}{}}
\LANGcs{
\Question{
V pravoúhlém trojúhelníku má přepona délku \( 50\,\mathrm{cm} \), obvod tohoto trojúhelníku je \( 12\,\mathrm{dm} \) a obsah \( 600\,\mathrm{cm}^2 \). Najděte velikost všech vnitřních úhlů trojúhelníku.}
{\( 90^{\circ};\ 36{,}87^{\circ};\ 53{,}13^{\circ} \)}{\( 90^{\circ};\ 30{,}96^{\circ};\ 59{,}04^{\circ} \)}{\( 90^{\circ};\ 38{,}65^{\circ};\ 51{,}35^{\circ} \)}{\( 90^{\circ};\ 33{,}13^{\circ};\ 56{,}87^{\circ} \)}{}{}}
\LANGsk{
\Question{
V pravouhlom trojuholníku  má prepona dĺžku \( 50\,\mathrm{cm} \), obvod tohto trojuholníka je \( 12\,\mathrm{dm} \) a obsah \( 600\,\mathrm{cm}^2 \). Nájdite veľkosti  všetkých vnútorných uhlov trojuholníka.}
{\( 90^{\circ};\ 36{,}87^{\circ};\ 53{,}13^{\circ} \)}{\( 90^{\circ};\ 30{,}96^{\circ};\ 59{,}04^{\circ} \)}{\( 90^{\circ};\ 38{,}65^{\circ};\ 51{,}35^{\circ} \)}{\( 90^{\circ};\ 33{,}13^{\circ};\ 56{,}87^{\circ} \)}{}{}}
\LANGes{
\Question{
Dado un triángulo rectángulo. La hipotenusa mide \( 50\,\mathrm{cm} \), el perímetro del triángulo mide \( 12\,\mathrm{dm} \) y su área es \( 600\,\mathrm{cm}^2 \). Calcula las medidas de todos los ángulos interiores.}
{\( 90^{\circ};\ 36.87^{\circ};\ 53.13^{\circ} \)}{\( 90^{\circ};\ 30.96^{\circ};\ 59.04^{\circ} \)}{\( 90^{\circ};\ 38.65^{\circ};\ 51.35^{\circ} \)}{\( 90^{\circ};\ 33.13^{\circ};\ 56.87^{\circ} \)}{}{}}
\End

\Data\ID{1003077010}
\Updated{2020-07-15 03:07:21}
\Level{B}
\SubArea{Triangles}
\LANGen{
\Question{
In an isosceles triangle \( ABC \) the base \( AB \) has length \( 12\,\mathrm{cm} \). The altitude to the base \( v_c=8\,\mathrm{cm} \). Determine the length of the median drawn from a vertex at the base to the side.}
{\( \sqrt{97}\,\mathrm{cm} \)}{\( \sqrt{93}\,\mathrm{cm} \)}{\( \sqrt{87}\,\mathrm{cm} \)}{\( \sqrt{83}\,\mathrm{cm} \)}{}{}}
\LANGpl{
\Question{
W trójkącie równoramiennym  \( ABC \) podstawa \( AB \) ma długość \( 12\,\mathrm{cm} \). Wysokość do podstawy \( v_c=8\,\mathrm{cm} \). 
Określ długość środkowej względem ramienia.}
{\( \sqrt{97}\,\mathrm{cm} \)}{\( \sqrt{93}\,\mathrm{cm} \)}{\( \sqrt{87}\,\mathrm{cm} \)}{\( \sqrt{83}\,\mathrm{cm} \)}{}{}}
\LANGcs{
\Question{
Rovnoramenný trojúhelník  \( ABC \) má základnu \( AB \) dlouhou \( 12\,\mathrm{cm} \). Výška na základnu \( v_c=8\,\mathrm{cm} \). Vypočítejte délku těžnice sestrojené na rameno trojúhelníku.}
{\( \sqrt{97}\,\mathrm{cm} \)}{\( \sqrt{93}\,\mathrm{cm} \)}{\( \sqrt{87}\,\mathrm{cm} \)}{\( \sqrt{83}\,\mathrm{cm} \)}{}{}}
\LANGsk{
\Question{
Rovnoramenný trojuholník  \( ABC \) má základňu  \( AB \) dlhú \( 12\,\mathrm{cm} \). Výška na základňu  \( v_c=8\,\mathrm{cm} \). Vypočítajte dĺžku ťažnice zostrojenej na rameno trojuholníka.}
{\( \sqrt{97}\,\mathrm{cm} \)}{\( \sqrt{93}\,\mathrm{cm} \)}{\( \sqrt{87}\,\mathrm{cm} \)}{\( \sqrt{83}\,\mathrm{cm} \)}{}{}}
\LANGes{
\Question{
La base \( AB \) del triángulo isósceles \( ABC \) mide \( 12\,\mathrm{cm} \). La altura sobre la base mide \( v_c=8\,\mathrm{cm} \). Calcula la longitud de la mediana dibujada desde un vértice de la base hacia un lado.}
{\( \sqrt{97}\,\mathrm{cm} \)}{\( \sqrt{93}\,\mathrm{cm} \)}{\( \sqrt{87}\,\mathrm{cm} \)}{\( \sqrt{83}\,\mathrm{cm} \)}{}{}}
\End

\Data\ID{1103021801}
\Updated{2022-04-23 05:04:05}
\Level{B}
\SubArea{Triangles}
\IMG{1103021801.pdf}
\LANGen{
\Question{
Given a right-angled triangle \( ABC \) with the right angle at vertex \( C \), calculate \( \sin\beta \) if \(\cos\alpha=\frac5{13} \).}
{\( \frac5{13} \)}{\( \frac{13}5 \)}{\( \frac{12}{13} \)}{\( \frac5{12} \)}{}{}}
\LANGpl{
\Question{
Dany jest trójkąt prostokątny  \( ABC \) o kącie prostym w wierzchołku \( C \), oblicz \( \sin\beta \) jeśli \(\cos\alpha=\frac5{13} \).}
{\( \frac5{13} \)}{\( \frac{13}5 \)}{\( \frac{12}{13} \)}{\( \frac5{12} \)}{}{}}
\LANGcs{
\Question{
Je dán pravoúhlý trojúhelník \( ABC \) s pravým úhlem u vrcholu \( C \). Vypočítejte \( \sin\beta \), jestliže \(\cos\alpha=\frac5{13} \).}
{\( \frac5{13} \)}{\( \frac{13}5 \)}{\( \frac{12}{13} \)}{\( \frac5{12} \)}{}{}}
\LANGsk{
\Question{
Daný je pravouhlý trojuholník  \( ABC \) s pravým uhlom pri vrchole \( C \). Vypočítajte \( \sin\beta \) ak \(\cos\alpha=\frac5{13} \).}
{\( \frac5{13} \)}{\( \frac{13}5 \)}{\( \frac{12}{13} \)}{\( \frac5{12} \)}{}{}}
\LANGes{
\Question{
Dado el triángulo rectángulo \( ABC \), siendo \( C \) el vértice del ángulo recto, calcula \( \sin\beta \) suponiendo que \(\cos\alpha=\frac5{13} \).}
{\( \frac5{13} \)}{\( \frac{13}5 \)}{\( \frac{12}{13} \)}{\( \frac5{12} \)}{}{}}
\End

\Data\ID{1103021802}
\Updated{2020-07-15 03:07:21}
\Level{B}
\SubArea{Triangles}
\IMG{1103021802.pdf}
\LANGen{
\Question{
The arms of a double ladder are \( 150\,\mathrm{cm} \) long. After opening the ladder (look at the picture) the arms contain an angle of \( 40^{\circ} \). Calculate the height of the opened ladder (the distance between the highest point of the ladder and the ground). Round the result to the nearest integer.}
{\( 141\,\mathrm{cm} \)}{\( 115\,\mathrm{cm} \)}{\( 51\,\mathrm{cm} \)}{\( 96\,\mathrm{cm} \)}{}{}}
\LANGpl{
\Question{
Ramiona podwójnej drabiny mają  \( 150\,\mathrm{cm} \) długości. Po rozłożeniu drabiny (patrz rysunek) jej ramiona tworzą kąt \( 40^{\circ} \). Oszacuj wysokość rozłożonej drabiny (odległość pomiędzy najwyższym punktem drabiny a ziemią). Wynik zaokrąglij do najbliższej liczby całkowitej.}
{\( 141\,\mathrm{cm} \)}{\( 115\,\mathrm{cm} \)}{\( 51\,\mathrm{cm} \)}{\( 96\,\mathrm{cm} \)}{}{}}
\LANGcs{
\Question{
Ramena dvojitého žebříku mají délku \( 150\,\mathrm{cm} \). Po rozevření žebříku (viz obrázek) ramena svírají úhel \( 40^{\circ} \). Určete zaokrouhleně na celé centimetry výšku takto rozevřeného žebříku (vzdálenost nejvyšších bodů žebříku od podlahy).}
{\( 141\,\mathrm{cm} \)}{\( 115\,\mathrm{cm} \)}{\( 51\,\mathrm{cm} \)}{\( 96\,\mathrm{cm} \)}{}{}}
\LANGsk{
\Question{
Ramená dvojitého rebríka majú dĺžku  \( 150\,\mathrm{cm} \). Po roztvorení rebríka (pozrite obrázok) ramená zvierajú uhol \( 40^{\circ} \). Určte zaokrúhlene na celé centimetre výšku takto roztvoreného rebríka (vzdialenosť najvyšších bodov rebríka od podlahy).}
{\( 141\,\mathrm{cm} \)}{\( 115\,\mathrm{cm} \)}{\( 51\,\mathrm{cm} \)}{\( 96\,\mathrm{cm} \)}{}{}}
\LANGes{
\Question{
Los lados de una escalera doble miden \( 150\,\mathrm{cm} \). Después de abrir la escalera (mira la imagen), los lados forman un ángulo de \( 40^{\circ} \). Calcula la altura de la escalera abierta (la distancia entre el punto más alto de la escalera y el suelo). Redondea el resultado al entero más cercano.}
{\( 141\,\mathrm{cm} \)}{\( 115\,\mathrm{cm} \)}{\( 51\,\mathrm{cm} \)}{\( 96\,\mathrm{cm} \)}{}{}}
\End

\Data\ID{1103021804}
\Updated{2021-04-19 06:04:10}
\Level{B}
\SubArea{Triangles}
\IMG{1103021804.pdf}
\LANGen{
\Question{
The roof shield has the shape of an isosceles triangle. The width of the shield is \( 12\,\mathrm{m} \) and the slope of the roof is \( 38^{\circ} \). Calculate the height of the shield. Round the result to two decimal places.}
{\( 4.69\,\mathrm{m} \)}{\( 7.39\,\mathrm{m} \)}{\( 9.46\,\mathrm{m} \)}{\( 3.70\,\mathrm{m} \)}{}{}}
\LANGpl{
\Question{
Szczyt dachu ma kształt trójkąta równoramiennego. Szerokość szczytu wynosi \( 12\,\mathrm{m} \), i nachylenie dachu jest \( 38^{\circ} \). Oblicz wysokość szczychu. Zaokrąglij wynik do dwóch miejsc dziesiętnych.}
{\( 4{,}69\,\mathrm{m} \)}{\( 7{,}39\,\mathrm{m} \)}{\( 9{,}46\,\mathrm{m} \)}{\( 3{,}70\,\mathrm{m} \)}{}{}}
\LANGcs{
\Question{
Štít střechy má tvar rovnoramenného trojúhelníku. Šířka štítu je  \( 12\,\mathrm{m} \) a sklon střechy je  \( 38^{\circ} \). Vypočítejte výšku štítu. Výsledek zaokrouhlete na dvě desetinná místa.}
{\( 4{,}69\,\mathrm{m} \)}{\( 7{,}39\,\mathrm{m} \)}{\( 9{,}46\,\mathrm{m} \)}{\( 3{,}70\,\mathrm{m} \)}{}{}}
\LANGsk{
\Question{
Štít strechy má tvar rovnoramenného trojuholníka.  Šírka štítu je  \( 12\,\mathrm{m} \) a sklon strechy je  \( 38^{\circ} \). Vypočítajte výšku štítu. Výsledok zaokrúhlite na dve desatinné miesta.}
{\( 4{,}69\,\mathrm{m} \)}{\( 7{,}39\,\mathrm{m} \)}{\( 9{,}46\,\mathrm{m} \)}{\( 3{,}70\,\mathrm{m} \)}{}{}}
\LANGes{
\Question{
La buhardilla del techo tiene forma del triángulo isósceles. La anchura de la buhardilla mide \( 12\,\mathrm{m} \) y la pendiente del techo es \( 38^{\circ} \). Calcula la altura de la buhardilla. Redondea el resultado a dos decimales.}
{\( 4.69\,\mathrm{m} \)}{\( 7.39\,\mathrm{m} \)}{\( 9.46\,\mathrm{m} \)}{\( 3.70\,\mathrm{m} \)}{}{}}
\End

\Data\ID{1103021806}
\Updated{2021-04-19 06:04:30}
\Level{B}
\SubArea{Triangles}
\IMG{1103021806.pdf}
\LANGen{
\Question{
The figure shows a water tower. From the site \( 85 \) meters away at a height of \( 1.2 \) meters, the measured angle of elevation of the top of the tower is \( 20^{\circ}30' \).  Find the height of the tower. Round the result to two decimal places.}
{\( 32.98\,\mathrm{m} \)}{\( 31.78\,\mathrm{m} \)}{\( 31.44\,\mathrm{m} \)}{\( 32.64\,\mathrm{m} \)}{}{}}
\LANGpl{
\Question{
Na rysunku pokazano wieżę ciśnień.  Z odległości \( 85 \) metrów na wysokości \( 1{,}2 \) metrów jest mierzony kąt podniesienia wierzchołka \( 20^{\circ}30' \).  Oblicz wysokość wieży. Zaokrąglij do dwóch miejsc dziesiętnych.}
{\( 32{,}98\,\mathrm{m} \)}{\( 31{,}78\,\mathrm{m} \)}{\( 31{,}44\,\mathrm{m} \)}{\( 32{,}64\,\mathrm{m} \)}{}{}}
\LANGcs{
\Question{
Vypočítejte výšku vodárenské věže, jejíž výška byla měřená přístrojem vysokým  \( 1{,}2 \) metru. Měřící přístroj je od paty věže vzdálený \( 85 \) metrů a naměřil výškový úhel \( 20^{\circ}30' \). Výsledek zaokrouhlete na dvě desetinná místa.}
{\( 32{,}98\,\mathrm{m} \)}{\( 31{,}78\,\mathrm{m} \)}{\( 31{,}44\,\mathrm{m} \)}{\( 32{,}64\,\mathrm{m} \)}{}{}}
\LANGsk{
\Question{
Vypočítajte výšku vodárenskej veže, ktorej výška bola meraná meracím prístrojom vysokým  \( 1{,}2 \) metra. Merací prístroj je od päty veže vzdialený \( 85 \) metrov a nameral výškový uhol \( 20^{\circ}30' \). Výsledok zaokrúhlite na dve desatinné miesta.}
{\( 32{,}98\,\mathrm{m} \)}{\( 31{,}78\,\mathrm{m} \)}{\( 31{,}44\,\mathrm{m} \)}{\( 32{,}64\,\mathrm{m} \)}{}{}}
\LANGes{
\Question{
La imagen representa una torre. Desde un sitio a \( 85 \) metros de distancia y a una altura de \( 1.2 \) metros, el ángulo de elevación respecto a la parte superior de la torre es \( 20^{\circ}30' \). Calcula la altura de la torre. Redondea el resultado a dos decimales.}
{\( 32.98\,\mathrm{m} \)}{\( 31.78\,\mathrm{m} \)}{\( 31.44\,\mathrm{m} \)}{\( 32.64\,\mathrm{m} \)}{}{}}
\End

\Data\ID{1103021807}
\Updated{2022-04-23 05:04:05}
\Level{B}
\SubArea{Triangles}
\IMG{1103021807.pdf}
\LANGen{
\Question{
An artillery battery is placed on a cliff. From the edge of the cliff \( 200\,\mathrm{m} \) high, the angle of depression to the ship on the sea is \( 10^{\circ} \). What is the distance \( d \) (see the picture) from the battery to the ship? Round the result to two decimal places.}
{\( 1151.75\,\mathrm{m} \)}{\( 203.09\,\mathrm{m} \)}{\( 35.27\,\mathrm{m} \)}{\( 1134.26\,\mathrm{m} \)}{}{}}
\LANGpl{
\Question{
Bateria artyleryjskia znajduje się na klifie. Od krawędzi urwiska \( 200\,\mathrm{m} \) wysokości kąt depresji do statku na morzu wynosi \( 10^{\circ} \). Jaka jest odległość \( d \) baterii od statku? Zaokrąglij wynik do dwóch miejsc dziesiętnych.}
{\( 1151{,}75\,\mathrm{m} \)}{\( 203{,}09\,\mathrm{m} \)}{\( 35{,}27\,\mathrm{m} \)}{\( 1134{,}26\,\mathrm{m} \)}{}{}}
\LANGcs{
\Question{
Dělostřelecká baterie je umístěná na útesu vysokém \( 200\,\mathrm{m} \). Určete vzdálenost \( d \) baterie od lodě, kterou z útesu pozorujeme v hloubkovém úhlu \( 10^{\circ} \). Výsledek zaokrouhlete na dvě desetinná místa.}
{\( 1151{,}75\,\mathrm{m} \)}{\( 203{,}09\,\mathrm{m} \)}{\( 35{,}27\,\mathrm{m} \)}{\( 1134{,}26\,\mathrm{m} \)}{}{}}
\LANGsk{
\Question{
Delostrelecká batéria je umiestnená na útese vysokom  \( 200\,\mathrm{m} \). Určte vzdialenosť  \( d \) batérie od lode, ktorú z útesu pozorujeme v hĺbkovom uhle \( 10^{\circ} \). Výsledok zaokrúhlite na dve desatinné miesta.}
{\( 1151{,}75\,\mathrm{m} \)}{\( 203{,}09\,\mathrm{m} \)}{\( 35{,}27\,\mathrm{m} \)}{\( 1134{,}26\,\mathrm{m} \)}{}{}}
\LANGes{
\Question{
Una batería de artillería está situada en un acantilado de \( 200\,\mathrm{m} \) de altitud. ¿Cuál es la distancia \( d \) de la batería al barco, observado desde el acantilado bajo el ángulo de depresión de \( 10^{\circ} \)? Redondea el resultado a dos decimales.}
{\( 1151.75\,\mathrm{m} \)}{\( 203.09\,\mathrm{m} \)}{\( 35.27\,\mathrm{m} \)}{\( 1134.26\,\mathrm{m} \)}{}{}}
\End

\Data\ID{1103021808}
\Updated{2021-04-19 06:04:03}
\Level{B}
\SubArea{Triangles}
\IMG{1103021808.pdf}
\LANGen{
\Question{
There is a cottage on the top of the mountain. From our site \( P \), \( 2\,\mathrm{km} \) as the crow flies from the cottage, we can observe the cottage to be at an angle of elevation of \( 30^{\circ} \). How many altitude meters do we still have to overcome to get to the cottage?}
{\( 1000\,\mathrm{m} \)}{\( 1732\,\mathrm{m} \)}{\( 2\,\mathrm{km} \)}{\( 1155\,\mathrm{m} \)}{}{}}
\LANGpl{
\Question{
Na szczycie góry znajduje się chatka. Z naszej pozycji \( P \) \( 2\,\mathrm{km} \) w linii powietrznej od chatki możemy obserwować chatkę pod kątem \( 30^{\circ} \). Ile metrów trzeba jeszcze pokonać, aby dostać się do chatki?}
{\( 1000\,\mathrm{m} \)}{\( 1732\,\mathrm{m} \)}{\( 2\,\mathrm{km} \)}{\( 1155\,\mathrm{m} \)}{}{}}
\LANGcs{
\Question{
Na vrcholku hory stojí chata. Z našeho místa \( P \) vzdáleného \( 2\,\mathrm{km} \) vzdušnou čarou od chaty vidíme chatu pod výškovým úhlem \( 30^{\circ} \). Kolik výškových metrů musíme překonat, abychom se dostali k chatě?}
{\( 1000\,\mathrm{m} \)}{\( 1732\,\mathrm{m} \)}{\( 2\,\mathrm{km} \)}{\( 1155\,\mathrm{m} \)}{}{}}
\LANGsk{
\Question{
Na vrchole hory je chata. Táto chata je podľa mapy, vzdušnou čiarou, od našej aktuálnej polohy \( P \) vzdialená  \( 2\,\mathrm{km} \) a vidíme ju pod výškovým uhlom  \( 30^{\circ} \). Koľko výškových metrov ešte musíme prekonať, ak sa chceme dostať k chate?}
{\( 1000\,\mathrm{m} \)}{\( 1732\,\mathrm{m} \)}{\( 2\,\mathrm{km} \)}{\( 1155\,\mathrm{m} \)}{}{}}
\LANGes{
\Question{
En la cima de una montaña hay una cabaña. Desde nuestro sitio \( P \), que está a \( 2\,\mathrm{km} \) en línea recta, podemos observar la cabaña bajo un ángulo de elevación de \( 30^{\circ} \). ¿Cuántos metros de altitud tenemos que superar para llegar a la cabaña?}
{\( 1000\,\mathrm{m} \)}{\( 1732\,\mathrm{m} \)}{\( 2\,\mathrm{km} \)}{\( 1155\,\mathrm{m} \)}{}{}}
\End

\Data\ID{1103021810}
\Updated{2021-04-19 06:04:40}
\Level{B}
\SubArea{Triangles}
\IMG{1103021810.pdf}
\LANGen{
\Question{
What is the height difference between the two stations of the cableway if its slope is \( 30^{\circ} \) and the distance between the stations is \( 1500\,\mathrm{m} \)?}
{\( 750\,\mathrm{m} \)}{\( 1299\,\mathrm{m} \)}{\( 866\,\mathrm{m} \)}{\( 890\,\mathrm{m} \)}{}{}}
\LANGpl{
\Question{
Jaka jest różnica wysokości między dwiema stacjami kolejki linowej, jeśli jej kąt nachylenia jest \( 30^{\circ} \) i odległość pomiędzy obydwiema stacjami wynosi \( 1500\,\mathrm{m} \)?}
{\( 750\,\mathrm{m} \)}{\( 1299\,\mathrm{m} \)}{\( 866\,\mathrm{m} \)}{\( 890\,\mathrm{m} \)}{}{}}
\LANGcs{
\Question{
Jaký je výškový rozdíl mezi dvěma stanicemi lanovky, jestliže její sklon je  \( 30^{\circ} \) a vzdálenost mezi stanicemi je \( 1500\,\mathrm{m} \)?}
{\( 750\,\mathrm{m} \)}{\( 1299\,\mathrm{m} \)}{\( 866\,\mathrm{m} \)}{\( 890\,\mathrm{m} \)}{}{}}
\LANGsk{
\Question{
Aký je výškový rozdiel dvoch staníc lanovky, keď jej stúpanie je  \( 30^{\circ} \) a vzdialenosť staníc je \( 1500\,\mathrm{m} \)?}
{\( 750\,\mathrm{m} \)}{\( 1299\,\mathrm{m} \)}{\( 866\,\mathrm{m} \)}{\( 890\,\mathrm{m} \)}{}{}}
\LANGes{
\Question{
¿Cuál es la diferencia de altitud entre dos estaciones de un teleférico si su pendiente es \( 30^{\circ} \) y la distancia entre las dos estaciones es de \( 1500\,\mathrm{m} \)?}
{\( 750\,\mathrm{m} \)}{\( 1299\,\mathrm{m} \)}{\( 866\,\mathrm{m} \)}{\( 890\,\mathrm{m} \)}{}{}}
\End

\Data\ID{1103076809}
\Updated{2021-04-19 06:04:42}
\Level{B}
\SubArea{Triangles}
\IMG{1103076809.pdf}
\LANGen{
\Question{
The diagram shows a square inscribed into an equilateral triangle with the side \( 4\,\mathrm{cm} \) long. Calculate the length of the side of the square. Round the result to two decimal places.}
{\( 1.86\,\mathrm{cm} \)}{\( 2.14\,\mathrm{cm} \)}{\( 3.12\,\mathrm{cm} \)}{\( 4.61\,\mathrm{cm} \)}{}{}}
\LANGpl{
\Question{
Rysunek przedstawia kwadrat wpisany w trójkąt równoboczny o boku długości \( 4\,\mathrm{cm} \). Oblicz długość boku kwadratu. Zaokrąglij wynik do drugiego miejsca po przecinku.}
{\( 1{,}86\,\mathrm{cm} \)}{\( 2{,}14\,\mathrm{cm} \)}{\( 3{,}12\,\mathrm{cm} \)}{\( 4{,}61\,\mathrm{cm} \)}{}{}}
\LANGcs{
\Question{
Do rovnostranného trojúhelníku se stranou dlouhou \( 4\,\mathrm{cm} \) je vepsaný čtverec. Vypočítejte délku strany tohoto čtverce. Výsledek uveďte s přesností na dvě desetinná místa.}
{\( 1{,}86\,\mathrm{cm} \)}{\( 2{,}14\,\mathrm{cm} \)}{\( 3{,}12\,\mathrm{cm} \)}{\( 4{,}61\,\mathrm{cm} \)}{}{}}
\LANGsk{
\Question{
Do rovnostranného trojuholníka so stranou dlhou  \( 4\,\mathrm{cm} \) je vpísaný štvorec. Vypočítajte dĺžku strany tohto štvorca. Výsledok uveďte s presnosťou na dve desatinné miesta.}
{\( 1{,}86\,\mathrm{cm} \)}{\( 2{,}14\,\mathrm{cm} \)}{\( 3{,}12\,\mathrm{cm} \)}{\( 4{,}61\,\mathrm{cm} \)}{}{}}
\LANGes{
\Question{
La imagen representa un cuadrado inscrito en un triángulo equilátero cuyo lado que mide \( 4\,\mathrm{cm} \). Calcula la longitud del lado del cuadrado. Redondea el resultado a dos decimales.}
{\( 1.86\,\mathrm{cm} \)}{\( 2.14\,\mathrm{cm} \)}{\( 3.12\,\mathrm{cm} \)}{\( 4.61\,\mathrm{cm} \)}{}{}}
\End

\Data\ID{1103076901}
\Updated{2020-07-15 03:07:21}
\Level{B}
\SubArea{Triangles}
\IMG{1103076901_0.pdf}
\LANGen{
\Question{
Given a right-angled triangle \( ABC \) with the right angle at vertex \( C \), calculate \( \mathrm{tg}\,\beta \) if \( \cos\alpha=\frac5{13} \).}
{\( \frac5{12} \)}{\( \frac5{13} \)}{\( \frac{13}5 \)}{\( \frac{12}{13} \)}{}{}}
\LANGpl{
\Question{
Dany jest trójkąt prostokątny \( ABC \) o kącie prostym w wierzchołku  \( C \), oblicz \( \mathrm{tg}\,\beta \) jeśli \( \cos\alpha=\frac5{13} \).}
{\( \frac5{12} \)}{\( \frac5{13} \)}{\( \frac{13}5 \)}{\( \frac{12}{13} \)}{}{}}
\LANGcs{
\Question{
Je dán pravoúhlý trojúhelník \( ABC \) s pravým úhlem u vrcholu \( C \). Vypočítejte \( \mathrm{tg}\,\beta \), jestliže \( \cos\alpha=\frac5{13} \).}
{\( \frac5{12} \)}{\( \frac5{13} \)}{\( \frac{13}5 \)}{\( \frac{12}{13} \)}{}{}}
\LANGsk{
\Question{
Daný je pravouhlý trojuholník  \( ABC \) s pravým uhlom pri vrchole  \( C \). Vypočítajte \( \mathrm{tg}\,\beta \) ak \( \cos\alpha=\frac5{13} \).}
{\( \frac5{12} \)}{\( \frac5{13} \)}{\( \frac{13}5 \)}{\( \frac{12}{13} \)}{}{}}
\LANGes{
\Question{
Dado el triángulo rectángulo \( ABC \), siendo \( C \) el vértice del ángulo recto, calcula \( \mathrm{tg}\,\beta \) si \( \cos\alpha=\frac5{13} \).}
{\( \frac5{12} \)}{\( \frac5{13} \)}{\( \frac{13}5 \)}{\( \frac{12}{13} \)}{}{}}
\End

\Data\ID{1103076902}
\Updated{2022-04-23 05:04:05}
\Level{B}
\SubArea{Triangles}
\IMG{1103076902.pdf}
\LANGen{
\Question{
Given a triangle \( ABC \), which of the following equalities is valid?}
{\( \frac a{\sin\alpha} = \frac b{\sin \beta} \)}{\( \frac ab = \frac{\sin \beta}{\sin \alpha} \)}{\( \frac a{\sin\alpha} =\frac{\sin\gamma}c \)}{\( \frac c{\sin\gamma} = \frac{\sin \alpha}a \)}{}{}}
\LANGpl{
\Question{
Dany jest trójkąt \( ABC \), która z poniższych równości jest prawdziwa?}
{\( \frac a{\sin\alpha} = \frac b{\sin \beta} \)}{\( \frac ab = \frac{\sin \beta}{\sin \alpha} \)}{\( \frac a{\sin\alpha} =\frac{\sin\gamma}c \)}{\( \frac c{\sin\gamma} = \frac{\sin \alpha}a \)}{}{}}
\LANGcs{
\Question{
Je dán pravoúhlý trojúhelník \( ABC \). Která z následujících rovností platí?}
{\( \frac a{\sin\alpha} = \frac b{\sin \beta} \)}{\( \frac ab = \frac{\sin \beta}{\sin \alpha} \)}{\( \frac a{\sin\alpha} =\frac{\sin\gamma}c \)}{\( \frac c{\sin\gamma} = \frac{\sin \alpha}a \)}{}{}}
\LANGsk{
\Question{
Daný je trojuholník  \( ABC \).  Vyberte pravdivé tvrdenie.}
{\( \frac a{\sin\alpha} = \frac b{\sin \beta} \)}{\( \frac ab = \frac{\sin \beta}{\sin \alpha} \)}{\( \frac a{\sin\alpha} =\frac{\sin\gamma}c \)}{\( \frac c{\sin\gamma} = \frac{\sin \alpha}a \)}{}{}}
\LANGes{
\Question{
Dado el triángulo \( ABC \), ¿Cuál de las siguientes igualdades es válida?}
{\( \frac a{\sin\alpha} = \frac b{\sin \beta} \)}{\( \frac ab = \frac{\sin \beta}{\sin \alpha} \)}{\( \frac a{\sin\alpha} =\frac{\sin\gamma}c \)}{\( \frac c{\sin\gamma} = \frac{\sin \alpha}a \)}{}{}}
\End

\Data\ID{1103076903}
\Updated{2022-04-23 05:04:05}
\Level{B}
\SubArea{Triangles}
\IMG{1103076903.pdf}
\LANGen{
\Question{
Which of the following formulas applies to the given triangle \( ABC \)? (See the picture.)}
{\( a^2 = b^2+c^2 -2bc\cos\alpha \)}{\( b^2 = a^2 + c^2 - 2bc\cos\alpha \)}{\( a^2 = b^2 +c^2-2bc\cos\gamma \)}{\( a^2 = b^2+c^2+2bc\cos\alpha \)}{}{}}
\LANGpl{
\Question{
Który z poniższych wzorów odnosi się do trójkąta \( ABC \)?}
{\( a^2 = b^2+c^2 -2bc\cos\alpha \)}{\( b^2 = a^2 + c^2 - 2bc\cos\alpha \)}{\( a^2 = b^2 +c^2-2bc\cos\gamma \)}{\( a^2 = b^2+c^2+2bc\cos\alpha \)}{}{}}
\LANGcs{
\Question{
Která z následujících rovností platí pro trojúhelník \( ABC \)?}
{\( a^2 = b^2+c^2 -2bc\cos\alpha \)}{\( b^2 = a^2 + c^2 - 2bc\cos\alpha \)}{\( a^2 = b^2 +c^2-2bc\cos\gamma \)}{\( a^2 = b^2+c^2+2bc\cos\alpha \)}{}{}}
\LANGsk{
\Question{
Ktorý z nasledujúcich vzorcov platí pre trojuholník \( ABC \)?}
{\( a^2 = b^2+c^2 -2bc\cos\alpha \)}{\( b^2 = a^2 + c^2 - 2bc\cos\alpha \)}{\( a^2 = b^2 +c^2-2bc\cos\gamma \)}{\( a^2 = b^2+c^2+2bc\cos\alpha \)}{}{}}
\LANGes{
\Question{
¿Cuál de las siguientes igualdades es válida para el triángulo \( ABC \)?}
{\( a^2 = b^2+c^2 -2bc\cos\alpha \)}{\( b^2 = a^2 + c^2 - 2bc\cos\alpha \)}{\( a^2 = b^2 +c^2-2bc\cos\gamma \)}{\( a^2 = b^2+c^2+2bc\cos\alpha \)}{}{}}
\End

\Data\ID{1103076904}
\Updated{2020-07-15 03:07:21}
\Level{B}
\SubArea{Triangles}
\IMG{1103076904.pdf}
\LANGen{
\Question{
The picture shows an acute triangle \( ABC \). Its area \( S=\frac{\sqrt3|AB|\cdot|AC|}4 \). What is the measure of the angle \( \alpha \)?}
{\( 60^{\circ} \)}{\( 45^{\circ} \)}{\( 30^{\circ} \)}{\( 90^{\circ} \)}{}{}}
\LANGpl{
\Question{
Rysunek przedstawia trójkąt ostrokątny \( ABC \). Jego pole powierzchni \( S=\frac{\sqrt3|AB|\cdot|AC|}4 \). Jake jest miara kąta \( \alpha \)?}
{\( 60^{\circ} \)}{\( 45^{\circ} \)}{\( 30^{\circ} \)}{\( 90^{\circ} \)}{}{}}
\LANGcs{
\Question{
Na obrázku je znázorněný ostroúhlý trojúhelník \( ABC \). Jeho obsah \( S=\frac{\sqrt3|AB|\cdot|AC|}4 \). Jakou velikost má úhel \( \alpha \)?}
{\( 60^{\circ} \)}{\( 45^{\circ} \)}{\( 30^{\circ} \)}{\( 90^{\circ} \)}{}{}}
\LANGsk{
\Question{
Na obrázku je znázornený ostrouhlý trojuholník  \( ABC \). Jeho obsah \( S=\frac{\sqrt3|AB|\cdot|AC|}4 \). Akú veľkosť má uhol  \( \alpha \)?}
{\( 60^{\circ} \)}{\( 45^{\circ} \)}{\( 30^{\circ} \)}{\( 90^{\circ} \)}{}{}}
\LANGes{
\Question{
La imagen representa un triángulo acutángulo \( ABC \). Su área equivale a \( S=\frac{\sqrt3|AB|\cdot|AC|}4 \). ¿Cuál es la medida del ángulo \( \alpha \)?}
{\( 60^{\circ} \)}{\( 45^{\circ} \)}{\( 30^{\circ} \)}{\( 90^{\circ} \)}{}{}}
\End

\Data\ID{1103076907}
\Updated{2020-07-15 03:07:21}
\Level{B}
\SubArea{Triangles}
\IMG{1103076907_0.pdf}
\LANGen{
\Question{
\( ABC \) is a triangle with the sides of lengths \( c=15 \), \( b=6 \), and the measure of \( \measuredangle CAB \) is \( 150^{\circ} \). Which of the numbers gives the size of the angle \( BCA \) most accurately?}
{\( 18.46^{\circ} \)}{\( 11.54^{\circ} \)}{\( 5.77^{\circ} \)}{\( 9.23^{\circ} \)}{}{}}
\LANGpl{
\Question{
\( ABC \) jest trójkątem o bokach długości \( c=15 \), \( b=6 \). Miara kąta \( CAB \) jest równa \( 150^{\circ} \). Która z podanych liczb jest najdokładniejszą miarą kąta \( BCA \)?}
{\( 18{,}46^{\circ} \)}{\( 11{,}54^{\circ} \)}{\( 5{,}77^{\circ} \)}{\( 9{,}23^{\circ} \)}{}{}}
\LANGcs{
\Question{
\( ABC \) je trojúhelník o délkách stran \( c=15 \), \( b=6 \). Velikost  \( \measuredangle CAB \) je \( 150^{\circ} \). Které z uvedených čísel nejpřesněji udává velikost úhlu \( BCA \) ?}
{\( 18{,}46^{\circ} \)}{\( 11{,}54^{\circ} \)}{\( 5{,}77^{\circ} \)}{\( 9{,}23^{\circ} \)}{}{}}
\LANGsk{
\Question{
\( ABC \) je trojuholník s dĺžkami strán \( c=15 \), \( b=6 \). Veľkosť  \( \measuredangle CAB \) je \( 150^{\circ} \). Ktoré z uvedených čísel najpresnejšie udáva veľkosť uhla \( BCA \) ?}
{\( 18{,}46^{\circ} \)}{\( 11{,}54^{\circ} \)}{\( 5{,}77^{\circ} \)}{\( 9{,}23^{\circ} \)}{}{}}
\LANGes{
\Question{
Dado el triángulo \( ABC \) cuyas longitudes de los lados son \( c=15 \), \( b=6 \), y la medida del \( \measuredangle CAB \) es \( 150^{\circ} \). ¿Cuál de los números da la medida del ángulo \( BCA \) con la mayor precisión?}
{\( 18.46^{\circ} \)}{\( 11.54^{\circ} \)}{\( 5.77^{\circ} \)}{\( 9.23^{\circ} \)}{}{}}
\End

\Data\ID{1103076908}
\Updated{2020-07-15 03:07:21}
\Level{B}
\SubArea{Triangles}
\IMG{1103076908.pdf}
\LANGen{
\Question{
The area of an obtuse triangle is \( 4\,\mathrm{cm}^2 \) and the lengths of the sides containing the obtuse angle are \( 2\,\mathrm{cm} \) and \( 8\,\mathrm{cm} \). Give the measure of this angle.}
{\( 150^{\circ} \)}{\( 120^{\circ} \)}{\( 135^{\circ} \)}{\( 105^{\circ} \)}{}{}}
\LANGpl{
\Question{
Pole powierzchni trójkąta rozwartokątnego wynosi \( 4\,\mathrm{cm}^2 \), a długości boków zawierające kąt rozwarty są \( 2\,\mathrm{cm} \) i \( 8\,\mathrm{cm} \). Podaj miarę tego kąta.}
{\( 150^{\circ} \)}{\( 120^{\circ} \)}{\( 135^{\circ} \)}{\( 105^{\circ} \)}{}{}}
\LANGcs{
\Question{
Tupoúhlý trojúhelník má obsah  \( 4\,\mathrm{cm}^2 \) a strany, které svírají tupý úhel, jsou dlouhé \( 2\,\mathrm{cm} \) a \( 8\,\mathrm{cm} \). Určete velikost tupého úhlu.}
{\( 150^{\circ} \)}{\( 120^{\circ} \)}{\( 135^{\circ} \)}{\( 105^{\circ} \)}{}{}}
\LANGsk{
\Question{
Tupouhlý trojuholník má obsah  \( 4\,\mathrm{cm}^2 \) a strany zvierajúce tupý uhol  sú dlhé \( 2\,\mathrm{cm} \) a \( 8\,\mathrm{cm} \). Určte veľkosť tohto tupého uhla.}
{\( 150^{\circ} \)}{\( 120^{\circ} \)}{\( 135^{\circ} \)}{\( 105^{\circ} \)}{}{}}
\LANGes{
\Question{
El área de un triángulo obtusángulo es \( 4\,\mathrm{cm}^2 \) y los lados, que forman el ánguo obtuso, miden \( 2\,\mathrm{cm} \) y \( 8\,\mathrm{cm} \). Halla la medida del ángulo obtuso.}
{\( 150^{\circ} \)}{\( 120^{\circ} \)}{\( 135^{\circ} \)}{\( 105^{\circ} \)}{}{}}
\End

\Data\ID{1103077003}
\Updated{2021-04-19 06:04:04}
\Level{B}
\SubArea{Triangles}
\IMG{1103077003_0.pdf}
\LANGen{
\Question{
In an isosceles triangle \( ABC \), the base \( AB = 6\,\mathrm{cm} \), the measure of the angle \( ABC \) is \( 20^{\circ} \). The angle bisector of the angle \( BAC \) intersects the side \( BC \) at a point \( K \). Find the length of the line segment \( BK \). Round the result to two decimal places.}
{\( 2.08\,\mathrm{cm} \)}{\( 5.64\,\mathrm{cm} \)}{\( 2.05\,\mathrm{cm} \)}{\( 2.18\,\mathrm{cm} \)}{}{}}
\LANGpl{
\Question{
Dany jest trójkąt równoramienny \( ABC \) o podstawie  \( AB = 6\,\mathrm{cm} \) i kącie \( ABC \) o mierze \( 20^{\circ} \). Dwusieczna kąta \( BAC \) przecina bok  \( BC \) w punkcie \( K \). Oblicz długość odcinka \( BK \). Zaokrąglij do dwóch miejsc dziesiętnych.}
{\( 2{,}08\,\mathrm{cm} \)}{\( 5{,}64\,\mathrm{cm} \)}{\( 2{,}05\,\mathrm{cm} \)}{\( 2{,}18\,\mathrm{cm} \)}{}{}}
\LANGcs{
\Question{
V rovnoramenném trojúhelníku \( ABC \) se základnou \( AB = 6\,\mathrm{cm} \) má úhel \( ABC \) velikost  \( 20^{\circ} \). Osa vnitřního úhlu \( BAC \) protíná stranu \( BC \) v bodě  \( K \). Vypočítejte délku úsečky \( BK \). Výsledek uveďte s přesností na 2 desetinná místa.}
{\( 2{,}08\,\mathrm{cm} \)}{\( 5{,}64\,\mathrm{cm} \)}{\( 2{,}05\,\mathrm{cm} \)}{\( 2{,}18\,\mathrm{cm} \)}{}{}}
\LANGsk{
\Question{
V rovnoramennom trojuholníku  \( ABC \) so základňou  \( AB = 6\,\mathrm{cm} \) má uhol \( ABC \) veľkosť  \( 20^{\circ} \). Os vnútorného uhla  \( BAC \) pretína stranu  \( BC \) v bode  \( K \). Vypočítajte dĺžku úsečky  \( BK \). Výsledok uveďte s presnosťou na 2 desatinné miesta.}
{\( 2{,}08\,\mathrm{cm} \)}{\( 5{,}64\,\mathrm{cm} \)}{\( 2{,}05\,\mathrm{cm} \)}{\( 2{,}18\,\mathrm{cm} \)}{}{}}
\LANGes{
\Question{
En un triángulo isósceles \( ABC \), con una base de \( AB = 6\,\mathrm{cm} \), el ángulo \( ABC \) mide \( 20^{\circ} \). La bisectriz angular del ángulo \( BAC \)  corta con el lado \( BC \) en el punto \( K \). Halla la longitud del segmento \( BK \). Redondea el resultado a dos decimales.}
{\( 2.08\,\mathrm{cm} \)}{\( 5.64\,\mathrm{cm} \)}{\( 2.05\,\mathrm{cm} \)}{\( 2.18\,\mathrm{cm} \)}{}{}}
\End

\Data\ID{1103077004}
\Updated{2022-04-23 05:04:05}
\Level{B}
\SubArea{Triangles}
\IMG{1103077004_0.pdf}
\LANGen{
\Question{
In a triangle \( ABC \), \( a=15\,\mathrm{cm} \), \( b=6\,\mathrm{cm} \) and the measure of the angle \( CAB \) is \( 120^{\circ} \). Which of the following numbers gives as accurately as possible the measure of the angle \( ABC \)?}
{\( 20.27^{\circ} \)}{\( 25.66^{\circ} \)}{\( 60^{\circ} \)}{\( 30^{\circ} \)}{}{}}
\LANGpl{
\Question{
W trójkącie  \( ABC \), \( a=15\,\mathrm{cm} \), \( b=6\,\mathrm{cm} \) , a miara kąta  \( CAB \) wynosi \( 120^{\circ} \). Która z poniższych liczb wskazuje tak dokładnie jak to możliwe miarę kąta \( ABC \)?}
{\( 20{,}27^{\circ} \)}{\( 25{,}66^{\circ} \)}{\( 60^{\circ} \)}{\( 30^{\circ} \)}{}{}}
\LANGcs{
\Question{
V trojúhelníku \( ABC \), \( a=15\,\mathrm{cm} \), \( b=6\,\mathrm{cm} \) a velikost úhlu \( CAB \) je \( 120^{\circ} \). Které z uvedených čísel nejpřesněji udává velikost úhlu \( ABC \)?}
{\( 20{,}27^{\circ} \)}{\( 25{,}66^{\circ} \)}{\( 60^{\circ} \)}{\( 30^{\circ} \)}{}{}}
\LANGsk{
\Question{
V trojuholníku  \( ABC \), \( a=15\,\mathrm{cm} \), \( b=6\,\mathrm{cm} \) a veľkosť uhla \( CAB \) je \( 120^{\circ} \). Ktoré z uvedených čísel najpresnejšie udáva veľkosť uhla \( ABC \)?}
{\( 20{,}27^{\circ} \)}{\( 25{,}66^{\circ} \)}{\( 60^{\circ} \)}{\( 30^{\circ} \)}{}{}}
\LANGes{
\Question{
Dado el triángulo \( ABC \), dos de sus lados miden \( a=15\,\mathrm{cm} \), \( b=6\,\mathrm{cm} \) y la medida del ángulo \( CAB \) es de \( 120^{\circ} \). ¿Cuál de los siguientes números se aproxima más a la medida del ángulo \( ABC \)?}
{\( 20.27^{\circ} \)}{\( 25.66^{\circ} \)}{\( 60^{\circ} \)}{\( 30^{\circ} \)}{}{}}
\End

\Data\ID{1103077005}
\Updated{2020-07-15 03:07:21}
\Level{B}
\SubArea{Triangles}
\IMG{1103077005_0.pdf}
\LANGen{
\Question{
In a triangle \( ABC \), \( c=2\,\mathrm{cm} \), \( a=3\,\mathrm{cm} \) and \( \beta=60^{\circ} \). What is the length of side \( b \)?}
{\( \sqrt7\,\mathrm{cm} \)}{\( 13-6\sqrt3\,\mathrm{cm} \)}{\( 19\,\mathrm{cm} \)}{\( 13+6\sqrt3\,\mathrm{cm} \)}{}{}}
\LANGpl{
\Question{
W trójkącie  \( ABC \), \( c=2\,\mathrm{cm} \), \( a=3\,\mathrm{cm} \) i \( \beta=60^{\circ} \). Jaka jest długość boku  \( b \)?}
{\( \sqrt7\,\mathrm{cm} \)}{\( 13-6\sqrt3\,\mathrm{cm} \)}{\( 19\,\mathrm{cm} \)}{\( 13+6\sqrt3\,\mathrm{cm} \)}{}{}}
\LANGcs{
\Question{
V trojúhelníku \( ABC \), \( c=2\,\mathrm{cm} \), \( a=3\,\mathrm{cm} \) a \( \beta=60^{\circ} \). Jaká je délka strany \( b \)?}
{\( \sqrt7\,\mathrm{cm} \)}{\( 13-6\sqrt3\,\mathrm{cm} \)}{\( 19\,\mathrm{cm} \)}{\( 13+6\sqrt3\,\mathrm{cm} \)}{}{}}
\LANGsk{
\Question{
V trojuholníku \( ABC \), \( c=2\,\mathrm{cm} \), \( a=3\,\mathrm{cm} \) a \( \beta=60^{\circ} \). Vypočítajte dĺžku strany \( b \).}
{\( \sqrt7\,\mathrm{cm} \)}{\( 13-6\sqrt3\,\mathrm{cm} \)}{\( 19\,\mathrm{cm} \)}{\( 13+6\sqrt3\,\mathrm{cm} \)}{}{}}
\LANGes{
\Question{
Dado el triángulo \( ABC \), \( c=2\,\mathrm{cm} \), \( a=3\,\mathrm{cm} \) y \( \beta=60^{\circ} \). ¿Cuánto mide el lado \( b \)?}
{\( \sqrt7\,\mathrm{cm} \)}{\( 13-6\sqrt3\,\mathrm{cm} \)}{\( 19\,\mathrm{cm} \)}{\( 13+6\sqrt3\,\mathrm{cm} \)}{}{}}
\End

\Data\ID{1103077007}
\Updated{2021-04-19 06:04:47}
\Level{B}
\SubArea{Triangles}
\IMG{1103077007.pdf}
\LANGen{
\Question{
In a triangle \( ABC \), \( a:b=1:2 \) and the measure of the angle \( BAC \) is \( 30^{\circ} \). Determine the size of the smallest interior angle of the triangle.}
{\( 30^{\circ} \)}{\( 20^{\circ} \)}{\( 10^{\circ} \)}{\( 90^{\circ} \)}{}{}}
\LANGpl{
\Question{
W trójkącie \( ABC \), \( a:b=1:2 \) a miara kąta \( BAC \) wynosi \( 30^{\circ} \). Oblicz miarę najmniejszego kąta wewnętrznego tego trójkąta.}
{\( 30^{\circ} \)}{\( 20^{\circ} \)}{\( 10^{\circ} \)}{\( 90^{\circ} \)}{}{}}
\LANGcs{
\Question{
V trojúhelníku \( ABC \), \( a:b=1:2 \)  a velikost úhlu \( BAC \) je \( 30^{\circ} \). Určete velikost nejmenšího vnitřního úhlu trojúhelníku.}
{\( 30^{\circ} \)}{\( 20^{\circ} \)}{\( 10^{\circ} \)}{\( 90^{\circ} \)}{}{}}
\LANGsk{
\Question{
V trojuholníku  \( ABC \), \( a:b=1:2 \) a veľkosť vnútorného uhla \( BAC \) je \( 30^{\circ} \). Nájdite veľkosť najmenšieho vnútorného uhla trojuholníka.}
{\( 30^{\circ} \)}{\( 20^{\circ} \)}{\( 10^{\circ} \)}{\( 90^{\circ} \)}{}{}}
\LANGes{
\Question{
Dado el triángulo \( ABC \), \( a:b=1:2 \) cuya medida del ángulo \( BAC \) es \( 30^{\circ} \). Calcula la medida del ángulo interior más pequeño del triángulo.}
{\( 30^{\circ} \)}{\( 20^{\circ} \)}{\( 10^{\circ} \)}{\( 90^{\circ} \)}{}{}}
\End

\Data\ID{1103077008}
\Updated{2020-07-15 03:07:21}
\Level{B}
\SubArea{Triangles}
\IMG{1103077008.pdf}
\LANGen{
\Question{
Given a triangle \( ABC \), the length of the median from \( C \) is \( 9\,\mathrm{cm} \) and the length of the median from \( B \) is \( 6\,\mathrm{cm} \). Let \( T \) be the centroid, and \( S \) be the midpoint of \( AC \). The measure of the angle \( BTC \) is \( 120^{\circ} \). Find the length of the side \( AC \).}
{\( 4\sqrt7\,\mathrm{cm} \)}{\( 2\sqrt7\,\mathrm{cm} \)}{\( 2\sqrt{13}\,\mathrm{cm} \)}{\( 4\sqrt{13}\,\mathrm{cm} \)}{}{}}
\LANGpl{
\Question{
Dany jest trójkąt  \( ABC \), długość środkowej z wierzchołka  \( C \) wynosi \( 9\,\mathrm{cm} \) i długość środkowej z wierzchołka \( B \) wynosi \( 6\,\mathrm{cm} \). Niech \( T \) będzie środkiem ciężkości trójkąta, a \( S \) środkiem \( AC \). Miara kąta \( BTC \) wynosi \( 120^{\circ} \). Wyznacz długość boku  \( AC \).}
{\( 4\sqrt7\,\mathrm{cm} \)}{\( 2\sqrt7\,\mathrm{cm} \)}{\( 2\sqrt{13}\,\mathrm{cm} \)}{\( 4\sqrt{13}\,\mathrm{cm} \)}{}{}}
\LANGcs{
\Question{
Je dán trojúhelník \( ABC \). Těžnice na stranu \( c \) měří \( 9\,\mathrm{cm} \) a na stranu \( b \) měří \( 6\,\mathrm{cm} \). Bod \( T \) je těžištěm trojúhelníku a bod \( S \) je středem strany \( AC \). Velikost úhlu  \( BTC \) je \( 120^{\circ} \). Vypočítejte velikost strany \( AC \).}
{\( 4\sqrt7\,\mathrm{cm} \)}{\( 2\sqrt7\,\mathrm{cm} \)}{\( 2\sqrt{13}\,\mathrm{cm} \)}{\( 4\sqrt{13}\,\mathrm{cm} \)}{}{}}
\LANGsk{
\Question{
Daný je trojuholník  \( ABC \). Ťažnica na stranu \( c \) meria \( 9\,\mathrm{cm} \) a na stranu \( b \) meria \( 6\,\mathrm{cm} \). Bod \( T \) je ťažisko trojuholníka a bod \( S \) je stred strany  \( AC \). Veľkosť uhla  \( BTC \) je \( 120^{\circ} \). Vypočítajte veľkosť strany  \( AC \).}
{\( 4\sqrt7\,\mathrm{cm} \)}{\( 2\sqrt7\,\mathrm{cm} \)}{\( 2\sqrt{13}\,\mathrm{cm} \)}{\( 4\sqrt{13}\,\mathrm{cm} \)}{}{}}
\LANGes{
\Question{
Dado el triángulo \( ABC \), la longitud de la mediana desde \( C \) es \( 9\,\mathrm{cm} \) y la longitud de la mediana desde \( B \) es \( 6\,\mathrm{cm} \). El punto \( T \) es el baricentro y el punto \( S \) es el centro de \( AC \). La medida del ángulo \( BTC \) es \( 120^{\circ} \). Halla la longitud del lado \( AC \).}
{\( 4\sqrt7\,\mathrm{cm} \)}{\( 2\sqrt7\,\mathrm{cm} \)}{\( 2\sqrt{13}\,\mathrm{cm} \)}{\( 4\sqrt{13}\,\mathrm{cm} \)}{}{}}
\End

\Data\ID{1103077009}
\Updated{2022-04-23 05:04:05}
\Level{B}
\SubArea{Triangles}
\IMG{1103077009.pdf}
\LANGen{
\Question{
The angles \( \alpha \), \(\beta \), \( \gamma \) of a right-angled triangle \( ABC \) are in the ratio \( 1:2:3 \) (see the picture). From the following ratios of sides select the one that is equal to \( \sqrt3:1 \).}
{\( b:a \)}{\( a:b \)}{\( c:b \)}{\( c:a \)}{}{}}
\LANGpl{
\Question{
Kąty \( \alpha \), \(\beta \), \( \gamma \) w trójkącie prostokątnym \( ABC \) są w stosunku  \( 1:2:3 \). Z poniższych stosunków boków wybierz ten, który  jest równy \( \sqrt3:1 \).}
{\( b:a \)}{\( a:b \)}{\( c:b \)}{\( c:a \)}{}{}}
\LANGcs{
\Question{
Úhly \( \alpha \), \(\beta \), \( \gamma \) v pravoúhlém trojúhelníku \( ABC \) jsou v poměru \( 1:2:3 \). Které dvě strany tohoto trojúhelníku jsou v poměru \( \sqrt3:1 \).}
{\( b:a \)}{\( a:b \)}{\( c:b \)}{\( c:a \)}{}{}}
\LANGsk{
\Question{
Uhly \( \alpha \), \(\beta \), \( \gamma \) v pravouhlom trojuholníku \( ABC \) sú v pomere \( 1:2:3 \). Ktoré dve strany tohto trojuholníka sú v pomere  \( \sqrt3:1 \).}
{\( b:a \)}{\( a:b \)}{\( c:b \)}{\( c:a \)}{}{}}
\LANGes{
\Question{
Los ángulos \( \alpha \), \(\beta \), \( \gamma \) del triángulo rectángulo \( ABC \) están en proporción \( 1:2:3 \). ¿Cuáles de los lados están en proporción \( \sqrt3:1 \)?}
{\( b:a \)}{\( a:b \)}{\( c:b \)}{\( c:a \)}{}{}}
\End

\Data\ID{1103077011}
\Updated{2021-04-19 06:04:51}
\Level{B}
\SubArea{Triangles}
\IMG{1103077011.pdf}
\LANGen{
\Question{
Consider a triangle \( ABC \) with \( a=1\,\mathrm{cm} \) and \( b = \sqrt3\,\mathrm{cm} \). The angle opposite the longer side is double the angle opposite the shorter side. Find the area of the triangle.}
{\( \frac{\sqrt3}2\,\mathrm{cm}^2 \)}{\( 2\sqrt3\,\mathrm{cm}^2 \)}{\( \sqrt3\,\mathrm{cm}^2 \)}{\( \frac{\sqrt3}4\,\mathrm{cm}^2 \)}{}{}}
\LANGpl{
\Question{
Rozważ trójkąt \( ABC \) o bokoch \( a=1\,\mathrm{cm} \) i \( b = \sqrt3\,\mathrm{cm} \).  Kąt naprzeciwko dłuższego boku tego trójkąta jest dwa razy większy od kąta naprzeciw krótszego boku. Oblicz pole powierzchni tego trójkąta.}
{\( \frac{\sqrt3}2\,\mathrm{cm}^2 \)}{\( 2\sqrt3\,\mathrm{cm}^2 \)}{\( \sqrt3\,\mathrm{cm}^2 \)}{\( \frac{\sqrt3}4\,\mathrm{cm}^2 \)}{}{}}
\LANGcs{
\Question{
Vypočítejte obsah trojúhelníku \( ABC \), ve kterém \( a=1\,\mathrm{cm} \) a \( b = \sqrt3\,\mathrm{cm} \). Vnitřní úhel naproti delší straně je dvojnásobkem úhlu naproti kratší straně.}
{\( \frac{\sqrt3}2\,\mathrm{cm}^2 \)}{\( 2\sqrt3\,\mathrm{cm}^2 \)}{\( \sqrt3\,\mathrm{cm}^2 \)}{\( \frac{\sqrt3}4\,\mathrm{cm}^2 \)}{}{}}
\LANGsk{
\Question{
Vypočítajte obsah trojuholníka  \( ABC \), v ktorom  \( a=1\,\mathrm{cm} \) a \( b =\sqrt3\,\mathrm{cm} \). Vnútorný uhol oproti dlhšej strane je dvojnásobkom uhla oproti kratšej strane.}
{\( \frac{\sqrt3}2\,\mathrm{cm}^2 \)}{\( 2\sqrt3\,\mathrm{cm}^2 \)}{\( \sqrt3\,\mathrm{cm}^2 \)}{\( \frac{\sqrt3}4\,\mathrm{cm}^2 \)}{}{}}
\LANGes{
\Question{
Dado el triángulo \( ABC \) con \( a=1\,\mathrm{cm} \) y \( b = \sqrt3\,\mathrm{cm} \).Se cumple que el ángulo opuesto al lado más largo es el doble del ángulo opuesto al lado más corto. Calcula el área del triángulo.}
{\( \frac{\sqrt3}2\,\mathrm{cm}^2 \)}{\( 2\sqrt3\,\mathrm{cm}^2 \)}{\( \sqrt3\,\mathrm{cm}^2 \)}{\( \frac{\sqrt3}4\,\mathrm{cm}^2 \)}{}{}}
\End

\Data\ID{2000005601}
\Updated{2021-12-01 04:12:49}
\Level{B}
\SubArea{Triangles}
\IMG{2100005601_4-figure0.pdf}
\LANGen{
\Question{
A right-angled triangle \(ABC\) with the right angle at the vertex \(C\) is given in the picture.  The length of the side \(c\) is \(5\,\mathrm{cm}\) and the measure of the angle \(\alpha\) is \(35^{\circ}\). What is the length of  the side \(a\)?}
{\(5\sin{35^{\circ}}\,\mathrm{cm}\)}{\(\frac{5}{\sin{35^{\circ}}}\,\mathrm{cm}\)}{\(5\cos{35^{\circ}}\,\mathrm{cm}\)}{\(\frac{5}{\cos{35^{\circ}}}\,\mathrm{cm}\)}{}{}}
\LANGpl{
\Question{
Na rysunku przedstawiono trójkąt prostokątny \(ABC\)  z kątem prostym przy wierzchołku \(C\).  Długość boku \(c\) wynosi \(5\,\mathrm{cm}\) a miara kąta \(\alpha\) jest \(35^{\circ}\). Oblicz długośc boku \(a\).}
{\(5\sin{35^{\circ}}\,\mathrm{cm}\)}{\(\frac{5}{\sin{35^{\circ}}}\,\mathrm{cm}\)}{\(5\cos{35^{\circ}}\,\mathrm{cm}\)}{\(\frac{5}{\cos{35^{\circ}}}\,\mathrm{cm}\)}{}{}}
\LANGcs{
\Question{
Na obrázku je znázorněn pravoúhlý trojúhelník \(ABC\) s pravým úhlem u vrcholu \(C\).  Délka strany \(c\) je \(5\,\mathrm{cm}\) a velikost úhlu \(\alpha\) je \(35^{\circ}\). Jaká je délka strany \(a\)?}
{\(5\sin{35^{\circ}}\,\mathrm{cm}\)}{\(\frac{5}{\sin{35^{\circ}}}\,\mathrm{cm}\)}{\(5\cos{35^{\circ}}\,\mathrm{cm}\)}{\(\frac{5}{\cos{35^{\circ}}}\,\mathrm{cm}\)}{}{}}
\LANGsk{
\Question{
Na obrázku je pravouhlý trojuholník \(ABC\) s pravým uhlom pri vrchole  \(C\).  Dĺžka strany  \(c\) je \(5\,\mathrm{cm}\) a veľkosť uhla \(\alpha\) je \(35^{\circ}\). Akú dĺžku má strana \(a\)?}
{\(5\sin{35^{\circ}}\,\mathrm{cm}\)}{\(\frac{5}{\sin{35^{\circ}}}\,\mathrm{cm}\)}{\(5\cos{35^{\circ}}\,\mathrm{cm}\)}{\(\frac{5}{\cos{35^{\circ}}}\,\mathrm{cm}\)}{}{}}
\LANGes{
\Question{
Dado el triángulo rectángulo \(ABC\) con el ángulo recto en el vértice \(C\).  La longitud del lado \(c\) es \(5\,\mathrm{cm}\) y la medida del ángulo \(\alpha\) es \(35^{\circ}\). Calcula la longitud del lado \(a\).}
{\(5\sin{35^{\circ}}\,\mathrm{cm}\)}{\(\frac{5}{\sin{35^{\circ}}}\,\mathrm{cm}\)}{\(5\cos{35^{\circ}}\,\mathrm{cm}\)}{\(\frac{5}{\cos{35^{\circ}}}\,\mathrm{cm}\)}{}{}}
\End

\Data\ID{2000005602}
\Updated{2021-12-01 04:12:55}
\Level{B}
\SubArea{Triangles}
\IMG{2100005601_4-figure1.pdf}
\LANGen{
\Question{
A right-angled triangle \(ABC\) is given in the picture. Its hypotenuse is \(10\,\mathrm{cm}\) long and the measure of its internal angle \(\alpha\) is \(30^{\circ}\).  What are the lengths of the legs in the triangle?}
{\( a=5\,\mathrm{cm}\), \( b=5\sqrt{3}\,\mathrm{cm}\)}{\( a=5\sqrt{3}\,\mathrm{cm}\), \( b=5\,\mathrm{cm}\)}{\(a=10\cos{30^{\circ}}\,\mathrm{cm}\), \(b=10\sin{35^{\circ}}\,\mathrm{cm}\)}{\(a=\frac{\sin{30^{\circ}}}{10}\,\mathrm{cm}\), \(b=\frac{\cos{30^{\circ}}}{10}\,\mathrm{cm}\)}{}{}}
\LANGpl{
\Question{
Na rysunku przedstawiono trójkąt prostokątny\(ABC\). Jego przeciwprostokątna ma  \(10\,\mathrm{cm}\) długości  a miara kąta wewnętrznego \(\alpha\) to \(30^{\circ}\). Oblicz długości ramion trójkąta.}
{\( a=5\,\mathrm{cm}\), \( b=5\sqrt{3}\,\mathrm{cm}\)}{\( a=5\sqrt{3}\,\mathrm{cm}\), \( b=5\,\mathrm{cm}\)}{\(a=10\cos{30^{\circ}}\,\mathrm{cm}\), \(b=10\sin{35^{\circ}}\,\mathrm{cm}\)}{\(a=\frac{\sin{30^{\circ}}}{10}\,\mathrm{cm}\), \(b=\frac{\cos{30^{\circ}}}{10}\,\mathrm{cm}\)}{}{}}
\LANGcs{
\Question{
Na obrázku je znázorněn pravoúhlý trojúhelník \(ABC\). Jeho přepona je dlouhá \(10\,\mathrm{cm}\) a velikost vnitřního úhlu \(\alpha\) je \(30^{\circ}\).  Jaká je délka odvěsen trojúhelníku?}
{\( a=5\,\mathrm{cm}\), \( b=5\sqrt{3}\,\mathrm{cm}\)}{\( a=5\sqrt{3}\,\mathrm{cm}\), \( b=5\,\mathrm{cm}\)}{\(a=10\cos{30^{\circ}}\,\mathrm{cm}\), \(b=10\sin{35^{\circ}}\,\mathrm{cm}\)}{\(a=\frac{\sin{30^{\circ}}}{10}\,\mathrm{cm}\), \(b=\frac{\cos{30^{\circ}}}{10}\,\mathrm{cm}\)}{}{}}
\LANGsk{
\Question{
Na obrázku je pravouhlý trojuholník  \(ABC\).  Prepona má dĺžku \(10\,\mathrm{cm}\) a jeho vnútorný uhol \(\alpha\)  má veľkosť  \(30^{\circ}\).  Akú dĺžku majú odvesny trojuholníka?}
{\( a=5\,\mathrm{cm}\), \( b=5\sqrt{3}\,\mathrm{cm}\)}{\( a=5\sqrt{3}\,\mathrm{cm}\), \( b=5\,\mathrm{cm}\)}{\(a=10\cos{30^{\circ}}\,\mathrm{cm}\), \(b=10\sin{35^{\circ}}\,\mathrm{cm}\)}{\(a=\frac{\sin{30^{\circ}}}{10}\,\mathrm{cm}\), \(b=\frac{\cos{30^{\circ}}}{10}\,\mathrm{cm}\)}{}{}}
\LANGes{
\Question{
Dado el triángulo rectángulo \(ABC\). Su hipotenusa mide \(10\,\mathrm{cm}\) y la medida del ángulo interior \(\alpha\) es \(30^{\circ}\).  Calcula las medidas de sus catetos.}
{\( a=5\,\mathrm{cm}\), \( b=5\sqrt{3}\,\mathrm{cm}\)}{\( a=5\sqrt{3}\,\mathrm{cm}\), \( b=5\,\mathrm{cm}\)}{\(a=10\cos{30^{\circ}}\,\mathrm{cm}\), \(b=10\sin{35^{\circ}}\,\mathrm{cm}\)}{\(a=\frac{\sin{30^{\circ}}}{10}\,\mathrm{cm}\), \(b=\frac{\cos{30^{\circ}}}{10}\,\mathrm{cm}\)}{}{}}
\End

\Data\ID{2000005603}
\Updated{2021-12-01 04:12:03}
\Level{B}
\SubArea{Triangles}
\IMG{2100005601_4-figure2.pdf}
\LANGen{
\Question{
A right-angled triangle \(ABC\) with the right angle at the vertex \(C\) is given in the picture. Calculate the length of the side \(b\), if \(a=20\,\mathrm{cm}\) and \(\beta=34^{\circ}\).}
{\(b=\frac{20}{ \mathop{\mathrm{tg}} {56^{\circ}}}\,\mathrm{cm}\)}{\(b=\frac{20}{ \mathop{\mathrm{tg}} {34^{\circ}}}\,\mathrm{cm}\)}{\( b=20\sin{34^{\circ}}\,\mathrm{cm}\)}{\( b=20\cos{34^{\circ}}\,\mathrm{cm}\)}{}{}}
\LANGpl{
\Question{
Na rysunku przedstawiono trójkąt prostokątny \(ABC\) z kątem prostym przy wierzchołku \(C\). Oblicz długość boku \(b\) jeśli \(a=20\,\mathrm{cm}\) i \(\beta=34^{\circ}\).}
{\(b=\frac{20}{ \mathop{\mathrm{tg}} {56^{\circ}}}\,\mathrm{cm}\)}{\(b=\frac{20}{ \mathop{\mathrm{tg}} {34^{\circ}}}\,\mathrm{cm}\)}{\( b=20\sin{34^{\circ}}\,\mathrm{cm}\)}{\( b=20\cos{34^{\circ}}\,\mathrm{cm}\)}{}{}}
\LANGcs{
\Question{
Na obrázku je znázorněn pravoúhlý trojúhelník \(ABC\) s pravým úhlem u vrcholu \(C\). Vypočtěte délku strany \(b\), pokud platí, že \(a=20\,\mathrm{cm}\) a \(\beta=34^{\circ}\).}
{\(b=\frac{20}{ \mathop{\mathrm{tg}} {56^{\circ}}}\,\mathrm{cm}\)}{\(b=\frac{20}{ \mathop{\mathrm{tg}} {34^{\circ}}}\,\mathrm{cm}\)}{\( b=20\sin{34^{\circ}}\,\mathrm{cm}\)}{\( b=20\cos{34^{\circ}}\,\mathrm{cm}\)}{}{}}
\LANGsk{
\Question{
Na obrázku je pravouhlý trojuholník \(ABC\) s pravým uhlom pri vrchole  \(C\). Vypočítajte dĺžku strany \(b\), ak  \(a=20\,\mathrm{cm}\) a \(\beta=34^{\circ}\).}
{\(b=\frac{20}{ \mathop{\mathrm{tg}} {56^{\circ}}}\,\mathrm{cm}\)}{\(b=\frac{20}{ \mathop{\mathrm{tg}} {34^{\circ}}}\,\mathrm{cm}\)}{\( b=20\sin{34^{\circ}}\,\mathrm{cm}\)}{\( b=20\cos{34^{\circ}}\,\mathrm{cm}\)}{}{}}
\LANGes{
\Question{
Dado el triángulo rectángulo \(ABC\) con el ángulo recto en el vértice \(C\). Calcula la longitud del lado \(b\), suponiendo que \(a=20\,\mathrm{cm}\) y \(\beta=34^{\circ}\).}
{\(b=\frac{20}{ \mathop{\mathrm{tg}} {56^{\circ}}}\,\mathrm{cm}\)}{\(b=\frac{20}{ \mathop{\mathrm{tg}} {34^{\circ}}}\,\mathrm{cm}\)}{\( b=20\sin{34^{\circ}}\,\mathrm{cm}\)}{\( b=20\cos{34^{\circ}}\,\mathrm{cm}\)}{}{}}
\End

\Data\ID{2000005604}
\Updated{2021-12-01 04:12:09}
\Level{B}
\SubArea{Triangles}
\IMG{2100005601_4-figure3.pdf}
\LANGen{
\Question{
A right-angled triangle \(ABC\) with the right angle at the vertex \(C\) is given in the picture. Calculate the height \(v_c\), if \(a=20\,\mathrm{cm}\) and \(\beta=50^{\circ}\).}
{\( 20\sin{50^{\circ}}\)}{\( 20\cos{50^{\circ}}\)}{\( 20 \mathop{\mathrm{tg}} {50^{\circ}}\)}{\( \frac{20}{\sin{50^{\circ}}}\)}{}{}}
\LANGpl{
\Question{
Na rysunku przedstawiono trójkąt prostokątny \(ABC\) z kątem prostym przy wierzchołku \(C\). Oblicz wyskość \(v_c\), jesli \(a=20\,\mathrm{cm}\) i \(\beta=50^{\circ}\).}
{\( 20\sin{50^{\circ}}\)}{\( 20\cos{50^{\circ}}\)}{\( 20\mathop{\mathrm{tg}}{50^{\circ}}\)}{\( \frac{20}{\sin{50^{\circ}}}\)}{}{}}
\LANGcs{
\Question{
Na obrázku je znázorněn pravoúhlý trojúhelník \(ABC\) s pravým úhlem u vrcholu \(C\). Vypočtěte výšku \(v_c\), pokud platí, že \(a=20\,\mathrm{cm}\) a \(\beta=50^{\circ}\).}
{\( 20\sin{50^{\circ}}\)}{\( 20\cos{50^{\circ}}\)}{\( 20\mathop{\mathrm{tg}}{50^{\circ}}\)}{\( \frac{20}{\sin{50^{\circ}}}\)}{}{}}
\LANGsk{
\Question{
Na obrázku je pravouhlý trojuholník  \(ABC\) s pravým uhlom pri vrchole \(C\). Vypočítajte výšku  \(v_c\), ak  \(a=20\,\mathrm{cm}\) a \(\beta=50^{\circ}\).}
{\( 20\sin{50^{\circ}}\)}{\( 20\cos{50^{\circ}}\)}{\( 20\mathop{\mathrm{tg}}{50^{\circ}}\)}{\( \frac{20}{\sin{50^{\circ}}}\)}{}{}}
\LANGes{
\Question{
Dado el triángulo rectángulo \(ABC\) con el ángulo recto en el vértice \(C\). Calcula la altura \(v_c\), suponiendo que \(a=20\,\mathrm{cm}\) y \(\beta=50^{\circ}\).}
{\( 20\sin{50^{\circ}}\)}{\( 20\cos{50^{\circ}}\)}{\( 20\mathop{\mathrm{tg}}{50^{\circ}}\)}{\( \frac{20}{\sin{50^{\circ}}}\)}{}{}}
\End

\Data\ID{2010015202}
\Updated{2022-08-25 10:08:01}
\Level{B}
\SubArea{Triangles}
\IMG{2010015201-figure1.pdf}
\LANGen{
\Question{
The ladder is leaning against the wall of a house. Its length is \( 5 \) meters. How high does the ladder reach if the angle between it and the wall is \( 45^{\circ} \)? (See the picture.)}
{\( \frac{5\sqrt2}{2}\,\mathrm{m} \)}{\( \frac{5}{2}\,\mathrm{m} \)}{\( \frac{5\sqrt3}{2}\,\mathrm{m} \)}{\( \frac{10}{\sqrt2}\,\mathrm{m} \)}{}{}}
\LANGpl{
\Question{
Drabina opiera się o ścianę domu. Jej długość to \( 5 \) metrów. Jak wysoko sięga drabina, jeśli kąt między nią a ścianą wynosi \( 45^{\circ} \)? (Popatrz na obrazek.)}
{\( \frac{5\sqrt2}{2}\,\mathrm{m} \)}{\( \frac{5}{2}\,\mathrm{m} \)}{\( \frac{5\sqrt3}{2}\,\mathrm{m} \)}{\( \frac{10}{\sqrt2}\,\mathrm{m} \)}{}{}}
\LANGcs{
\Question{
O stěnu domu je opřený žebřík dlouhý \( 5 \) metrů. Jak vysoko žebřík dosáhne, pokud se stěnou domu svírá úhel \( 45^{\circ} \)? (Viz obrázek.)}
{\( \frac{5\sqrt2}{2}\,\mathrm{m} \)}{\( \frac{5}{2}\,\mathrm{m} \)}{\( \frac{5\sqrt3}{2}\,\mathrm{m} \)}{\( \frac{10}{\sqrt2}\,\mathrm{m} \)}{}{}}
\LANGsk{
\Question{
Rebrík je opretý o stenu domu. Jeho dĺžka je \( 5 \) metrov. Ako vysoko dosiahne rebrík, ak uhol medzi ním a stenou je \( 45^{\circ} \)? (Pozri obrázok.)}
{\( \frac{5\sqrt2}{2}\,\mathrm{m} \)}{\( \frac{5}{2}\,\mathrm{m} \)}{\( \frac{5\sqrt3}{2}\,\mathrm{m} \)}{\( \frac{10}{\sqrt2}\,\mathrm{m} \)}{}{}}
\LANGes{
\Question{
Una escalera se apoya en la pared de una casa. La escalera mide \( 5 \) metros. ¿A qué altura llega la escalera si forma un ángulo de \( 45^{\circ} \) con la pared?}
{\( \frac{5\sqrt2}{2}\,\mathrm{m} \)}{\( \frac{5}{2}\,\mathrm{m} \)}{\( \frac{5\sqrt3}{2}\,\mathrm{m} \)}{\( \frac{10}{\sqrt2}\,\mathrm{m} \)}{}{}}
\End

\Data\ID{2010015203}
\Updated{2022-08-25 10:08:04}
\Level{B}
\SubArea{Triangles}
\LANGen{
\Question{
In a triangle with interior angles \( 30^{\circ} \), \( 60^{\circ} \) and \( 90^{\circ} \), the shortest  side is \( 10\,\mathrm{cm} \) long. Find the length of its longest side.}
{\( 20\,\mathrm{cm} \)}{\( \frac{20}{\sqrt3}\,\mathrm{cm} \)}{\( 20\sqrt3\,\mathrm{cm} \)}{\(15\,\mathrm{cm} \)}{}{}}
\LANGpl{
\Question{
W trójkącie o kątach wewnętrznych \( 30^{\circ} \), \( 60^{\circ} \) i \( 90^{\circ} \), najkrótszy bok ma \( 10\,\mathrm{cm} \) długości. Oblicz długość jego najdłuższego boku.}
{\( 20\,\mathrm{cm} \)}{\( \frac{20}{\sqrt3}\,\mathrm{cm} \)}{\( 20\sqrt3\,\mathrm{cm} \)}{\(15\,\mathrm{cm} \)}{}{}}
\LANGcs{
\Question{
Velikosti vnitřních úhlů trojúhelníku jsou \( 30^{\circ} \), \( 60^{\circ} \) a \( 90^{\circ} \). Jeho nejkratší strana má velikost \( 10\,\mathrm{cm} \). Jaká je velikost jeho nejdelší strany?}
{\( 20\,\mathrm{cm} \)}{\( \frac{20}{\sqrt3}\,\mathrm{cm} \)}{\( 20\sqrt3\,\mathrm{cm} \)}{\(15\,\mathrm{cm} \)}{}{}}
\LANGsk{
\Question{
V trojuholníku s vnútornými uhlami \( 30^{\circ} \), \( 60^{\circ} \) a \( 90^{\circ} \), je najkrašia strana  dlhá \( 10\,\mathrm{cm} \). Nájdite dĺžku jeho najdlhšej strany.}
{\( 20\,\mathrm{cm} \)}{\( \frac{20}{\sqrt3}\,\mathrm{cm} \)}{\( 20\sqrt3\,\mathrm{cm} \)}{\(15\,\mathrm{cm} \)}{}{}}
\LANGes{
\Question{
En un triángulo con ángulos interiores de \( 30^{\circ} \), \( 60^{\circ} \) y \( 90^{\circ} \), el lado más corto mide \( 10\,\mathrm{cm} \). Calcula la longitud de su lado más largo.}
{\( 20\,\mathrm{cm} \)}{\( \frac{20}{\sqrt3}\,\mathrm{cm} \)}{\( 20\sqrt3\,\mathrm{cm} \)}{\(15\,\mathrm{cm} \)}{}{}}
\End

\Data\ID{2010015206}
\Updated{2022-08-25 10:08:04}
\Level{B}
\SubArea{Triangles}
\LANGen{
\Question{
The lengths of the sides in a triangle are \( a \), \( b \), \( c \) and the opposite angles are \( \alpha \), \( \beta \), \( \gamma \). Give the measure of \( \beta \) if \( b^2=a^2+c^2+ac\sqrt3 \).}
{\( 150^{\circ}\)}{\( 30^{\circ}\)}{\( 60^{\circ}\)}{\( 120^{\circ}\)}{}{}}
\LANGpl{
\Question{
Długości boków trójkąta to \( a \), \( b \), \( c \), a przeciwległe kąty to \( \alpha \), \( \beta \), \( \gamma \). Podaj miarę \( \beta \) jeśli \( b^2=a^2+c^2+ac\sqrt3 \).}
{\( 150^{\circ}\)}{\( 30^{\circ}\)}{\( 60^{\circ}\)}{\( 120^{\circ}\)}{}{}}
\LANGcs{
\Question{
Velikosti stran trojúhelníku jsou \( a \), \( b \), \( c \) a velikosti protilehlých úhlů \( \alpha \), \( \beta \), \( \gamma \). Vypočtěte velikost úhlu \( \beta \), pokud platí, že \( b^2=a^2+c^2+ac\sqrt3 \).}
{\( 150^{\circ}\)}{\( 30^{\circ}\)}{\( 60^{\circ}\)}{\( 120^{\circ}\)}{}{}}
\LANGsk{
\Question{
Dĺžky strán trojuholníka sú \( a \), \( b \), \( c \) a vnútorné uhly  \( \alpha \), \( \beta \), \( \gamma \). Vypočítajte veľkosť uhla \( \beta \) ak \( b^2=a^2+c^2+ac\sqrt3 \).}
{\( 150^{\circ}\)}{\( 30^{\circ}\)}{\( 60^{\circ}\)}{\( 120^{\circ}\)}{}{}}
\LANGes{
\Question{
En un triángulo, las longitudes de los lados son \( a \), \( b \), \( c \) y los ángulos opuestos son \( \alpha \), \( \beta \), \( \gamma \). Calcula la medida del ángulo \( \beta \) si \( b^2=a^2+c^2+ac\sqrt3 \).}
{\( 150^{\circ}\)}{\( 30^{\circ}\)}{\( 60^{\circ}\)}{\( 120^{\circ}\)}{}{}}
\End

\Data\ID{2010015207}
\Updated{2022-08-25 10:08:04}
\Level{B}
\SubArea{Triangles}
\LANGen{
\Question{
The lengths of the sides of a triangle are \( 4\,\mathrm{cm} \), \( 6\,\mathrm{cm} \) and \( 8\,\mathrm{cm} \). Calculate cosine of its greatest interior angle.}
{\( -\frac14\)}{\( 104.47^{\circ}\)}{\( 28.96^{\circ}\)}{\(\frac78\)}{}{}}
\LANGpl{
\Question{
Długości boków trójkąta to \( 4\,\mathrm{cm} \), \( 6\,\mathrm{cm} \) i \( 8\,\mathrm{cm} \). Oblicz cosinus jego największego kąta wewnętrznego.}
{\( -\frac14\)}{\( 104{,}47^{\circ}\)}{\( 28{,}96^{\circ}\)}{\(\frac78\)}{}{}}
\LANGcs{
\Question{
Velikosti stran trojúhelníka jsou \( 4\,\mathrm{cm} \), \( 6\,\mathrm{cm} \) a \( 8\,\mathrm{cm} \). Vypočtěte kosinus jeho největšího vntřního úhlu.}
{\( -\frac14\)}{\( 104{,}47^{\circ}\)}{\( 28{,}96^{\circ}\)}{\(\frac78\)}{}{}}
\LANGsk{
\Question{
Dĺžky strán trojuholníka sú \( 4\,\mathrm{cm} \), \( 6\,\mathrm{cm} \) a \( 8\,\mathrm{cm} \). Vypočítajte kosínus jeho najväčšieho vnútorného uhla.}
{\( -\frac14\)}{\( 104{,}47^{\circ}\)}{\( 28{,}96^{\circ}\)}{\(\frac78\)}{}{}}
\LANGes{
\Question{
Las longitudes de un triángulo son \( 4\,\mathrm{cm} \), \( 6\,\mathrm{cm} \) y \( 8\,\mathrm{cm} \). Calcula el valor del coseno de su ángulo interior más grande.}
{\( -\frac14\)}{\( 104.47^{\circ}\)}{\( 28.96^{\circ}\)}{\(\frac78\)}{}{}}
\End

\Data\ID{2010015210}
\Updated{2022-08-25 10:08:04}
\Level{B}
\SubArea{Triangles}
\IMG{2010015201-figure4.pdf}
\LANGen{
\Question{
Given a right-angled triangle \( ABC \) with the right angle at vertex \( C \), calculate \( \cos\beta \) if \(\sin\alpha=\frac4{7} \).}
{\( \frac4{7} \)}{\( \frac7{4} \)}{\( \frac{\sqrt{33}}{7} \)}{\( \frac7{\sqrt{33}} \)}{}{}}
\LANGpl{
\Question{
Mając trójkąt prostokątny \( ABC \)  z kątem prostym w wierzchołku \( C \), oblicz \( \cos\beta \),  jeśli \(\sin\alpha=\frac4{7} \).}
{\( \frac4{7} \)}{\( \frac7{4} \)}{\( \frac{\sqrt{33}}{7} \)}{\( \frac7{\sqrt{33}} \)}{}{}}
\LANGcs{
\Question{
Je dán pravoúhlý trojúhelník \( ABC \) s pravým úhlem u vrcholu \( C \). Čemu se rovná \( \cos\beta \), pokud \(\sin\alpha=\frac4{7} \)?}
{\( \frac4{7} \)}{\( \frac7{4} \)}{\( \frac{\sqrt{33}}{7} \)}{\( \frac7{\sqrt{33}} \)}{}{}}
\LANGsk{
\Question{
Daný je pravouhlý trojuholník  \( ABC \) s pravým uhlom pri vrchole  \( C \).  Vypočítajte \( \cos\beta \) ak \(\sin\alpha=\frac4{7} \).}
{\( \frac4{7} \)}{\( \frac7{4} \)}{\( \frac{\sqrt{33}}{7} \)}{\( \frac7{\sqrt{33}} \)}{}{}}
\LANGes{
\Question{
Dado el triángulo rectángulo \( ABC \), siendo \( C \) el vértice del ángulo recto, calcula \( \cos\beta \) suponiendo que \(\sin\alpha=\frac4{7} \).}
{\( \frac4{7} \)}{\( \frac7{4} \)}{\( \frac{\sqrt{33}}{7} \)}{\( \frac7{\sqrt{33}} \)}{}{}}
\End

\Data\ID{2010015301}
\Updated{2022-08-25 10:08:04}
\Level{B}
\SubArea{Triangles}
\IMG{2010015301-figure0.pdf}
\LANGen{
\Question{
An artillery battery is placed on a cliff. From the edge of the cliff \( 200\,\mathrm{m} \) high, the angle of depression to the ship on the sea is \( 10^{\circ} \). What is the distance \( d \) (See the picture.) from the cliff to the ship? Round the result to two decimal places.}
{\( 1134.26\,\mathrm{m} \)}{\( 1151.75\,\mathrm{m} \)}{\( 35.27\,\mathrm{m} \)}{\( 203.09\,\mathrm{m} \)}{}{}}
\LANGpl{
\Question{
Bateria artylerii jest umieszczona na klifie. Od krawędzi klifu o wysokości \( 200\,\mathrm{m} \) kąt obniżenia do statku na morzu wynosi \( 10^{\circ} \). Jaka jest odległość \( d \) (Popatrz na rysunek.) od klifu do statku? Zaokrąglij wynik do dwóch miejsc po przecinku.}
{\( 1134{,}26\,\mathrm{m} \)}{\( 1151{,}75\,\mathrm{m} \)}{\( 35{,}27\,\mathrm{m} \)}{\( 203{,}09\,\mathrm{m} \)}{}{}}
\LANGcs{
\Question{
Na okraji útesu vysokého \( 200\,\mathrm{m} \) je umístěna dělostřelecká baterie. Hloubkový úhel k lodi na moři je \( 10^{\circ} \). Jaká je vzdálenost \( d \) (Viz obrázek.) od útesu k lodi? Výsledek zaokrouhlete na dvě desetinná místa.}
{\( 1134{,}26\,\mathrm{m} \)}{\( 1151{,}75\,\mathrm{m} \)}{\( 35{,}27\,\mathrm{m} \)}{\( 203{,}09\,\mathrm{m} \)}{}{}}
\LANGsk{
\Question{
Delostrelecká batéria je umiestnená na kolmom útese vysokom  \( 200\,\mathrm{m} \).  Z útesu pozorujeme  loď v hĺbkovom uhle \( 10^{\circ} \). Vypočítajte vzdialenosť  \( d \) - lode od útesu. (Pozri obrázok.) Výsledok zaokrúhlite na dve desatinné miesta.}
{\( 1134{,}26\,\mathrm{m} \)}{\( 1151{,}75\,\mathrm{m} \)}{\( 35{,}27\,\mathrm{m} \)}{\( 203{,}09\,\mathrm{m} \)}{}{}}
\LANGes{
\Question{
Una batería de artillería está situada en un acantilado de \( 200\,\mathrm{m} \) de altitud. ¿Cuál es la distancia \( d \) del acantilado al barco, observado desde el acantilado bajo el ángulo de depresión de \( 10^{\circ} \)?}
{\( 1134.26\,\mathrm{m} \)}{\( 1151.75\,\mathrm{m} \)}{\( 35.27\,\mathrm{m} \)}{\( 203.09\,\mathrm{m} \)}{}{}}
\End

\Data\ID{2010015303}
\Updated{2022-08-25 10:08:18}
\Level{B}
\SubArea{Triangles}
\IMG{2010015301-figure2.pdf}
\LANGen{
\Question{
Triangle \( ABC \) is given (See the picture.). Choose the correct statement, if \(r\) is a radius of its circumcircle.}
{\( \frac{b}{\sin\beta} = 2r \)}{\( \frac{a}{\sin \alpha}= \frac{\sin\beta}{b}\)}{\( c \sin \alpha = b \sin \gamma \)}{\( \frac{a}{\sin \alpha}= r\)}{}{}}
\LANGpl{
\Question{
Dany jest trójkąt \( ABC \). Wybierz prawidłowe stwierdzenie, jeśli \(r\) jest promieniem jego okręgu opisanego.}
{\( \frac{b}{\sin\beta} = 2r \)}{\( \frac{a}{\sin \alpha}= \frac{\sin\beta}{b}\)}{\( c \sin \alpha = b \sin \gamma \)}{\( \frac{a}{\sin \alpha}= r\)}{}{}}
\LANGcs{
\Question{
Je dán trojúhelník \( ABC \) (Viz obrázek.). Vyberte pravdivé tvrzení, pokud víme, že \(r\) je poloměrem kružnice trojuhelníku opsané.}
{\( \frac{b}{\sin\beta} = 2r \)}{\( \frac{a}{\sin \alpha}= \frac{\sin\beta}{b}\)}{\( c \sin \alpha = b \sin \gamma \)}{\( \frac{a}{\sin \alpha}= r\)}{}{}}
\LANGsk{
\Question{
Na obrázku je trojuholník \( ABC \). Vyberte správne tvrdenie, ak \(r\) je polomer jeho opísanej kružnice.}
{\( \frac{b}{\sin\beta} = 2r \)}{\( \frac{a}{\sin \alpha}= \frac{\sin\beta}{b}\)}{\( c \sin \alpha = b \sin \gamma \)}{\( \frac{a}{\sin \alpha}= r\)}{}{}}
\LANGes{
\Question{
Dado el triángulo \( ABC \). Elige la proposición lógica, suponiendo que \(r\) es un radio de su circunferencia circunscrita.}
{\( \frac{b}{\sin\beta} = 2r \)}{\( \frac{a}{\sin \alpha}= \frac{\sin\beta}{b}\)}{\( c \sin \alpha = b \sin \gamma \)}{\( \frac{a}{\sin \alpha}= r\)}{}{}}
\End

\Data\ID{2010015304}
\Updated{2022-08-25 10:08:18}
\Level{B}
\SubArea{Triangles}
\IMG{2010015301-figure2.pdf}
\LANGen{
\Question{
Which of the following formulas applies to the given triangle \( ABC \)? (See the picture.)}
{\( c^2 = a^2+b^2 -2ab\cos\gamma \)}{\( c^2 = a^2+b^2 +2ab\cos\gamma \)}{\( c^2 = a^2+b^2 -2ab\cos\alpha \)}{\( c^2 = a^2+b^2 -2ab\sin\beta \)}{}{}}
\LANGpl{
\Question{
Który z poniższych wzorów ma zastosowanie do danego trójkąta \( ABC \)? (Popatrz na rysunek.)}
{\( c^2 = a^2+b^2 -2ab\cos\gamma \)}{\( c^2 = a^2+b^2 +2ab\cos\gamma \)}{\( c^2 = a^2+b^2 -2ab\cos\alpha \)}{\( c^2 = a^2+b^2 -2ab\sin\beta \)}{}{}}
\LANGcs{
\Question{
Která z uvedených rovnic platí pro daný trojúhelník \( ABC \)? (Viz obrázek.)}
{\( c^2 = a^2+b^2 -2ab\cos\gamma \)}{\( c^2 = a^2+b^2 +2ab\cos\gamma \)}{\( c^2 = a^2+b^2 -2ab\cos\alpha \)}{\( c^2 = a^2+b^2 -2ab\sin\beta \)}{}{}}
\LANGsk{
\Question{
Ktorý z nasledujúcich vzorcov platí pre daný trojuholník \( ABC \)? (Pozri obrázok.)}
{\( c^2 = a^2+b^2 -2ab\cos\gamma \)}{\( c^2 = a^2+b^2 +2ab\cos\gamma \)}{\( c^2 = a^2+b^2 -2ab\cos\alpha \)}{\( c^2 = a^2+b^2 -2ab\sin\beta \)}{}{}}
\LANGes{
\Question{
¿Cuál de las siguientes igualdades es válida para el triángulo \( ABC \)?}
{\( c^2 = a^2+b^2 -2ab\cos\gamma \)}{\( c^2 = a^2+b^2 +2ab\cos\gamma \)}{\( c^2 = a^2+b^2 -2ab\cos\alpha \)}{\( c^2 = a^2+b^2 -2ab\sin\beta \)}{}{}}
\End

\Data\ID{2010015305}
\Updated{2022-08-25 10:08:18}
\Level{B}
\SubArea{Triangles}
\IMG{2010015301-figure3.pdf}
\LANGen{
\Question{
In a triangle \( ABC \), \( a=15\,\mathrm{cm} \), \( c=8\,\mathrm{cm} \) and the measure of the angle \( CAB \) is \( 120^{\circ} \). Which of the following numbers gives as accurately as possible the measure of the angle \( BCA \)?}
{\( 27.51^{\circ} \)}{\( 16.12^{\circ} \)}{\( 30.13^{\circ} \)}{\( 12.45^{\circ} \)}{}{}}
\LANGpl{
\Question{
W trójkącie \( ABC \), \( a=15\,\mathrm{cm} \), \( c=8\,\mathrm{cm} \) a miara kąta \( CAB \) wynosi \( 120^{\circ} \). Która z poniższych liczb podaje możliwie najdokładniej miarę kąta \( BCA \)?}
{\( 27{,}51^{\circ} \)}{\( 16{,}12^{\circ} \)}{\( 30{,}13^{\circ} \)}{\( 12{,}45^{\circ} \)}{}{}}
\LANGcs{
\Question{
Je dán trojúhelník \( ABC \) se stranami \( a=15\,\mathrm{cm} \), \( c=8\,\mathrm{cm} \) a velikostí úhlu \( CAB \) \( 120^{\circ} \). Které z následujících čísel nejpřesněji vyjadřije velikost úhlu \( BCA \)?}
{\( 27{,}51^{\circ} \)}{\( 16{,}12^{\circ} \)}{\( 30{,}13^{\circ} \)}{\( 12{,}45^{\circ} \)}{}{}}
\LANGsk{
\Question{
V trojuholníku  \( ABC \), \( a=15\,\mathrm{cm} \), \( c=8\,\mathrm{cm} \) a veľkosť uhla \( CAB \) je \( 120^{\circ} \). Ktoré z uvedených čísel najpresnejšie udáva veľkosť uhla \( BCA \)?}
{\( 27{,}51^{\circ} \)}{\( 16{,}12^{\circ} \)}{\( 30{,}13^{\circ} \)}{\( 12{,}45^{\circ} \)}{}{}}
\LANGes{
\Question{
Dado el triángulo \( ABC \), dos de sus lados miden \( a=15\,\mathrm{cm} \) y \( c=8\,\mathrm{cm} \) y la medida del ángulo \( CAB \) es de \( 120^{\circ} \). ¿Cuál de los siguientes números se aproxima más a la medida del ángulo \( BCA \)?}
{\( 27.51^{\circ} \)}{\( 16.12^{\circ} \)}{\( 30.13^{\circ} \)}{\( 12.45^{\circ} \)}{}{}}
\End

\Data\ID{2010015306}
\Updated{2022-08-25 10:08:18}
\Level{B}
\SubArea{Triangles}
\IMG{2010015301-figure2.pdf}
\LANGen{
\Question{
The angles \( \alpha \), \(\beta \), \( \gamma \) of a right-angled triangle \( ABC \) are in the ratio \( 1:2:3 \) (See the picture.). From the following ratios of sides select the one that is equal to \( 1:\sqrt3 \).}
{\( a:b \)}{\( b:a\)}{\( c:b \)}{\( a:c \)}{}{}}
\LANGpl{
\Question{
Kąty \( \alpha \), \(\beta \), \( \gamma \) trójkąta prostokątnego \( ABC \) są w stosunku \( 1:2:3 \) (Popatrz na rysunek.). Z poniższych stosunków boków wybierz ten, który jest równy \( 1:\sqrt3 \).}
{\( a:b \)}{\( b:a\)}{\( c:b \)}{\( a:c \)}{}{}}
\LANGcs{
\Question{
Úhly \( \alpha \), \(\beta \), \( \gamma \) pravoúhlého trojúhelníku \( ABC \) jsou v poměru \( 1:2:3 \) (Viz obrázek.). Které dvě strany jsou v poměru \( 1:\sqrt3 \)?}
{\( a:b \)}{\( b:a\)}{\( c:b \)}{\( a:c \)}{}{}}
\LANGsk{
\Question{
Uhly \( \alpha \), \(\beta \), \( \gamma \) v pravouhlom trojuholníku \( ABC \) sú v pomere \( 1:2:3 \) (Pozri obrázok.).  Ktoré dve strany tohto trojuholníka sú v pomere \( 1:\sqrt3 \).}
{\( a:b \)}{\( b:a\)}{\( c:b \)}{\( a:c \)}{}{}}
\LANGes{
\Question{
Los ángulos \( \alpha \), \(\beta \), \( \gamma \) del triángulo rectángulo \( ABC \) están en proporción \( 1:2:3 \). ¿Cuáles de los lados están en proporción \( 1:\sqrt3 \)?}
{\( a:b \)}{\( b:a\)}{\( c:b \)}{\( a:c \)}{}{}}
\End

\Data\ID{2010015307}
\Updated{2022-08-25 10:08:18}
\Level{B}
\SubArea{Triangles}
\IMG{2010015301-figure4.pdf}
\LANGen{
\Question{
The angle of elevation of a straight road is \(9^{\circ }\). The distance between two places measured along the road is \(2\, \mathrm{km}\). For these places, find the difference in altitudes, i.e. the vertical distance, and round the result to the nearest meter. (See the picture.)}
{\(313\, \mathrm{m}\)}{\(1975\, \mathrm{m}\)}{\(317\, \mathrm{m}\)}{\(78\, \mathrm{m}\)}{}{}}
\LANGpl{
\Question{
Kąt wzniesienia prostej drogi wynosi \(9^{\circ }\).  Odległość między dwoma miejscami mierzona wzdłuż drogi wynosi \(2\, \mathrm{km}\). Znajdź różnicę wysokości tych miejsc, czyli odległość w pionie, i zaokrąglij wynik do najbliższego metra. (Popatrz rysunek.)}
{\(313\, \mathrm{m}\)}{\(1975\, \mathrm{m}\)}{\(317\, \mathrm{m}\)}{\(78\, \mathrm{m}\)}{}{}}
\LANGcs{
\Question{
Rovná silnice má stoupání o velikosti \(9^{\circ }\). Jaký je rozdíl v nadmořské výšce dvou míst, která jsou od sebe vzdálená \(2\, \mathrm{km}\)? Výsledek zaokrouhli na metry (Viz obrázek.).}
{\(313\, \mathrm{m}\)}{\(1975\, \mathrm{m}\)}{\(317\, \mathrm{m}\)}{\(78\, \mathrm{m}\)}{}{}}
\LANGsk{
\Question{
Cesta má klesanie  \(9^{\circ }\). O koľko metrov sa líši nadmorská výška dvoch miest,  ktoré sú od seba po ceste vzdialené  \(2\, \mathrm{km}\)? (Pozri obrázok.) Výsledok zaokrúhlite na celé metre.}
{\(313\, \mathrm{m}\)}{\(1975\, \mathrm{m}\)}{\(317\, \mathrm{m}\)}{\(78\, \mathrm{m}\)}{}{}}
\LANGes{
\Question{
Una carretera recta tiene un ángulo de elevación de \(9^{\circ }\) respecto al plano horizontal. La distancia entre dos lugares (medida a lo largo de la carretera) es \(2\, \mathrm{km}\). Halla la diferencia en altitudes (es decir la distancia vertical) para estos lugares y expresa el resultado en metros.}
{\(313\, \mathrm{m}\)}{\(1975\, \mathrm{m}\)}{\(317\, \mathrm{m}\)}{\(78\, \mathrm{m}\)}{}{}}
\End

\Data\ID{2010015308}
\Updated{2022-08-25 10:08:18}
\Level{B}
\SubArea{Triangles}
\IMG{2010015301-figure6.pdf}
\LANGen{
\Question{
A roof gable has the shape of an isosceles triangle with the base of \(14\, \mathrm{m}\) and the height \(6\,\mathrm{m}\). What is the angle between the roof and the horizontal direction? (See the picture.) Round your result to one decimal place.}
{\(40.6^{\circ}\)}{\(49.4^{\circ}\)}{\(59.0^{\circ}\)}{\(31.0^{\circ}\)}{}{}}
\LANGpl{
\Question{
Dach szczytowy ma kształt trójkąta równoramiennego o podstawie \(14\, \mathrm{m}\) i wysokości \(6\,\mathrm{m}\). Jaki jest kąt między dachem a kierunkiem poziomym? (Popatrz rysunek.) Zaokrąglij wynik do jednego miejsca po przecinku.}
{\(40{,}6^{\circ}\)}{\(49{,}4^{\circ}\)}{\(59{,}0^{\circ}\)}{\(31{,}0^{\circ}\)}{}{}}
\LANGcs{
\Question{
Střecha má tvar rovnoramenného trojúhelníku s délkou přepony \(14\, \mathrm{m}\) a výškou \(6\,\mathrm{m}\). Jaký je sklon střechy? (Viz obrázek.) Výsledek zaokrouhlete na jedno desetinné místo.}
{\(40{,}6^{\circ}\)}{\(49{,}4^{\circ}\)}{\(59{,}0^{\circ}\)}{\(31{,}0^{\circ}\)}{}{}}
\LANGsk{
\Question{
Strecha má tvar rovnoramenného trojuholníka. Jeho šírka je  \(14\, \mathrm{m}\) a výška \(6\,\mathrm{m}\). Aký sklon má strecha?  (Pozri obrázok.) Výsledok zaokrúhlite na jedno desatinné miesto.}
{\(40{,}6^{\circ}\)}{\(49{,}4^{\circ}\)}{\(59{,}0^{\circ}\)}{\(31{,}0^{\circ}\)}{}{}}
\LANGes{
\Question{
Una buhardilla tiene una forma de triángulo isósceles con una base de \(14\, \mathrm{m}\) y la altura de \(6\,\mathrm{m}\). Calcula el ángulo entre el techo y el plano horizontal. Redondea el resultado a los decimales de metros.}
{\(40.6^{\circ}\)}{\(49.4^{\circ}\)}{\(59.0^{\circ}\)}{\(31.0^{\circ}\)}{}{}}
\End

\Data\ID{2010015309}
\Updated{2022-08-25 10:08:18}
\Level{B}
\SubArea{Triangles}
\IMG{2010015301-figure8.pdf}
\LANGen{
\Question{
The right angled triangle \(ABC\) is given (See the picture.). Choose the correct statement.}
{\(\mathrm{tg}\,\alpha = \frac{a} 
{b}\)}{\(\sin \beta= \frac{a} 
{c}\)}{\(\cos \beta = \frac{a} 
{b}\)}{\(\mathrm{tg}\,\beta = \frac{a} 
{b}\)}{}{}}
\LANGpl{
\Question{
Dany jest trójkąt prostokątny \(ABC\)  (Popatrz na obrazek,). Wybierz poprawne stwierdzenie.}
{\(\mathrm{tg}\,\alpha = \frac{a} 
{b}\)}{\(\sin \beta= \frac{a} 
{c}\)}{\(\cos \beta = \frac{a} 
{b}\)}{\(\mathrm{tg}\,\beta = \frac{a} 
{b}\)}{}{}}
\LANGcs{
\Question{
Je dán pravoúhlý trojíuhelník \(ABC\) (Viz obrázek.). Vyberte pravdivé tvrzení.}
{\(\mathrm{tg}\,\alpha = \frac{a} 
{b}\)}{\(\sin \beta= \frac{a} 
{c}\)}{\(\cos \beta = \frac{a} 
{b}\)}{\(\mathrm{tg}\,\beta = \frac{a} 
{b}\)}{}{}}
\LANGsk{
\Question{
Na obrázku je pravouhlý trojuholník  \(ABC\). Vyberte správne tvrdenie.}
{\(\mathrm{tg}\,\alpha = \frac{a} 
{b}\)}{\(\sin \beta= \frac{a} 
{c}\)}{\(\cos \beta = \frac{a} 
{b}\)}{\(\mathrm{tg}\,\beta = \frac{a} 
{b}\)}{}{}}
\LANGes{
\Question{
Dado un triángulo rectángulo \(ABC\) (Mira la imagen,). Halla la relación válida entre los ángulos y los lados del triángulo.}
{\(\mathrm{tg}\,\alpha = \frac{a} 
{b}\)}{\(\sin \beta= \frac{a} 
{c}\)}{\(\cos \beta = \frac{a} 
{b}\)}{\(\mathrm{tg}\,\beta = \frac{a} 
{b}\)}{}{}}
\End

\Data\ID{9000035001}
\Updated{2022-04-23 05:04:05}
\Level{B}
\SubArea{Triangles}
\IMG{2010015301-figure5.pdf}
\LANGen{
\Question{
The angle of elevation of a straight road is \(3^{\circ }30'\).
The distance between two places measured along the road is
\(2\, \mathrm{km}\). For these places, find the difference in altitudes, i.e. the vertical distance, and round the result to the nearest meter. (See the picture.)}
{\(122\, \mathrm{m}\)}{\(276\, \mathrm{m}\)}{\(98\, \mathrm{m}\)}{\(49\, \mathrm{m}\)}{}{}}
\LANGpl{
\Question{
Kąt podniesienia drogi wynosi  \(3^{\circ }30'\).
Odległość pomiędzy dwoma miejscami (mierzona wzdłuż drogi) jest równa
\(2\, \mathrm{km}\). Znajdź
różnicę wysokości (tj. odległość w pionie) między tymi miejscami i zaokrąglij do pełnych metrów.}
{\(122\, \mathrm{m}\)}{\(276\, \mathrm{m}\)}{\(98\, \mathrm{m}\)}{\(49\, \mathrm{m}\)}{}{}}
\LANGcs{
\Question{
Silnice má stoupání \(3^{\circ }30'\).
O kolik metrů se liší nadmořská výška dvou míst, která jsou od sebe po silnici
vzdálená \(2\, \mathrm{km}\)?
(Výsledek zaokrouhlete na celé metry.)}
{\(122\, \mathrm{m}\)}{\(276\, \mathrm{m}\)}{\(98\, \mathrm{m}\)}{\(49\, \mathrm{m}\)}{}{}}
\LANGsk{
\Question{
Cesta má stúpanie \(3^{\circ }30'\).
O koľko metrov sa líši nadmorská výška dvoch miest, ktoré sú od seba po ceste
vzdialené \(2\, \mathrm{km}\)?
(Výsledok zaokrúhlite na celé metre.)}
{\(122\, \mathrm{m}\)}{\(276\, \mathrm{m}\)}{\(98\, \mathrm{m}\)}{\(49\, \mathrm{m}\)}{}{}}
\LANGes{
\Question{
Una carretera recta tiene un ángulo de elevación de \(3^{\circ }30'\) respecto al plano horizontal. La distancia entre dos lugares (medida a lo largo de la carretera) es
\(2\, \mathrm{km}\). Halla la diferencia en altitudes (es decir. la distancia vertical) para estos lugares y expresa el resultado en metros.}
{\(122\, \mathrm{m}\)}{\(276\, \mathrm{m}\)}{\(98\, \mathrm{m}\)}{\(49\, \mathrm{m}\)}{}{}}
\End

\Data\ID{9000035003}
\Updated{2021-04-19 06:04:38}
\Level{B}
\SubArea{Triangles}
\LANGen{
\Question{
The tree of the height \(12\, \mathrm{m}\) 
is observed from the place horizontal with the base of the tree. The angle of elevation
is \(10^{\circ }\).
Find the distance of the observer from the base and round to the nearest meters.}
{\(68\, \mathrm{m}\)}{\(2\, \mathrm{m}\)}{\(12\, \mathrm{m}\)}{\(48\, \mathrm{m}\)}{}{}}
\LANGpl{
\Question{
Drzewo o wysokości równej  \(12\, \mathrm{m}\) 
jest obserwowane z miejsca poziomego względem jego podstawy. Wzniesienia kąt wynosi \(10^{\circ }\).
Powiedz w jakiej odległości od podstawy tego drzewa znajduje się obserwator i zaokrąglij wynik do pełnych metrów.}
{\(68\, \mathrm{m}\)}{\(2\, \mathrm{m}\)}{\(12\, \mathrm{m}\)}{\(48\, \mathrm{m}\)}{}{}}
\LANGcs{
\Question{
Strom vysoký \(12\) 
metrů pozorujeme z místa, které je ve vodorovné rovině s patou stromu. Vidíme ho
pod úhlem \(10^{\circ }\).
V jaké vzdálenosti od paty stojíme? (Výsledek zaokrouhlete na celé metry.)}
{\(68\, \mathrm{m}\)}{\(2\, \mathrm{m}\)}{\(12\, \mathrm{m}\)}{\(48\, \mathrm{m}\)}{}{}}
\LANGsk{
\Question{
Strom vysoký \(12\) 
metrov pozorujeme z miesta, ktoré je vo vodorovnej rovine s pätou stromu. Vidíme ho
pod uhlom \(10^{\circ }\).
V akej vzdialenosti od päty stojíme? (Výsledok zaokrúhlite na celé metre.)}
{\(68\, \mathrm{m}\)}{\(2\, \mathrm{m}\)}{\(12\, \mathrm{m}\)}{\(48\, \mathrm{m}\)}{}{}}
\LANGes{
\Question{
Observamos un árbol, que mide \(12\, \mathrm{m}\), desde un lugar horizontal a la base del árbol. El ángulo de elevación es \(10^{\circ }\). ¿A qué distancia estamos del árbol? Expresa el resultado en metros.}
{\(68\, \mathrm{m}\)}{\(2\, \mathrm{m}\)}{\(12\, \mathrm{m}\)}{\(48\, \mathrm{m}\)}{}{}}
\End

\Data\ID{9000035004}
\Updated{2020-07-15 03:07:34}
\Level{B}
\SubArea{Triangles}
\LANGen{
\Question{
The triangle \(ABC\) 
has the angle \(\beta = 59^{\circ }\) 
and the side \(a = 14\, \mathrm{cm}\).
Find the altitude \(v_{c}\) 
(the line segment which is perpendicular to the side
\(c\) and joins the
vertex \(C\) with
the side \(c\))
and round to the nearest centimeters.}
{\(12\, \mathrm{cm}\)}{\(7\, \mathrm{cm}\)}{\(10\, \mathrm{cm}\)}{\(23\, \mathrm{cm}\)}{}{}}
\LANGpl{
\Question{
Trójkąt  \(ABC\) 
posiada kąt  \(\beta = 59^{\circ }\) 
i bok  \(a = 14\, \mathrm{cm}\).
Znajdź wysokość  \(v_{c}\) 
(odcinek prostopadły do boku 
\(c\), który łączy wierzchołek \(C\) z 
bokiem  \(c\))
i zaokrąglij do pełnych centymetrów.}
{\(12\, \mathrm{cm}\)}{\(7\, \mathrm{cm}\)}{\(10\, \mathrm{cm}\)}{\(23\, \mathrm{cm}\)}{}{}}
\LANGcs{
\Question{
Vypočítejte výšku \(v_{c}\) 
v trojúhelníku \(ABC\),
je-li úhel \(\beta = 59^{\circ }\) 
a strana \(a = 14\, \mathrm{cm}\).
(Výsledek zaokrouhlete na celé centimetry.)}
{\(12\, \mathrm{cm}\)}{\(7\, \mathrm{cm}\)}{\(10\, \mathrm{cm}\)}{\(23\, \mathrm{cm}\)}{}{}}
\LANGsk{
\Question{
Vypočítajte výšku \(v_{c}\) 
v trojuholníku \(ABC\),
ak je uhol \(\beta = 59^{\circ }\) 
a strana \(a = 14\, \mathrm{cm}\).
(Výsledok zaokrúhlite na celé centimetre.)}
{\(12\, \mathrm{cm}\)}{\(7\, \mathrm{cm}\)}{\(10\, \mathrm{cm}\)}{\(23\, \mathrm{cm}\)}{}{}}
\LANGes{
\Question{
Calcula la altura \(v_{c}\)  del triángulo \(ABC\), si uno de los ángulos mide \(\beta = 59^{\circ }\) 
y el lado \(a = 14\, \mathrm{cm}\). (Redondea el resultado a centímetros.)}
{\(12\, \mathrm{cm}\)}{\(7\, \mathrm{cm}\)}{\(10\, \mathrm{cm}\)}{\(23\, \mathrm{cm}\)}{}{}}
\End

\Data\ID{9000035006}
\Updated{2020-07-15 03:07:34}
\Level{B}
\SubArea{Triangles}
\LANGen{
\Question{
A ladder of the length \(15\, \mathrm{m}\) 
leans against a wall. The angle between the ladder and the horizontal direction is
\(70^{\circ }\). Find
the height of the top of the ladder and round your answer to the nearest meters.}
{\(14\, \mathrm{m}\)}{\(13\, \mathrm{m}\)}{\(16\, \mathrm{m}\)}{\(15\, \mathrm{m}\)}{}{}}
\LANGpl{
\Question{
Drabina o długości  \(15\, \mathrm{m}\) 
jest oparta o ścianę. Kąt pomiędzy drabiną a poziomym kierunkiem wynosi 
\(70^{\circ }\). Znajdź wysokość szczytu drabiny i zaokrąglij wynik  do pełnych metrów.}
{\(14\, \mathrm{m}\)}{\(13\, \mathrm{m}\)}{\(16\, \mathrm{m}\)}{\(15\, \mathrm{m}\)}{}{}}
\LANGcs{
\Question{
Jak vysoko dosahuje žebřík, který je dlouhý
\(15\, \mathrm{m}\), svírá-li s vodorovnou
rovinou úhel \(70^{\circ }\)?
(Výsledek zaokrouhlete na celé metry.)}
{\(14\, \mathrm{m}\)}{\(13\, \mathrm{m}\)}{\(16\, \mathrm{m}\)}{\(15\, \mathrm{m}\)}{}{}}
\LANGsk{
\Question{
Ako vysoko dosiahne rebrík, ktorý je dlhý
\(15\, \mathrm{m}\), ak zviera s vodorovnou
rovinou uhol \(70^{\circ }\)?
(Výsledok zaokrúhlite na celé metre.)}
{\(14\, \mathrm{m}\)}{\(13\, \mathrm{m}\)}{\(16\, \mathrm{m}\)}{\(15\, \mathrm{m}\)}{}{}}
\LANGes{
\Question{
Una escalera, que tiene una longitud de \(15\, \mathrm{m}\), se apoya en una pared. El ángulo entre la escalera y el plano horizontal mide
\(70^{\circ }\). Calcula la altura de la parte superior de la escalera y redondea el resultado a metros.}
{\(14\, \mathrm{m}\)}{\(13\, \mathrm{m}\)}{\(16\, \mathrm{m}\)}{\(15\, \mathrm{m}\)}{}{}}
\End

\Data\ID{9000035007}
\Updated{2022-04-23 05:04:05}
\Level{B}
\SubArea{Triangles}
\IMG{2010015301-figure7.pdf}
\LANGen{
\Question{
A roof gable has the shape of an isosceles triangle with the base of \(14\, \mathrm{m}\).
The angle between the roof and the horizontal direction is
\(31^{\circ }\). Find the height of the gable. Round your result to one decimal place.}
{\(4.2\, \mathrm{m}\)}{\(5.9\, \mathrm{m}\)}{\(3.6\, \mathrm{m}\)}{\(11.2\, \mathrm{m}\)}{}{}}
\LANGpl{
\Question{
Szczyt dachu ma kształt trójkąta równoramiennego (trójkąt mający dwa boki równej długości) z podstawą 
o długości \(14\, \mathrm{m}\).
Kąt pomiędzy dachem a poziomym kierunkiem wynosi 
\(31^{\circ }\). Znajdź wysokość szczytu i zaokrąglij swoją odpowiedź do jednego miejsca po przecinku.}
{\(4.2\, \mathrm{m}\)}{\(5.9\, \mathrm{m}\)}{\(3.6\, \mathrm{m}\)}{\(11.2\, \mathrm{m}\)}{}{}}
\LANGcs{
\Question{
Štít střechy má tvar rovnoramenného trojúhelníka. Jeho šířka je
\(14\, \mathrm{m}\), sklon
střechy je \(31^{\circ }\).
Jaká je výška štítu v metrech? (Výsledek zaokrouhlete na jedno desetinné
místo.)}
{\(4{,}2\, \mathrm{m}\)}{\(5{,}9\, \mathrm{m}\)}{\(3{,}6\, \mathrm{m}\)}{\(11{,}2\, \mathrm{m}\)}{}{}}
\LANGsk{
\Question{
Štít strechy má tvar rovnoramenného trojuholníka. Jeho šírka je
\(14\, \mathrm{m}\), sklon
strechy je \(31^{\circ }\).
Aká je výška štítu v metroch? (Výsledok zaokrúhlite na jedno desatinné
miesto.)}
{\(4{,}2\, \mathrm{m}\)}{\(5{,}9\, \mathrm{m}\)}{\(3{,}6\, \mathrm{m}\)}{\(11{,}2\, \mathrm{m}\)}{}{}}
\LANGes{
\Question{
Una buhardilla tiene una forma de triángulo isósceles con una base de \(14\, \mathrm{m}\).
El ángulo entre el techo y el plano horizontal mide \(31^{\circ }\). Halla la altura de la buhardilla y redondea el resultado a los decimales de metros.}
{\(4.2\, \mathrm{m}\)}{\(5.9\, \mathrm{m}\)}{\(3.6\, \mathrm{m}\)}{\(11.2\, \mathrm{m}\)}{}{}}
\End

\Data\ID{9000035008}
\Updated{2021-04-24 04:04:57}
\Level{B}
\SubArea{Triangles}
\LANGen{
\Question{
Sun shines to the road at the angle \(53^{\circ }22'\).
An electric column near the road casts the shadow of the length
\(4.5\, \mathrm{m}\). Find
the height of the column and round your answer to the nearest meters.}
{\(6\, \mathrm{m}\)}{\(3\, \mathrm{m}\)}{\(4\, \mathrm{m}\)}{\(5\, \mathrm{m}\)}{}{}}
\LANGpl{
\Question{
Promienie słoneczne padają na drogę pod kątem  \(53^{\circ }22'\).
Słup elektryczny rzuca na drogę cień o długości 
\(4.5\, \mathrm{m}\). Wyznacz wysokość słupa i zaokrąglij swoją odpowiedź do pełnych metrów.}
{\(6\, \mathrm{m}\)}{\(3\, \mathrm{m}\)}{\(4\, \mathrm{m}\)}{\(5\, \mathrm{m}\)}{}{}}
\LANGcs{
\Question{
Sluneční paprsky dopadají na silnici pod úhlem
\(53^{\circ }22'\).
Určete, jak vysoký je sloup, který vrhá na silnici stín dlouhý
\(4{,}5\, \mathrm{m}\).
(Výsledek zaokrouhlete na celé metry.)}
{\(6\, \mathrm{m}\)}{\(3\, \mathrm{m}\)}{\(4\, \mathrm{m}\)}{\(5\, \mathrm{m}\)}{}{}}
\LANGsk{
\Question{
Slnečné lúče dopadajú na cestu pod uhlom
\(53^{\circ }22'\).
Určte, aký vysoký je stĺp, ktorý vrhá na cestu tieň dlhý
\(4{,}5\, \mathrm{m}\).
(Výsledok zaokrúhlite na celé metre.)}
{\(6\, \mathrm{m}\)}{\(3\, \mathrm{m}\)}{\(4\, \mathrm{m}\)}{\(5\, \mathrm{m}\)}{}{}}
\LANGes{
\Question{
El ángulo con el cual los rayos solares caen sobre la Tierra es de \(53^{\circ }22'\). Una columna eléctrica cerca de la carretera proyecta una sombra de \(4.5\, \mathrm{m}\). Calcula la altura de la columna y expresa el resultado en metros.}
{\(6\, \mathrm{m}\)}{\(3\, \mathrm{m}\)}{\(4\, \mathrm{m}\)}{\(5\, \mathrm{m}\)}{}{}}
\End

\Data\ID{9000035009}
\Updated{2021-04-24 04:04:55}
\Level{B}
\SubArea{Triangles}
\LANGen{
\Question{
Two forces act on the body at one point. The force
\(F_{1} = 760\, \mathrm{N}\) 
acts horizontally from left to the right and the force
\(F_{2} = 28.8\, \mathrm{N}\) acts
vertically from the top to the bottom. Find the angle between the horizontal direction
and the direction of the resulting force and round your answer to the nearest degrees
and minutes.}
{\(2^{\circ }10'\)}{\(3^{\circ }10'\)}{\(2^{\circ }20'\)}{\(3^{\circ }20'\)}{}{}}
\LANGpl{
\Question{
Dwie siły działają na ciało w jednym punkcie. Siła 
\(F_{1} = 760\, \mathrm{N}\) 
działa poziomo od lewej do prawej, i siła 
\(F_{2} = 28.8\, \mathrm{N}\) działa pionowo od góry do dołu. Znajdź miarę kąta pomiędzy poziomym kierunkiem a kierunkiem siły wynikowej i zaokrąglij swoją odpowiedź do pełnych stopni i minut.}
{\(2^{\circ }10'\)}{\(3^{\circ }10'\)}{\(2^{\circ }20'\)}{\(3^{\circ }20'\)}{}{}}
\LANGcs{
\Question{
Na těleso působí v jednom bodě dvě síly: síla
\(F_{1}\) o
velikosti \(760\, \mathrm{N}\) 
působí ve vodorovném směru (zleva doprava) a síla
\(F_{2}\) o
velikosti \(28{,}8\, \mathrm{N}\) 
působí ve směru svislém (shora dolů). Těleso se vlivem těchto dvou sil dá
do pohybu. Určete odchylku trajektorie tělesa od vodorovného směru.
(Výsledek zaokrouhlete na celé stupně a minuty.)}
{\(2^{\circ }10'\)}{\(3^{\circ }10'\)}{\(2^{\circ }20'\)}{\(3^{\circ }20'\)}{}{}}
\LANGsk{
\Question{
Na teleso pôsobia v jednom bode dve sily: sila
\(F_{1}\) o
velikosti \(760\, \mathrm{N}\) 
pôsobí vo vodorovnom smere (zľava doprava) a sila
\(F_{2}\) o
veľkosti \(28{,}8\, \mathrm{N}\) 
pôsobí v smere zvislom (zhora dole). Teleso sa vplyvom týchto dvoch síl dá
do pohybu. Určte odchýlku trajektórie telesa od vodorovného smeru.
(Výsledok zaokrúhlite na celé stupne a minúty.)}
{\(2^{\circ }10'\)}{\(3^{\circ }10'\)}{\(2^{\circ }20'\)}{\(3^{\circ }20'\)}{}{}}
\LANGes{
\Question{
Dos fuerzas actúan sobre el cuerpo en un punto. La fuerza \(F_{1} = 760\, \mathrm{N}\) actúa horizontalmente desde la izquiera a la derecha y la fuerza \(F_{2} = 28.8\, \mathrm{N}\) actúa verticalmente desde arriba hacia abajo. Calcula la medida del ángulo entre la dirección horizontal y la dirección de la fueza resultante y redondeando el resultado a los grados y minutos más cercanos.}
{\(2^{\circ }10'\)}{\(3^{\circ }10'\)}{\(2^{\circ }20'\)}{\(3^{\circ }20'\)}{}{}}
\End

\Data\ID{9000045701}
\Updated{2022-04-23 05:04:05}
\Level{B}
\SubArea{Triangles}
\IMG{000457_01-figure0.pdf}
\LANGen{
\Question{
The right angled triangle \(ABC\) is given (see the picture). Choose the correct statement.}
{\(\cos \beta = \frac{a} 
{c}\)}{\(\cos \beta = \frac{b} 
{c}\)}{\(\mathop{\mathrm{tg}}\nolimits \alpha = \frac{b} 
{a}\)}{\(\sin \alpha = \frac{c} 
{a}\)}{}{}}
\LANGpl{
\Question{
Dany jest trójkąt prostokątny \(ABC\) (zobacz rysunek). Znajdź prawidłową zależność między kątami i bokami trójkąta.}
{\(\cos \beta = \frac{a} 
{c}\)}{\(\cos \beta = \frac{b} 
{c}\)}{\(\mathop{\mathrm{tg}}\nolimits \alpha = \frac{b} 
{a}\)}{\(\sin \alpha = \frac{c} 
{a}\)}{}{}}
\LANGcs{
\Question{
Je dán pravoúhlý trojúhelník \(ABC\) 
(viz obrázek). Vyberte správné vyjádření hodnoty goniometrické funkce
ostrého úhlu pomocí poměru délek stran.}
{\(\cos \beta = \frac{a} 
{c}\)}{\(\cos \beta = \frac{b} 
{c}\)}{\(\mathop{\mathrm{tg}}\nolimits \alpha = \frac{b} 
{a}\)}{\(\sin \alpha = \frac{c} 
{a}\)}{}{}}
\LANGsk{
\Question{
Je daný pravouhlý trojuholník \(ABC\) 
(viď obrázok). Vyberte správne vyjadrenie hodnoty goniometrickej funkcie
ostrého uhla pomocou pomeru dĺžok strán.}
{\(\cos \beta = \frac{a} 
{c}\)}{\(\cos \beta = \frac{b} 
{c}\)}{\(\mathop{\mathrm{tg}}\nolimits \alpha = \frac{b} 
{a}\)}{\(\sin \alpha = \frac{c} 
{a}\)}{}{}}
\LANGes{
\Question{
Dado un triángulo rectángulo \(ABC\) (mira la imagen). Halla la relación válida entre los ángulos y los lados del triángulo.}
{\(\cos \beta = \frac{a} 
{c}\)}{\(\cos \beta = \frac{b} 
{c}\)}{\(\mathop{\mathrm{tg}}\nolimits \alpha = \frac{b} 
{a}\)}{\(\sin \alpha = \frac{c} 
{a}\)}{}{}}
\End

\Data\ID{9000045702}
\Updated{2020-07-15 03:07:34}
\Level{B}
\SubArea{Triangles}
\IMG{000457_02-figure0.pdf}
\LANGen{
\Question{
Given a right triangle \(ABC\) (see the picture). Find the
valid relation between the angle \(\alpha \) 
and the sides of the triangle.}
{\(\mathop{\mathrm{tg}}\nolimits \alpha = \frac{a} 
{c}\)}{\(\sin \alpha = \frac{a} 
{c}\)}{\(\cos \alpha = \frac{b} 
{a}\)}{\(\mathop{\mathrm{cotg}}\nolimits \alpha = \frac{b} 
{a}\)}{}{}}
\LANGpl{
\Question{
Dany jest trójkąt prostokątny \(ABC\) (zobacz rysunek). Znajdź prawidłową zależność pomiędzy kątem  \(\alpha \) 
a bokami tego trójkąta.}
{\(\mathop{\mathrm{tg}}\nolimits \alpha = \frac{a} 
{c}\)}{\(\sin \alpha = \frac{a} 
{c}\)}{\(\cos \alpha = \frac{b} 
{a}\)}{\(\mathop{\mathrm{cotg}}\nolimits \alpha = \frac{b} 
{a}\)}{}{}}
\LANGcs{
\Question{
Je dán pravoúhlý trojúhelník \(ABC\) 
(viz obrázek). Vyberte správné vyjádření hodnoty goniometrické funkce
ostrého úhlu pomocí poměru délek stran.}
{\(\mathop{\mathrm{tg}}\nolimits \alpha = \frac{a} 
{c}\)}{\(\sin \alpha = \frac{a} 
{c}\)}{\(\cos \alpha = \frac{b} 
{a}\)}{\(\mathop{\mathrm{cotg}}\nolimits \alpha = \frac{b} 
{a}\)}{}{}}
\LANGsk{
\Question{
Je daný pravouhlý trojuholník \(ABC\) 
(viď obrázok). Vyberte správne vyjadrenie hodnoty goniometrickej funkcie
ostrého uhla pomocou pomeru dĺžok strán.}
{\(\mathop{\mathrm{tg}}\nolimits \alpha = \frac{a} 
{c}\)}{\(\sin \alpha = \frac{a} 
{c}\)}{\(\cos \alpha = \frac{b} 
{a}\)}{\(\mathop{\mathrm{cotg}}\nolimits \alpha = \frac{b} 
{a}\)}{}{}}
\LANGes{
\Question{
Dado un triángulo rectángulo \(ABC\) (mira la imagen). Halla la relación válida entre el ángulo \(\alpha \) 
y los lados del triángulo.}
{\(\mathop{\mathrm{tg}}\nolimits \alpha = \frac{a} 
{c}\)}{\(\sin \alpha = \frac{a} 
{c}\)}{\(\cos \alpha = \frac{b} 
{a}\)}{\(\mathop{\mathrm{cotg}}\nolimits \alpha = \frac{b} 
{a}\)}{}{}}
\End

\Data\ID{9000045703}
\Updated{2020-07-15 03:07:34}
\Level{B}
\SubArea{Triangles}
\IMG{000457_03-figure0.pdf}
\LANGen{
\Question{
Given a right triangle \(ABC\) with the right angle at $C$ and an altitude $v$ (see the picture). Find the
valid relation between the angle \(\alpha \) 
and the lengths in the triangle.}
{\(\sin \alpha = \frac{v} 
{b}\)}{\(\sin \alpha = \frac{v} 
{c}\)}{\(\sin \alpha = \frac{a} 
{v}\)}{\(\sin \alpha = \frac{c} 
{a}\)}{}{}}
\LANGpl{
\Question{
Dany jest trójkąt prostokątny \(ABC\) z kątem prostym przy wierzchołku $C$ i wysokością $v$ (patrz rysunek). Znajdź prawidłową zależność pomiędzy kątem \(\alpha \) 
i długościami w tym trójkącie.}
{\(\sin \alpha = \frac{v} 
{b}\)}{\(\sin \alpha = \frac{v} 
{c}\)}{\(\sin \alpha = \frac{a} 
{v}\)}{\(\sin \alpha = \frac{c} 
{a}\)}{}{}}
\LANGcs{
\Question{
Je dán pravoúhlý trojúhelník \(ABC\) 
s pravým úhlem při vrcholu C a výškou
\(v\) (viz
obrázek). Vyberte správné vyjádření hodnoty goniometrické funkce
ostrého úhlu.}
{\(\sin \alpha = \frac{v} 
{b}\)}{\(\sin \alpha = \frac{v} 
{c}\)}{\(\sin \alpha = \frac{a} 
{v}\)}{\(\sin \alpha = \frac{c} 
{a}\)}{}{}}
\LANGsk{
\Question{
Je daný pravouhlý trojuholník \(ABC\) 
s pravým uhlom pri vrchole C a výškou
\(v\) (viď
obrázok). Vyberte správne vyjadrenie hodnoty goniometrickej funkcie
ostrého uhla.}
{\(\sin \alpha = \frac{v} 
{b}\)}{\(\sin \alpha = \frac{v} 
{c}\)}{\(\sin \alpha = \frac{a} 
{v}\)}{\(\sin \alpha = \frac{c} 
{a}\)}{}{}}
\LANGes{
\Question{
Dado un triángulo rectángulo \(ABC\), siendo $C$ el vértice del ángulo recto, con la altura $v$ (mira la imagen). Halla la relación válida entre el ángulo \(\alpha \) 
y las longitudes en el triángulo.}
{\(\sin \alpha = \frac{v} 
{b}\)}{\(\sin \alpha = \frac{v} 
{c}\)}{\(\sin \alpha = \frac{a} 
{v}\)}{\(\sin \alpha = \frac{c} 
{a}\)}{}{}}
\End

\Data\ID{9000045704}
\Updated{2020-07-15 03:07:34}
\Level{B}
\SubArea{Triangles}
\IMG{000457_04-figure0.pdf}
\LANGen{
\Question{
Given a right triangle \(ABC\) with the right angle at $C$ and an altitude $v$ (see the picture). Find the
valid relation between the angle \(\beta \) 
and the lengths in the triangle.}
{\(\sin \beta = \frac{v} 
{a}\)}{\(\mathop{\mathrm{tg}}\nolimits \beta = \frac{a} 
{v}\)}{\(\cos \beta = \frac{v} 
{a}\)}{\(\mathop{\mathrm{tg}}\nolimits \beta = \frac{v} 
{a}\)}{}{}}
\LANGpl{
\Question{
Dany jest trójkąt prostokątny \(ABC\) z kątem prostym przy wierzchołku $C$ i wysokością $v$ (patrz rysunek). Znajdź prawidłową zależność pomiędzy kątem  \(\beta \) 
i długościami w tym trójkącie.}
{\(\sin \beta = \frac{v} 
{a}\)}{\(\mathop{\mathrm{tg}}\nolimits \beta = \frac{a} 
{v}\)}{\(\cos \beta = \frac{v} 
{a}\)}{\(\mathop{\mathrm{tg}}\nolimits \beta = \frac{v} 
{a}\)}{}{}}
\LANGcs{
\Question{
Je dán pravoúhlý trojúhelník \(ABC\) 
s pravým úhlem při vrcholu C a výškou
\(v\) (viz
obrázek). Vyberte správné vyjádření hodnoty goniometrické funkce
ostrého úhlu.}
{\(\sin \beta = \frac{v} 
{a}\)}{\(\mathop{\mathrm{tg}}\nolimits \beta = \frac{a} 
{v}\)}{\(\cos \beta = \frac{v} 
{a}\)}{\(\mathop{\mathrm{tg}}\nolimits \beta = \frac{v} 
{a}\)}{}{}}
\LANGsk{
\Question{
Je daný pravouhlý trojuholník \(ABC\) 
s pravým uhlom pri vrchole C a výškou
\(v\) (viď
obrázok). Vyberte správne vyjadrenie hodnoty goniometrickej funkcie
ostrého uhla.}
{\(\sin \beta = \frac{v} 
{a}\)}{\(\mathop{\mathrm{tg}}\nolimits \beta = \frac{a} 
{v}\)}{\(\cos \beta = \frac{v} 
{a}\)}{\(\mathop{\mathrm{tg}}\nolimits \beta = \frac{v} 
{a}\)}{}{}}
\LANGes{
\Question{
Dado un triángulo rectángulo \(ABC\), siendo $C$ el vértice del ángulo recto, con la altura $v$ (mira la imagen). Halla la relación válida entre el ángulo \(\beta \) 
y las longitudes en el triángulo.}
{\(\sin \beta = \frac{v} 
{a}\)}{\(\mathop{\mathrm{tg}}\nolimits \beta = \frac{a} 
{v}\)}{\(\cos \beta = \frac{v} 
{a}\)}{\(\mathop{\mathrm{tg}}\nolimits \beta = \frac{v} 
{a}\)}{}{}}
\End

\Data\ID{9000046403}
\Updated{2020-07-15 03:07:34}
\Level{B}
\SubArea{Triangles}
\LANGen{
\Question{
Consider an isosceles triangle, i.e. the triangle with two
sides of equal length. The length of the third side is
\(4\, \mathrm{cm}\). One of the
interior angles is \(120^{\circ }\).
Find the area of this triangle.}
{\(\frac{4\sqrt{3}}
 {3} \, \mathrm{cm}^{2}\)}{\(4\sqrt{3}\, \mathrm{cm}^{2}\)}{\(\frac{8\sqrt{3}}
 {3} \, \mathrm{cm}^{2}\)}{}{}{}}
\LANGpl{
\Question{
Rozważ trójkąt równoramienny, tzn. taki, w którym dwa boki są jednakowej długości. Długość trzeciego boku wynosi 
\(4\, \mathrm{cm}\). Jeden z kątów wewnętrznych jest równy \(120^{\circ }\).
Oblicz pole powierzchni tego trójkąta.}
{\(\frac{4\sqrt{3}}
 {3} \, \mathrm{cm}^{2}\)}{\(4\sqrt{3}\, \mathrm{cm}^{2}\)}{\(\frac{8\sqrt{3}}
 {3} \, \mathrm{cm}^{2}\)}{}{}{}}
\LANGcs{
\Question{
Určete obsah rovnoramenného trojúhelníku se základnou délky
\(4\, \mathrm{cm}\) a vnitřním
úhlem \(120^{\circ }\).}
{\(\frac{4\sqrt{3}}
 {3} \, \mathrm{cm}^{2}\)}{\(4\sqrt{3}\, \mathrm{cm}^{2}\)}{\(\frac{8\sqrt{3}}
 {3} \, \mathrm{cm}^{2}\)}{}{}{}}
\LANGsk{
\Question{
Určte obsah rovnoramenného trojuholníka so základňou dĺžky
\(4\, \mathrm{cm}\) a vnútorným
uhlom \(120^{\circ }\).}
{\(\frac{4\sqrt{3}}
 {3} \, \mathrm{cm}^{2}\)}{\(4\sqrt{3}\, \mathrm{cm}^{2}\)}{\(\frac{8\sqrt{3}}
 {3} \, \mathrm{cm}^{2}\)}{}{}{}}
\LANGes{
\Question{
Dado un triángulo isósceles. El tercer lado mide
\(4\, \mathrm{cm}\). Uno de los ángulos interiores mide \(120^{\circ }\).
Calcula el área del triángulo.}
{\(\frac{4\sqrt{3}}
 {3} \, \mathrm{cm}^{2}\)}{\(4\sqrt{3}\, \mathrm{cm}^{2}\)}{\(\frac{8\sqrt{3}}
 {3} \, \mathrm{cm}^{2}\)}{}{}{}}
\End

\Data\ID{1003021902}
\Updated{2022-08-25 10:08:24}
\Level{C}
\SubArea{Triangles}
\LANGen{
\Question{
What is the width of a computer screen if the ratio of its width and height is \( 16:9 \) and the computer has \( 23 \)-inch monitor? Round the result to two decimal places. (\( 1 \) inch=\( 2.54\,\mathrm{cm} \))}
{\( 50.92\,\mathrm{cm} \)}{\( 20.05\,\mathrm{cm} \)}{\( 11.28\,\mathrm{cm} \)}{\( 28.65\,\mathrm{cm} \)}{}{}}
\LANGpl{
\Question{
Jaka jest szerokość  ekranu komputera, jeśli stosunek jego szerokości i wysokości wynosi \( 16:9 \), a komputer ma  \( 23 \) calowy monitor? Zaokrąglij do dwóch miejsc po przecinku. (\( 1 \) inch=\( 2{,}54\,\mathrm{cm} \))}
{\( 50{,}92\,\mathrm{cm} \)}{\( 20{,}05\,\mathrm{cm} \)}{\( 11{,}28\,\mathrm{cm} \)}{\( 28{,}65\,\mathrm{cm} \)}{}{}}
\LANGcs{
\Question{
Jakou šířku má monitor počítače, jestliže je poměr šířky a výšky monitoru \( 16:9 \) a monitor je \( 23 \) palcový? Výsledek zaokrouhlete na dvě desetinná místa. (\( 1 \) palec=\( 2{,}54\,\mathrm{cm} \))}
{\( 50{,}92\,\mathrm{cm} \)}{\( 20{,}05\,\mathrm{cm} \)}{\( 11{,}28\,\mathrm{cm} \)}{\( 28{,}65\,\mathrm{cm} \)}{}{}}
\LANGsk{
\Question{
Akú šírku má monitor počítača, ak pomer šírky a výšky monitora je  \( 16:9 \) a monitor je  \( 23 \) palcový? Výsledok zaokrúhlite na dve desatinné miesta. (\( 1 \) inch=\( 2{,}54\,\mathrm{cm} \))}
{\( 50{,}92\,\mathrm{cm} \)}{\( 20{,}05\,\mathrm{cm} \)}{\( 11{,}28\,\mathrm{cm} \)}{\( 28{,}65\,\mathrm{cm} \)}{}{}}
\LANGes{
\Question{
¿Cuál es el ancho de la pantalla de un ordenador si su altura y anchura están en proporción \( 16:9 \) y el tamaño de la pantalla son \( 23 \) pulgadas en diagonal? Redondea el resultado a dos decimales. (\( 1 \) pulgada=\( 2.54\,\mathrm{cm} \))}
{\( 50.92\,\mathrm{cm} \)}{\( 20.05\,\mathrm{cm} \)}{\( 11.28\,\mathrm{cm} \)}{\( 28.65\,\mathrm{cm} \)}{}{}}
\End

\Data\ID{1003021905}
\Updated{2022-08-25 10:08:29}
\Level{C}
\SubArea{Triangles}
\LANGen{
\Question{
Calculate the height between two floors if you know that the number of stairs between the floors is \( 16 \), the slope of the staircase is \( 30^{\circ} \) and the depth of a stair is \( 25\,\mathrm{cm} \).}
{\( \frac{400}{\sqrt3}\,\mathrm{cm} \)}{\( \frac{25}{\sqrt3}\,\mathrm{cm} \)}{\( 200\,\mathrm{cm} \)}{\( 400\,\mathrm{cm} \)}{}{}}
\LANGpl{
\Question{
Oblicz wysokość pomiędzy dwiema piętrami wiedząc, że liczba schodów pomiędzy piętrami wynosi  \( 16 \), nachylenie klatki schodowej jest równe  \( 30^{\circ} \), a głębokość schodów ma \( 25\,\mathrm{cm} \).}
{\( \frac{400}{\sqrt3}\,\mathrm{cm} \)}{\( \frac{25}{\sqrt3}\,\mathrm{cm} \)}{\( 200\,\mathrm{cm} \)}{\( 400\,\mathrm{cm} \)}{}{}}
\LANGcs{
\Question{
Určete výšku mezi dvěma poschodími, jestliže víte, že počet schodů mezi dvěma poschodími je  \( 16 \), sklon stoupání \( 30^{\circ} \) a hloubka schodu \( 25\,\mathrm{cm} \).}
{\( \frac{400}{\sqrt3}\,\mathrm{cm} \)}{\( \frac{25}{\sqrt3}\,\mathrm{cm} \)}{\( 200\,\mathrm{cm} \)}{\( 400\,\mathrm{cm} \)}{}{}}
\LANGsk{
\Question{
Určte výšku medzi dvoma poschodiami, ak viete, že počet schodov medzi dvoma poschodiami je  \( 16 \), sklon stúpania  \( 30^{\circ} \) a šírka schodu  \( 25\,\mathrm{cm} \).}
{\( \frac{400}{\sqrt3}\,\mathrm{cm} \)}{\( \frac{25}{\sqrt3}\,\mathrm{cm} \)}{\( 200\,\mathrm{cm} \)}{\( 400\,\mathrm{cm} \)}{}{}}
\LANGes{
\Question{
Calcula la altura entre dos pisos si sabemos que hay \( 16 \) escalones entre piso y piso, la pendiente de la escalera es \( 30^{\circ} \) y la profundidad de un escalón mide \( 25\,\mathrm{cm} \).}
{\( \frac{400}{\sqrt3}\,\mathrm{cm} \)}{\( \frac{25}{\sqrt3}\,\mathrm{cm} \)}{\( 200\,\mathrm{cm} \)}{\( 400\,\mathrm{cm} \)}{}{}}
\End

\Data\ID{1003076708}
\Updated{2021-08-29 09:08:25}
\Level{C}
\SubArea{Triangles}
\LANGen{
\Question{
The measures of the interior angles of a  triangle are \( 30^{\circ} \), \( 45^{\circ} \) and \( 105^{\circ} \). The length of its longest side is \( 10\,\mathrm{cm} \). The length of its shortest side is:}
{\( 5.18\,\mathrm{cm} \)}{\( 7.33\,\mathrm{cm} \)}{\( 5.01\,\mathrm{cm} \)}{\( 7.07\,\mathrm{cm} \)}{}{}}
\LANGpl{
\Question{
Miary kątów wewnętrznych trójkąta są równe  \( 30^{\circ} \), \( 45^{\circ} \) i \( 105^{\circ} \). Długość najdłuższego boku wynosi \( 10\,\mathrm{cm} \). Długość najkrótszego boku jest równa:}
{\( 5{,}18\,\mathrm{cm} \)}{\( 7{,}33\,\mathrm{cm} \)}{\( 5{,}01\,\mathrm{cm} \)}{\( 7{,}07\,\mathrm{cm} \)}{}{}}
\LANGcs{
\Question{
Vnitřní úhly trojúhelníku mají velikost \( 30^{\circ} \), \( 45^{\circ} \) a \( 105^{\circ} \). Jeho nejdelší strana měří \( 10\,\mathrm{cm} \). Nejkratší strana trojúhelníku měří:}
{\( 5{,}18\,\mathrm{cm} \)}{\( 7{,}33\,\mathrm{cm} \)}{\( 5{,}01\,\mathrm{cm} \)}{\( 7{,}07\,\mathrm{cm} \)}{}{}}
\LANGsk{
\Question{
Vnútorné uhly trojuholníka majú veľkosti  \( 30^{\circ} \), \( 45^{\circ} \) a \( 105^{\circ} \). Jeho najdlhšia strana meria  \( 10\,\mathrm{cm} \). Najkratšia strana trojuholníka meria:}
{\( 5{,}18\,\mathrm{cm} \)}{\( 7{,}33\,\mathrm{cm} \)}{\( 5{,}01\,\mathrm{cm} \)}{\( 7{,}07\,\mathrm{cm} \)}{}{}}
\LANGes{
\Question{
Las medidas de los ángulos interiores de un triángulo son  \( 30^{\circ} \), \( 45^{\circ} \) y \( 105^{\circ} \).  La longitud de su lado más largo es\( 10\,\mathrm{cm} \). La longitud de su lado más corto es:}
{\( 5.18\,\mathrm{cm} \)}{\( 7.33\,\mathrm{cm} \)}{\( 5.01\,\mathrm{cm} \)}{\( 7.07\,\mathrm{cm} \)}{}{}}
\End

\Data\ID{1003076710}
\Updated{2021-08-29 09:08:52}
\Level{C}
\SubArea{Triangles}
\LANGen{
\Question{
\( ABC \) is a triangle where the side \( b \) is \( 74\,\mathrm{cm} \) long and the angle \( \alpha = 60^{\circ} \). Find the length of its side \( c \)  if you know that the area of the triangle is \( 720.9\,\mathrm{cm}^2 \).}
{\( 22.5\,\mathrm{cm} \)}{\( 37.56\,\mathrm{cm} \)}{\( 38.97\,\mathrm{cm} \)}{\( 24.54\,\mathrm{cm} \)}{}{}}
\LANGpl{
\Question{
Dany jest trójkąt \( ABC \) o boku  \( b \) równym \( 74\,\mathrm{cm} \) i kącie \( \alpha = 60^{\circ} \). Oblicz długość boku \( c \)  wiedząc, że pole powierzchni trójkąta wynosi \( 720{,}9\,\mathrm{cm}^2 \).}
{\( 22{,}5\,\mathrm{cm} \)}{\( 37{,}56\,\mathrm{cm} \)}{\( 38{,}97\,\mathrm{cm} \)}{\( 24{,}54\,\mathrm{cm} \)}{}{}}
\LANGcs{
\Question{
Jaká je velikost strany c v trojúhelníku \( ABC \), jestliže jeho obsah je  \( 720{,}9\,\mathrm{cm}^2 \), délka strany \( b \) je \( 74\,\mathrm{cm} \) a úhel \( \alpha = 60^{\circ} \)?}
{\( 22{,}5\,\mathrm{cm} \)}{\( 37{,}56\,\mathrm{cm} \)}{\( 38{,}97\,\mathrm{cm} \)}{\( 24{,}54\,\mathrm{cm} \)}{}{}}
\LANGsk{
\Question{
Aká je veľkosť strany c v trojuholníku \( ABC \), ak jeho obsah je  \( 720{,}9\,\mathrm{cm}^2 \), dĺžka strany \( b \) je \( 74\,\mathrm{cm} \) a uhol \( \alpha = 60^{\circ} \)?}
{\( 22{,}5\,\mathrm{cm} \)}{\( 37{,}56\,\mathrm{cm} \)}{\( 38{,}97\,\mathrm{cm} \)}{\( 24{,}54\,\mathrm{cm} \)}{}{}}
\LANGes{
\Question{
\( ABC \) es un triángulo donde el lado \( b \)  es  \( 74\,\mathrm{cm} \) largo y el ángulo \( \alpha = 60^{\circ} \). Calcula la longitud de su lado \( c \)  si sabes que el área del triángulo es \( 720.9\,\mathrm{cm}^2 \).}
{\( 22.5\,\mathrm{cm} \)}{\( 37.56\,\mathrm{cm} \)}{\( 38.97\,\mathrm{cm} \)}{\( 24.54\,\mathrm{cm} \)}{}{}}
\End

\Data\ID{1103021901}
\Updated{2022-04-23 05:04:05}
\Level{C}
\SubArea{Triangles}
\IMG{1103021901.pdf}
\LANGen{
\Question{
What is the slope of a staircase if the height of a stair is \( 15\,\mathrm{cm} \) and the depth of a stair is \( 30\,\mathrm{cm} \)?}
{\( 26.57^{\circ} \)}{\( 63.43^{\circ} \)}{\( 30^{\circ} \)}{\( 60^{\circ} \)}{}{}}
\LANGpl{
\Question{
Jakie jest nachylenie klatki schodowej, jeśli wysokość każdego schodu jest \( 15\,\mathrm{cm} \) i głębokość \( 30\,\mathrm{cm} \)?}
{\( 26{,}57^{\circ} \)}{\( 63{,}43^{\circ} \)}{\( 30^{\circ} \)}{\( 60^{\circ} \)}{}{}}
\LANGcs{
\Question{
Pod jakým úhlem stoupá schodiště, jehož schody jsou \( 30\,\mathrm{cm} \) hluboké a  \( 15\,\mathrm{cm} \) vysoké?}
{\( 26{,}57^{\circ} \)}{\( 63{,}43^{\circ} \)}{\( 30^{\circ} \)}{\( 60^{\circ} \)}{}{}}
\LANGsk{
\Question{
Pod akým uhlom stúpa schodište, ktorého schody sú \( 30\,\mathrm{cm} \) široké a  \( 15\,\mathrm{cm} \) vysoké?}
{\( 26{,}57^{\circ} \)}{\( 63{,}43^{\circ} \)}{\( 30^{\circ} \)}{\( 60^{\circ} \)}{}{}}
\LANGes{
\Question{
¿Cuál es la pendiente de una escalera si la altura de un escalón es \( 15\,\mathrm{cm} \) y la profundidad es \( 30\,\mathrm{cm} \)?}
{\( 26.57^{\circ} \)}{\( 63.43^{\circ} \)}{\( 30^{\circ} \)}{\( 60^{\circ} \)}{}{}}
\End

\Data\ID{1103021903}
\Updated{2021-04-19 06:04:08}
\Level{C}
\SubArea{Triangles}
\IMG{1103021903.pdf}
\LANGen{
\Question{
An observer was watching an approaching plane flying at a height of \( 3000\,\mathrm{m} \) in a straight line with constant velocity.  At the first moment of measurement the observer saw the plane to be at an angle of elevation of \( 25^{\circ} \). After \( 10 \) seconds the angle of elevation changed to \( 35^{\circ} \). What was the speed of the plane? Round the result to ones.}
{\( 215\,\mathrm{m}\cdot\mathrm{s}^{-1} \)}{\( 2149\,\mathrm{m}\cdot\mathrm{s}^{-1} \)}{\( 6576\,\mathrm{m}\cdot\mathrm{s}^{-1} \)}{\( 658\,\mathrm{m}\cdot\mathrm{s}^{-1} \)}{}{}}
\LANGpl{
\Question{
Obserwator patrzy na zbliżający się samolot lecący na wysokości \( 3000\,\mathrm{m} \) w linii prostej ze stałą prędkością. W pierwszej chwili zauważył, ze kąt wzniesienia samolotu wynosi \( 25^{\circ} \). Po \( 10 \) sekundach kąt wzniesienia zmienił się na \( 35^{\circ} \). Jaka była prędkość samolotu? Zaokrąglij do jedności.}
{\( 215\,\mathrm{m}\cdot\mathrm{s}^{-1} \)}{\( 2149\,\mathrm{m}\cdot\mathrm{s}^{-1} \)}{\( 6576\,\mathrm{m}\cdot\mathrm{s}^{-1} \)}{\( 658\,\mathrm{m}\cdot\mathrm{s}^{-1} \)}{}{}}
\LANGcs{
\Question{
Pozorovatel sleduje blížící se letadlo letící konstantní rychlostí po přímce ve výšce \( 3000\,\mathrm{m} \).  V prvním okamžiku spatří pozorovatel letadlo ve výškovém úhlu \( 25^{\circ} \). Po \( 10 \) sekundách se výškový úhel změní na \( 35^{\circ} \). Jakou rychlostí se letadlo pohybovalo? Zaokrouhlete na jednotky.}
{\( 215\,\mathrm{m}\cdot\mathrm{s}^{-1} \)}{\( 2149\,\mathrm{m}\cdot\mathrm{s}^{-1} \)}{\( 6576\,\mathrm{m}\cdot\mathrm{s}^{-1} \)}{\( 658\,\mathrm{m}\cdot\mathrm{s}^{-1} \)}{}{}}
\LANGsk{
\Question{
Pozorovateľ sleduje približujúce sa lietalo, ktoré letí vo výške  \( 3000\,\mathrm{m} \) rovnomerným priamočiarym pohybom. Pri prvom meraní pozorovateľ videl lietadlo pod výškovým uhlom \( 25^{\circ} \). Po uplynutí  \( 10 \) sekúnd sa výškový uhol zmenil na   \( 35^{\circ} \). Akou rýchlosťou letelo lietadlo? Výsledok zaokrúhlite na jednotky.}
{\( 215\,\mathrm{m}\cdot\mathrm{s}^{-1} \)}{\( 2149\,\mathrm{m}\cdot\mathrm{s}^{-1} \)}{\( 6576\,\mathrm{m}\cdot\mathrm{s}^{-1} \)}{\( 658\,\mathrm{m}\cdot\mathrm{s}^{-1} \)}{}{}}
\LANGes{
\Question{
Un observador observaba un avión que se aproximaba volando, a una altura de \( 3000\,\mathrm{m} \), en línea recta con velocidad constante.  En un primer momento, el observador vio que el avión estaba en un ángulo de elevación de \( 25^{\circ} \). Después de \( 10 \) segundos, el ángulo de elevación cambió a \( 35^{\circ} \). Calcula la velocidad del avión. Redondea el resultado a las unidades.}
{\( 215\,\mathrm{m}\cdot\mathrm{s}^{-1} \)}{\( 2149\,\mathrm{m}\cdot\mathrm{s}^{-1} \)}{\( 6576\,\mathrm{m}\cdot\mathrm{s}^{-1} \)}{\( 658\,\mathrm{m}\cdot\mathrm{s}^{-1} \)}{}{}}
\End

\Data\ID{1103021904}
\Updated{2021-04-19 06:04:16}
\Level{C}
\SubArea{Triangles}
\IMG{1103021904.pdf}
\LANGen{
\Question{
From the highest window of Orava Castle, the angles of depression to the banks of the Orava river are \( 60^{\circ} \) and \( 20^{\circ} \). The height of the window above the river is \( 50\,\mathrm{m} \). What is the width of the river?}
{\( 108.5\,\mathrm{m} \)}{\( 137.4\,\mathrm{m} \)}{\( 100.5\,\mathrm{m} \)}{\( 125.4\,\mathrm{m} \)}{}{}}
\LANGpl{
\Question{
Z najwyższego okna Zamku Orawskiego możemy obserwować brzegi rzeki Orawy pod kątami depresji \( 60^{\circ} \) i \( 20^{\circ} \). Wysokość okna nad rzeką jest równa \( 50\,\mathrm{m} \). Jaka jest szerokość rzeki?}
{\( 108{,}5\,\mathrm{m} \)}{\( 137{,}4\,\mathrm{m} \)}{\( 100{,}5\,\mathrm{m} \)}{\( 125{,}4\,\mathrm{m} \)}{}{}}
\LANGcs{
\Question{
Z nejvyššího okna Oravského hradu je vidět na břehy řeky Oravy v hloubkových úhlech \( 60^{\circ} \) a \( 20^{\circ} \). Výška okna nad hladinou Oravy je \( 50\,\mathrm{m} \). Jak je řeka široká?}
{\( 108{,}5\,\mathrm{m} \)}{\( 137{,}4\,\mathrm{m} \)}{\( 100{,}5\,\mathrm{m} \)}{\( 125{,}4\,\mathrm{m} \)}{}{}}
\LANGsk{
\Question{
Z najvyššieho okna Oravského hradu vidno brehy rieky Orava v hĺbkových uhloch  \( 60^{\circ} \) a \( 20^{\circ} \). Výška okna nad hladinou Oravy je \( 50\,\mathrm{m} \). Akú šírku má rieka?}
{\( 108{,}5\,\mathrm{m} \)}{\( 137{,}4\,\mathrm{m} \)}{\( 100{,}5\,\mathrm{m} \)}{\( 125{,}4\,\mathrm{m} \)}{}{}}
\LANGes{
\Question{
Desde la ventana más alta del castillo de Orava, los ángulos de depresión a las orillas del río Oravice son \( 60^{\circ} \) y \( 20^{\circ} \). La altura de la ventana sobre el río es \( 50\,\mathrm{m} \). ¿Cuál es el ancho del río?}
{\( 108.5\,\mathrm{m} \)}{\( 137.4\,\mathrm{m} \)}{\( 100.5\,\mathrm{m} \)}{\( 125.4\,\mathrm{m} \)}{}{}}
\End

\Data\ID{1103021906}
\Updated{2021-04-19 06:04:58}
\Level{C}
\SubArea{Triangles}
\IMG{1103021906.pdf}
\LANGen{
\Question{
The distance between the places \( A \) and \( C \) on a straight road is \( 300\,\mathrm{m} \). There is a balloon \( B \) above the road between the places \( A \) and \( C \). See the picture. The angles of elevation from the places \( A \) and \( C \) to the balloon \( B \) are \( 20^{\circ} \) and \( 40^{\circ} \) respectively. What is the height \( h \) of the balloon?}
{\( 76\,\mathrm{m} \)}{\( 168\,\mathrm{m} \)}{\( 488\,\mathrm{m} \)}{\( 523\,\mathrm{m} \)}{}{}}
\LANGpl{
\Question{
Odległość pomiędzy miejscami  \( A \) i \( C \) w drodze prostej wynosi \( 300\,\mathrm{m} \). Nad drogą znajduje się balon \( B \) pomiędzy miejscami \( A \) i \( C \). Patrrz rysunek. Kąty wzniesienia pomiędzy miejscami \( A \) i \( C \) oraz balonu \( B \) mają kolejno \( 20^{\circ} \) i \( 40^{\circ} \). Jaka jest wysokość  \( h \) na której znajduje się balon?}
{\( 76\,\mathrm{m} \)}{\( 168\,\mathrm{m} \)}{\( 488\,\mathrm{m} \)}{\( 523\,\mathrm{m} \)}{}{}}
\LANGcs{
\Question{
Vzdálenost míst \( A \) a \( C \) na rovné cestě je  \( 300\,\mathrm{m} \). Mezi místy \( A \) a \( C \) se nad cestou vznáší balón \( B \) (viz obrázek). Z místa \( A \) je možné pozorovat balón  \( B \) pod výškovým úhlem \( 20^{\circ} \), z místa \( C \)  pod výškovým úhlem \( 40^{\circ} \).  Určete, v jaké výšce \( h \), zaokrouhleno na celé metry, se vznáší balón nad cestou.}
{\( 76\,\mathrm{m} \)}{\( 168\,\mathrm{m} \)}{\( 488\,\mathrm{m} \)}{\( 523\,\mathrm{m} \)}{}{}}
\LANGsk{
\Question{
Vzdialenosť miest  \( A \) a \( C \) na rovnej ceste je  \( 300\,\mathrm{m} \). Medzi miestami \( A \) a \( C \) sa nad cestou vznáša balón \( B \). Pozri obrázok. Z miesta  \( A \) je možné pozorovať balón  \( B \) pod výškovým uhlom \( 20^{\circ} \), z miesta \( C \)  pod výškovým uhlom \( 40^{\circ} \).  Určte zaokrúhlene na celé metre, v akej výške \( h \)  sa vznáša balón nad cestou.}
{\( 76\,\mathrm{m} \)}{\( 168\,\mathrm{m} \)}{\( 488\,\mathrm{m} \)}{\( 523\,\mathrm{m} \)}{}{}}
\LANGes{
\Question{
La distancia entre los puntos \( A \) y \( C \) en línea recta es \( 300\,\mathrm{m} \). Hay un globo \( B \) sobre el camino entre los puntos \( A \) y \( C \). Mira la imagen. Los ángulos de elevación desde los puntos \( A \) y \( C \) respecto al globo \( B \) son \( 20^{\circ} \) y \( 40^{\circ} \) respectivamente. ¿Cuál es la altura \( h \) del globo?}
{\( 76\,\mathrm{m} \)}{\( 168\,\mathrm{m} \)}{\( 488\,\mathrm{m} \)}{\( 523\,\mathrm{m} \)}{}{}}
\End

\Data\ID{1103021907}
\Updated{2021-04-19 06:04:14}
\Level{C}
\SubArea{Triangles}
\IMG{1103021907.pdf}
\LANGen{
\Question{
A plane is flying at a speed of \( 900\,\mathrm{km}\cdot\mathrm{h}^{-1} \) and according to the compass the axis of the airplane is still heading westbound. What angle will the flight path of the airplane make with the east-west direction if the south wind starts to blow at a speed of \( 10\,\mathrm{m}\cdot\mathrm{s}^{-1} \)? Round the result to two decimal places.}
{\( 2.29^{\circ} \)}{\( 0.64^{\circ} \)}{\( 0.01^{\circ} \)}{\( 87.71^{\circ} \)}{}{}}
\LANGpl{
\Question{
Samolot porusza się z prędkością \( 900\,\mathrm{km}\cdot\mathrm{h}^{-1} \) i  w odniesieniu do wskazówki kompasu zmierza na zachód. Jaki kąt utworzy trasa lotu samolotu z kierunkiem wschód-zachód, jeśli południowy wiatr zacznie wiać z prędkością \( 10\,\mathrm{m}\cdot\mathrm{s}^{-1} \)? Zaokrąglij do dwóch miejsc po przecinku.}
{\( 2{,}29^{\circ} \)}{\( 0{,}64^{\circ} \)}{\( 0{,}01^{\circ} \)}{\( 87{,}71^{\circ} \)}{}{}}
\LANGcs{
\Question{
Letadlo letí rychlostí \( 900\,\mathrm{km}\cdot\mathrm{h}^{-1} \) a podle kompasu směřuje osa letadla na západ. Jaký úhel bude svírat dráha letadla ke směru východ-západ, jestliže začne foukat jižní vítr rychlostí \( 10\,\mathrm{m}\cdot\mathrm{s}^{-1} \)? Výsledek zaokrouhlete na dvě desetinná místa.}
{\( 2{,}29^{\circ} \)}{\( 0{,}64^{\circ} \)}{\( 0{,}01^{\circ} \)}{\( 87{,}71^{\circ} \)}{}{}}
\LANGsk{
\Question{
Lietadlo letí rýchlosťou  \( 900\,\mathrm{km}\cdot\mathrm{h}^{-1} \) a podľa kompasu os lietadla smeruje stále na západ. Aký uhol bude zvierať dráha lietadla so smerom východ-západ, ak začne fúkať južný vietor rýchlosťou \( 10\,\mathrm{m}\cdot\mathrm{s}^{-1} \)? Výsledok zaokrúhlite na dve desatinné miesta.}
{\( 2{,}29^{\circ} \)}{\( 0{,}64^{\circ} \)}{\( 0{,}01^{\circ} \)}{\( 87{,}71^{\circ} \)}{}{}}
\LANGes{
\Question{
Un avión vuela a una velocidad de \( 900\,\mathrm{km}\cdot\mathrm{h}^{-1} \) y, según la brújula. el avión se dirige hacia el oeste. ¿Qué ángulo formará la trayectoria del vuelo del avión con la dirección este-oeste si el viento del sur comienza a soplar a una velocidad de \( 10\,\mathrm{m}\cdot\mathrm{s}^{-1} \)? Redondea el resultado a dos decimales.}
{\( 2.29^{\circ} \)}{\( 0.64^{\circ} \)}{\( 0.01^{\circ} \)}{\( 87.71^{\circ} \)}{}{}}
\End

\Data\ID{2010015204}
\Updated{2022-08-25 10:08:04}
\Level{C}
\SubArea{Triangles}
\LANGen{
\Question{
What is the height of a computer screen if the ratio of its width and height is \( 16:9 \) and the computer has \( 23 \)-inch monitor? Round the result to two decimal places. (\( 1 \) inch=\( 2.54\,\mathrm{cm} \))}
{\( 28.64\,\mathrm{cm} \)}{\(50.92\,\mathrm{cm} \)}{\( 20.05\,\mathrm{cm} \)}{\(11.28\,\mathrm{cm} \)}{}{}}
\LANGpl{
\Question{
Jaka jest wysokość ekranu komputera, jeśli stosunek jego szerokości do wysokości wynosi \( 16:9 \), a komputer ma \( 23 \)-calowy monitor? Zaokrąglij wynik do dwóch miejsc po przecinku. (\( 1 \) cal=\( 2{,}54\,\mathrm{cm} \))}
{\( 28{,}64\,\mathrm{cm} \)}{\(50{,}92\,\mathrm{cm} \)}{\( 20{,}05\,\mathrm{cm} \)}{\(11{,}28\,\mathrm{cm} \)}{}{}}
\LANGcs{
\Question{
Jaká je výška obrazovky monitoru s poměrem stran \( 16:9 \) a úhlopříčkou o délce \( 23 \) palců? Výsledek zaokrouhlete na 2 desetinná místa. (\( 1 \) palec=\( 2{,}54\,\mathrm{cm} \))}
{\( 28{,}64\,\mathrm{cm} \)}{\(50{,}92\,\mathrm{cm} \)}{\( 20{,}05\,\mathrm{cm} \)}{\(11{,}28\,\mathrm{cm} \)}{}{}}
\LANGsk{
\Question{
Aká je výška obrazovky monitora, ak pomer jej šírky a výšky je \( 16:9 \), uhlopriečka má dĺžku \( 23 \) palcov? Výsledok zaokrúhlite na dve desatinné miesta. (\( 1 \) palec=\( 2{,}54\,\mathrm{cm} \))}
{\( 28{,}64\,\mathrm{cm} \)}{\(50{,}92\,\mathrm{cm} \)}{\( 20{,}05\,\mathrm{cm} \)}{\(11{,}28\,\mathrm{cm} \)}{}{}}
\LANGes{
\Question{
¿Cuál es la altura de la pantalla de un ordenador si su altura y anchura están en proporción \( 16:9 \) y el tamaño de la pantalla son \( 23 \) pulgadas en diagonal? Redondea el resultado a dos decimales. (\( 1 \) pulgada=\( 2.54\,\mathrm{cm} \))}
{\( 28.64\,\mathrm{cm} \)}{\(50.92\,\mathrm{cm} \)}{\( 20.05\,\mathrm{cm} \)}{\(11.28\,\mathrm{cm} \)}{}{}}
\End

\Data\ID{2010015302}
\Updated{2022-08-25 10:08:04}
\Level{C}
\SubArea{Triangles}
\IMG{2010015301-figure1.pdf}
\LANGen{
\Question{
The slope of a staircase is \(30^{\circ}\).  What is the depth of a stair if its height is \( 15\,\mathrm{cm} \)? (See the picture.)}
{\(15\sqrt3\,\mathrm{cm} \)}{\(30\,\mathrm{cm} \)}{\( \frac{15}{\sqrt3}\,\mathrm{cm} \)}{\( 28\,\mathrm{cm} \)}{}{}}
\LANGpl{
\Question{
Nachylenie schodów wynosi \(30^{\circ}\). Jaka jest głębokość schodów, jeśli ich wysokość to \( 15\,\mathrm{cm} \)? (Popatrz na rysunek.)}
{\(15\sqrt3\,\mathrm{cm} \)}{\(30\,\mathrm{cm} \)}{\( \frac{15}{\sqrt3}\,\mathrm{cm} \)}{\( 28\,\mathrm{cm} \)}{}{}}
\LANGcs{
\Question{
Sklon schodiště je \(30^{\circ}\).  Jaká je hloubka jednoho schodu, jestliže je \( 15\,\mathrm{cm} \) vysoký? (Viz obrázek.)}
{\(15\sqrt3\,\mathrm{cm} \)}{\(30\,\mathrm{cm} \)}{\( \frac{15}{\sqrt3}\,\mathrm{cm} \)}{\( 28\,\mathrm{cm} \)}{}{}}
\LANGsk{
\Question{
Schodište stúpa pod uhlom \(30^{\circ}\).  Akú šírku majú schody, ak ich výška je \( 15\,\mathrm{cm} \)? (Pozri obrázok.)}
{\(15\sqrt3\,\mathrm{cm} \)}{\(30\,\mathrm{cm} \)}{\( \frac{15}{\sqrt3}\,\mathrm{cm} \)}{\( 28\,\mathrm{cm} \)}{}{}}
\LANGes{
\Question{
La pendiente de una escalera es \(30^{\circ}\).  ¿Cuál es la profundidad de un escalón si su altura es \( 15\,\mathrm{cm} \)?}
{\(15\sqrt3\,\mathrm{cm} \)}{\(30\,\mathrm{cm} \)}{\( \frac{15}{\sqrt3}\,\mathrm{cm} \)}{\( 28\,\mathrm{cm} \)}{}{}}
\End

\Data\ID{2010015310}
\Updated{2022-08-25 10:08:18}
\Level{C}
\SubArea{Triangles}
\LANGen{
\Question{
The shadow of a \(32\,\mathrm{m}\) high tree is \(40\) meters. At the same moment, the shadow of a person is \(200\,\mathrm{cm}\) long. What is the height of the person?}
{\(160\,\mathrm{cm}\)}{\(170\,\mathrm{cm}\)}{\(180\,\mathrm{cm}\)}{\(185\,\mathrm{cm}\)}{}{}}
\LANGpl{
\Question{
Cień drzewa o wysokości \(32\,\mathrm{m}\) wynosi \(40\) metrów. W tym samym momencie cień osoby ma długość \(200\,\mathrm{cm}\). Jaki wzrost ma ta osoba?}
{\(160\,\mathrm{cm}\)}{\(170\,\mathrm{cm}\)}{\(180\,\mathrm{cm}\)}{\(185\,\mathrm{cm}\)}{}{}}
\LANGcs{
\Question{
Stín stromu vysokého \(32\,\mathrm{m}\) měří \(40\) metrů. Ve stejný okamžik měří stín stojícího člověka \(200\,\mathrm{cm}\). Jak je tento člověk vysoký?}
{\(160\,\mathrm{cm}\)}{\(170\,\mathrm{cm}\)}{\(180\,\mathrm{cm}\)}{\(185\,\mathrm{cm}\)}{}{}}
\LANGsk{
\Question{
Strom s výškou  \(32\,\mathrm{m}\) vrhá  \(40\) metrový  tieň.  Vypočítajte výšku osoby, ktorá v rovnakom okamihu vrhá \(200\,\mathrm{cm}\) tieň.}
{\(160\,\mathrm{cm}\)}{\(170\,\mathrm{cm}\)}{\(180\,\mathrm{cm}\)}{\(185\,\mathrm{cm}\)}{}{}}
\LANGes{
\Question{
Un árbol de \(32\,\mathrm{m}\) proyecta una sombra de \(40\) metros. Al mismo tiempo, un hombre proyecta una sombra de \(200\,\mathrm{cm}\). Calcula la altura del hombre.}
{\(160\,\mathrm{cm}\)}{\(170\,\mathrm{cm}\)}{\(180\,\mathrm{cm}\)}{\(185\,\mathrm{cm}\)}{}{}}
\End

\Data\ID{9000036101}
\Updated{2021-04-24 04:04:40}
\Level{C}
\SubArea{Triangles}
\IMG{000361_01-figure0.pdf}
\LANGen{
\Question{
A \(3\, \mathrm{m}\) long
rod is in a slant position with respect to the observer's eye: one end is in the distance
\(20\, \mathrm{m}\) and the
other one \(18\, \mathrm{m}\).
Find the visual angle of the rod (the angle between the lines which connect the
observer's eye and the ends of the rod) and round to the nearest degrees.}
{\(7^{\circ }\)}{\(3^{\circ }\)}{\(45^{\circ }\)}{\(83^{\circ }\)}{}{}}
\LANGpl{
\Question{
Pręt o długości \(3\) jest nachylony względem oka obserwatora: jeden koniec jest w odległości \(20\, \mathrm{m}\), a drugi \(18\, \mathrm{m}\).  Znajdź kąt widzenia pręta (kąt między liniami, które łączą oko obserwatora i końce pręta) i zaokrąglij do pełnego stopnia.}
{\(7^{\circ }\)}{\(3^{\circ }\)}{\(45^{\circ }\)}{\(83^{\circ }\)}{}{}}
\LANGcs{
\Question{
V jakém zorném úhlu se jeví pozorovateli tyč dlouhá
\(3\, \mathrm{m}\), je-li od jednoho jejího
konce vzdálen \(20\, \mathrm{m}\) a
od druhého konce \(18\, \mathrm{m}\)?
Výsledek zaokrouhlete na celé stupně.}
{\(7^{\circ }\)}{\(3^{\circ }\)}{\(45^{\circ }\)}{\(83^{\circ }\)}{}{}}
\LANGsk{
\Question{
V akom zornom uhle sa javí pozorovateľovi tyč dlhá
\(3\, \mathrm{m}\), ak je od jedného jeho
konca vzdialený \(20\, \mathrm{m}\) a
od druhého konca \(18\, \mathrm{m}\)?
Výsledok zaokrúhlite na celé stupne.}
{\(7^{\circ }\)}{\(3^{\circ }\)}{\(45^{\circ }\)}{\(83^{\circ }\)}{}{}}
\LANGes{
\Question{
Una barra de \(3\, \mathrm{m}\) está inclinada respecto al ojo de un observador: un extremo está a una distancia de
\(20\, \mathrm{m}\) y el otro a una de \(18\, \mathrm{m}\).
Halla el ángulo visual de la barra (el ángulo entre las líneas que conectan el ojo del observador con los extremos de la barra) y redondea el resultado a grados.}
{\(7^{\circ }\)}{\(3^{\circ }\)}{\(45^{\circ }\)}{\(83^{\circ }\)}{}{}}
\End

\Data\ID{9000036102}
\Updated{2021-04-24 04:04:38}
\Level{C}
\SubArea{Triangles}
\LANGen{
\Question{
Three forces act on the same body in the same point and the total
force on the body is zero (the forces cancel). The first two forces are
\(8\, \mathrm{N}\) and
\(10\, \mathrm{N}\) and the angle between
these forces is \(55^{\circ }\).
Find the third force.}
{\(16\, \mathrm{N}\)}{\(15\, \mathrm{N}\)}{\(17\, \mathrm{N}\)}{\(18\, \mathrm{N}\)}{}{}}
\LANGpl{
\Question{
Trzy siły działają na to samo ciało w tym samym punkcie, a całkowita siła działająca na ciało jest równa zeru
(siły równoważą się). Dwie pierwsze siły wynoszą 
\(8\, \mathrm{N}\) i
\(10\, \mathrm{N}\), a kąt między nimi jest równy
 \(55^{\circ }\).
Oblicz wartość trzeciej siły.}
{\(16\, \mathrm{N}\)}{\(15\, \mathrm{N}\)}{\(17\, \mathrm{N}\)}{\(18\, \mathrm{N}\)}{}{}}
\LANGcs{
\Question{
V jednom bodě působí síly \(F_{1}\) 
a \(F_{2}\) o
velikostech \(8\, \mathrm{N}\) a
\(10\, \mathrm{N}\) a svírají spolu
úhel \(55^{\circ }\). Vypočítejte
velikost síly \(F_{3}\),
která působí ve stejném bodě a svými účinky ruší působení sil
\(F_{1}\) a
\(F_{2}\).}
{\(16\, \mathrm{N}\)}{\(15\, \mathrm{N}\)}{\(17\, \mathrm{N}\)}{\(18\, \mathrm{N}\)}{}{}}
\LANGsk{
\Question{
V jednom bode pôsobia sily \(F_{1}\) 
a \(F_{2}\) o
veľkostiach \(8\, \mathrm{N}\) a
\(10\, \mathrm{N}\) a zvierajú spolu
uhol \(55^{\circ }\). Vypočítajte
veľkosť sily \(F_{3}\),
ktorá pôsobí v rovnakom bode a svojimi účinkami ruší pôsobenie síl
\(F_{1}\) a
\(F_{2}\).}
{\(16\, \mathrm{N}\)}{\(15\, \mathrm{N}\)}{\(17\, \mathrm{N}\)}{\(18\, \mathrm{N}\)}{}{}}
\LANGes{
\Question{
Tres fuerzas actúan sobre el mismo cuerpo en el mismo punto y la fuerza total sobre el cuerpo es nula (las fuerzas se cancelan). Las dos primeras fuerzas son de
\(8\, \mathrm{N}\) y \(10\, \mathrm{N}\) y el ángulo entre ellas mide \(55^{\circ }\). Halla la tercera fuerza.}
{\(16\, \mathrm{N}\)}{\(15\, \mathrm{N}\)}{\(17\, \mathrm{N}\)}{\(18\, \mathrm{N}\)}{}{}}
\End

\Data\ID{9000036103}
\Updated{2021-04-24 04:04:29}
\Level{C}
\SubArea{Triangles}
\LANGen{
\Question{
Three forces \(F_{1}\),
\(F_{2}\) and
\(F_{3}\) act on the
same body in the same point and the total force on the body is zero (the forces cancel). The first
two forces are \(F_{1} = 8\, \mathrm{N}\) 
and \(F_{2} = 10\, \mathrm{N}\) and the
angle between \(F_{1}\) 
and \(F_{2}\) is
\(55^{\circ }\). Find the
angle between \(F_{3}\) 
and \(F_{1}\).
Round your answer to the nearest degrees.}
{\(149^{\circ }\)}{\(125^{\circ }\)}{\(55^{\circ }\)}{\(30^{\circ }\)}{}{}}
\LANGpl{
\Question{
Trzy siły \(F_{1}\), \(F_{2}\) i \(F_{3}\) działają na to samo ciało w tym samym punkcie, a całkowita siła działąjaca na ciało wynosi zero (siły równoważą się).  Dwie pierwsze siły  są:
 \(F_{1} = 8\, \mathrm{N}\) i \(F_{2} = 10\, \mathrm{N}\), a kąt pomiędzy  \(F_{1}\) 
i \(F_{2}\) wynosi \(55^{\circ }\). Oblicz kąt pomiędzy \(F_{3}\)  i \(F_{1}\). Zaokrąglij odpowiedź do pełnych stopni.}
{\(149^{\circ }\)}{\(125^{\circ }\)}{\(55^{\circ }\)}{\(30^{\circ }\)}{}{}}
\LANGcs{
\Question{
V jednom bodě působí síly \(F_{1}\) 
a \(F_{2}\) o
velikostech \(8\, \mathrm{N}\) a
\(10\, \mathrm{N}\) a svírají spolu
úhel \(55^{\circ }\). Ve stejném
bodě působí síla \(F_{3}\),
která svými účinky ruší působení sil
\(F_{1}\) a
\(F_{2}\). Určete úhel,
který spolu svírá \(F_{3}\) 
a \(F_{1}\).
Výsledek zaokrouhlete na celé stupně.}
{\(149^{\circ }\)}{\(125^{\circ }\)}{\(55^{\circ }\)}{\(30^{\circ }\)}{}{}}
\LANGsk{
\Question{
V jednom bode pôsobia sily \(F_{1}\) 
a \(F_{2}\) o
veľkostiach \(8\, \mathrm{N}\) a
\(10\, \mathrm{N}\) a zvierajú spolu
uhol \(55^{\circ }\). V rovnakom
bode pôsobí sila \(F_{3}\),
ktorá svojimi účinkami ruší pôsobenie síl
\(F_{1}\) a
\(F_{2}\). Určte uhol,
ktorý spolu zviera \(F_{3}\) 
a \(F_{1}\).
Výsledok zaokrúhlite na celé stupne.}
{\(149^{\circ }\)}{\(125^{\circ }\)}{\(55^{\circ }\)}{\(30^{\circ }\)}{}{}}
\LANGes{
\Question{
Tres fuerzas \(F_{1}\), \(F_{2}\) y \(F_{3}\) actúan sobre el mismo cuerpo en el mismo punto y la fuerza total sobre ele cuerpo es nula (las fuerzas se cancelan). Las dos primeras fuerzas son de \(F_{1} = 8\, \mathrm{N}\) y \(F_{2} = 10\, \mathrm{N}\) y el ángulo entre \(F_{1}\) y \(F_{2}\) mide \(55^{\circ }\). Halla el ángulo entre \(F_{3}\) y \(F_{1}\). Redondea el resultado a los grados más cercanos.}
{\(149^{\circ }\)}{\(125^{\circ }\)}{\(55^{\circ }\)}{\(30^{\circ }\)}{}{}}
\End

\Data\ID{9000036106}
\Updated{2021-04-24 05:04:08}
\Level{C}
\SubArea{Triangles}
\LANGen{
\Question{
Two straight roads go off from the crossing. The angle between directions of the roads is
\(52^{\circ }18'\).
A significant tree is on the first road in the distance
\(250\, \mathrm{m}\) from
the crossing. A rock with a beautiful view is on the second road in the distance
\(380\, \mathrm{m}\) from
the crossing. Find the direct distance (length of a line segment) from the rock to the
tree and round your answer to nearest meters.}
{\(301\, \mathrm{m}\)}{\(411\, \mathrm{m}\)}{\(568\, \mathrm{m}\)}{\(629\, \mathrm{m}\)}{}{}}
\LANGpl{
\Question{
Dwie proste drogi wychodzą ze skrzyżowania. Kąt między kierunkami dróg wynosi \(52^{\circ }18'\). Na pierwszej drodze w odległości \(250\, \mathrm{m}\) od skrzyżowania znajduje się drzewo. Na drugiej drodze w odległości \(380\, \mathrm{m}\) od skrzyżowania znajduje się skala z pięknym widokiem. Znajdź bezpośrednią odległość (długość linii) od skały do drzewa i zaokrąglij swoją odpowiedź do pełnych metrów.}
{\(301\, \mathrm{m}\)}{\(411\, \mathrm{m}\)}{\(568\, \mathrm{m}\)}{\(629\, \mathrm{m}\)}{}{}}
\LANGcs{
\Question{
Dvě přímé cesty vycházejí z rozcestníku
\(R\) a svírají úhel
\(52^{\circ }18'\). Na jedné z těchto
cest ve vzdálenosti \(250\, \mathrm{m}\) 
od rozcestníku \(R\) je
místo \(A\), na druhé
ve vzdálenosti \(380\, \mathrm{m}\) 
od rozcestníku \(R\) je
místo \(B\). Vypočítejte
vzdálenost míst \(A\) 
a \(B\) (tzn. délku
úsečky \(AB\)).
Výsledek zaokrouhlete na celé metry.}
{\(301\, \mathrm{m}\)}{\(411\, \mathrm{m}\)}{\(568\, \mathrm{m}\)}{\(629\, \mathrm{m}\)}{}{}}
\LANGsk{
\Question{
Dve priame cesty vychádzajú z rázcestia
\(R\) a zvierajú uhol
\(52^{\circ }18'\). Na jednej z týchto
ciest vo vzdialenosti \(250\, \mathrm{m}\) 
od rázcestia \(R\) je
miesto \(A\), na druhej
vo vzdialenosti \(380\, \mathrm{m}\) 
od rázcestia \(R\) je
miesto \(B\). Vypočítajte
vzdialenosť miest \(A\) 
a \(B\) (tzn. dĺžku
úsečky \(AB\)).
Výsledok zaokrúhlite na celé metre.}
{\(301\, \mathrm{m}\)}{\(411\, \mathrm{m}\)}{\(568\, \mathrm{m}\)}{\(629\, \mathrm{m}\)}{}{}}
\LANGes{
\Question{
Dos caminos rectos salen de un poste indicador \(R\) y forman un ángulo \(52^{\circ }18'\). En uno de los caminos a una distancia de \(250\, \mathrm{m}\) 
desde el poste indicador \(R\) hay un punto \(A\), en otro camino a una distancia de \(380\, \mathrm{m}\) 
desde el poste indicador \(R\) hay un punto \(B\). Calcula la distancia entre los puntos \(A\) y \(B\) (es decir, el segmento \(AB\)). Redondea el resultado a metros.}
{\(301\, \mathrm{m}\)}{\(411\, \mathrm{m}\)}{\(568\, \mathrm{m}\)}{\(629\, \mathrm{m}\)}{}{}}
\End

\Data\ID{9000036107}
\Updated{2021-04-24 05:04:05}
\Level{C}
\SubArea{Triangles}
\LANGen{
\Question{
There are three information panels \(A\),
\(B\) and
\(C\) in the park. The direct
distance between \(B\) 
and \(C\) is
\(150\, \mathrm{m}\). The visual angle of this
distance from the panel \(A\) 
is \(55^{\circ }\). The visual angle
of the distance \(AC\) 
from the panel \(B\) is
\(39^{\circ }\). Find the direct distance
between the panels \(A\) 
and \(B\) 
and round your answer to nearest meters.}
{\(183\, \mathrm{m}\)}{\(147\, \mathrm{m}\)}{\(195\, \mathrm{m}\)}{\(218\, \mathrm{m}\)}{}{}}
\LANGpl{
\Question{
W parku znajdują się trzy tablice informacyjne: \(A\),
\(B\) i \(C\).  Bezpośrednia odległość między \(B\)  i \(C\) wynosi \(150\, \mathrm{m}\). Kąt widzenia tej odległości od tablicy \(A\) 
wynosi \(55^{\circ }\). Kąt widzenia odległości  \(AC\)  od panelu \(B\) wynosi
\( 39^{\circ }\). Znajdź bezpośrednią odległość między panelami A i B i zaokrąglij swoją odpowiedź do pełnych metrów.}
{\(183\, \mathrm{m}\)}{\(147\, \mathrm{m}\)}{\(195\, \mathrm{m}\)}{\(218\, \mathrm{m}\)}{}{}}
\LANGcs{
\Question{
V parku jsou tři informační tabule
\(A\),
\(B\) a
\(C\). Přímá
vzdálenost tabulí \(B\) 
a \(C\) je
\(150\, \mathrm{m}\). Od tabule
\(A\) vidíme
tabule \(B\) a
\(C\) pod zorným úhlem
\(55^{\circ }\) a od tabule
\(B\) vidíme
tabule \(A\) a
\(C\) pod zorným úhlem
\( 39^{\circ }\). Jaká je přímá
vzdálenost tabulí \(A\) 
a \(B\)?
Výsledek zaokrouhlete na celé metry.}
{\(183\, \mathrm{m}\)}{\(147\, \mathrm{m}\)}{\(195\, \mathrm{m}\)}{\(218\, \mathrm{m}\)}{}{}}
\LANGsk{
\Question{
V parku sú tri informačné tabule
\(A\),
\(B\) a
\(C\). Priama
vzdialenosť tabúľ \(B\) 
a \(C\) je
\(150\, \mathrm{m}\). Od tabule
\(A\) vidíme
tabuľu \(B\) a
\(C\) pod zorným uhlom
\( 55^{\circ }\) a od tabule
\(B\) vidíme
tabuľu \(A\) a
\(C\) pod zorným uhlom
\( 39^{\circ }\). Aká je priama
vzdialenosť tabúľ \(A\) 
a \(B\)?
Výsledok zaokrúhlite na celé metre.}
{\(183\, \mathrm{m}\)}{\(147\, \mathrm{m}\)}{\(195\, \mathrm{m}\)}{\(218\, \mathrm{m}\)}{}{}}
\LANGes{
\Question{
En un parque hay tres paneles de información \(A\), \(B\) y \(C\). La distancia en linea recta entre \(B\) y \(C\) es \(150\, \mathrm{m}\). El ángulo visual de esta distancia desde el panel \(A\) es \(55^{\circ }\). El ángulo visual de la distancia \(AC\) desde el panel \(B\) es \(39^{\circ }\). Clacula la distancia en linea recta entre los paneles \(A\) y \(B\) y redondea el resultado a los metros más cercanos.}
{\(183\, \mathrm{m}\)}{\(147\, \mathrm{m}\)}{\(195\, \mathrm{m}\)}{\(218\, \mathrm{m}\)}{}{}}
\End

\Data\ID{9000036108}
\Updated{2021-04-24 05:04:13}
\Level{C}
\SubArea{Triangles}
\IMG{000361_08-figure0.pdf}
\LANGen{
\Question{
The center of a spherical balloon is at a height of \(500\, \mathrm{m}\) 
height. The visual angle of the balloon is
\(1^{\circ }30'\). The elevation angle of the
center of the balloon is \(42^{\circ }50'\).
Find the diameter of the balloon in meters and round to nearest one decimal.}
{\(19.3\, \mathrm{m}\)}{\(18.2\, \mathrm{m}\)}{\(18.9\, \mathrm{m}\)}{\(19.5\, \mathrm{m}\)}{}{}}
\LANGpl{
\Question{
Środek kulistego balonu znajduje się na wysokośći \(500\, \mathrm{m}\). Kąt widzenia balonu wynosi \(1^{\circ }30'\). Elewacja środka balonu wynosi \(42^{\circ }50'\). Znajdź średnicę balonu w metrach i zaokrąglij do miejsc dziesiętnych po przecinku.}
{\(19.3\, \mathrm{m}\)}{\(18.2\, \mathrm{m}\)}{\(18.9\, \mathrm{m}\)}{\(19.5\, \mathrm{m}\)}{}{}}
\LANGcs{
\Question{
Horkovzdušný balón tvaru koule má střed ve výšce
\(500\, \mathrm{m}\) nad
zemí. Pozorujeme ho z místa na zemi, z něhož ho vidíme v zorném úhlu
\(1^{\circ }30'\).
Z místa pozorování má střed balónu výškový úhel
\(42^{\circ }50'\).
Vypočítejte průměr balónu v metrech. Výsledek zaokrouhlete na jedno
desetinné místo.}
{\(19{,}3\, \mathrm{m}\)}{\(18{,}2\, \mathrm{m}\)}{\(18{,}9\, \mathrm{m}\)}{\(19{,}5\, \mathrm{m}\)}{}{}}
\LANGsk{
\Question{
Teplovzdušný balón tvaru gule má stred vo výške
\(500\, \mathrm{m}\) nad
zemou. Pozorujeme ho z miesta na zemi, z neho ho vidíme v zornom uhle
\(1^{\circ }30'\).
Z miesta pozorovania má stred balónu výškový uhol
\(42^{\circ }50'\).
Vypočítajte priemer balóna v metroch. Výsledok zaokrúhlite na jedno
desatinné miesto.}
{\(19{,}3\, \mathrm{m}\)}{\(18{,}2\, \mathrm{m}\)}{\(18{,}9\, \mathrm{m}\)}{\(19{,}5\, \mathrm{m}\)}{}{}}
\LANGes{
\Question{
El centro de un globo esférico está a una altura de \(500\, \mathrm{m}\). El ángulo visual del globo es de \(1^{\circ }30'\). El ángulo de elevación del centro del globo es de \(42^{\circ }50'\). Calcula el diámetro del globo en metros y redondea el resultado al decimal más cercano.}
{\(19.3\, \mathrm{m}\)}{\(18.2\, \mathrm{m}\)}{\(18.9\, \mathrm{m}\)}{\(19.5\, \mathrm{m}\)}{}{}}
\End

\Data\ID{9000036109}
\Updated{2021-04-24 05:04:12}
\Level{C}
\SubArea{Triangles}
\LANGen{
\Question{
The point \(A\) is
located \(20\, \mathrm{cm}\) from a
mirror and the point \(B\) 
is located \(50\, \mathrm{cm}\) 
from the same mirror. The direct distance between
\(A\) and
\(B\) (the length of
the segment \(AB\))
is \(70\, \mathrm{cm}\).
Find the angle of incidence of the ray through the point
\(A\) which is reflected
to the point \(B\) 
and round your answer to nearest degrees. (The angle of incidence is the angle
between the incident ray and the normal to the mirror.)}
{\(42^{\circ }\)}{\(37^{\circ }\)}{\(38^{\circ }\)}{\(48^{\circ }\)}{}{}}
\LANGpl{
\Question{
Punkt   \(A\) znajduje się \(20\, \mathrm{cm}\)  od lustra, a punkt \(B\)  znajduje się \(50\, \mathrm{cm}\)  od tego samego lustra. Bezpośrednia odległość między \(A\) i
\(B\)  (długość odcinka \(AB\))
wynosi  \(70\, \mathrm{cm}\). Znajdź kąt padania promienia przez punkt \(A\), który jest odbity do punktu \(B\)  i zaokrąglij swoją odpowiedź do pełnych stopni. (Kąt padania to kąt pomiędzy promieniem padającym a normalną.)}
{\(42^{\circ }\)}{\(37^{\circ }\)}{\(38^{\circ }\)}{\(48^{\circ }\)}{}{}}
\LANGcs{
\Question{
Jaký je úhel dopadu paprsku, který projde bodem
\(A\) a po odrazu od zrcadla
projde bodem \(B\)? Bod
\(A\) je ve vzdálenosti
\(20\, \mathrm{cm}\) od zrcadla a bod
\(B\) ve vzdálenosti
\(50\, \mathrm{cm}\) od zrcadla.
Vzdálenost \(|AB| = 70\, \mathrm{cm}\).
(Pozn.: úhel dopadu paprsku je úhel mezi kolmicí dopadu a dopadajícím
paprskem.) Výsledek zaokrouhlete na celé stupně.}
{\(42^{\circ }\)}{\(37^{\circ }\)}{\(38^{\circ }\)}{\(48^{\circ }\)}{}{}}
\LANGsk{
\Question{
Aký je uhol dopadu lúča, ktorý prejde bodom
\(A\) a po odraze od zrkadla
prejde bodom \(B\)? Bod
\(A\) je vo vzdialenosti
\(20\, \mathrm{cm}\) od zrkadla a bod
\(B\) vo vzdialenosti
\(50\, \mathrm{cm}\) od zrkadla.
Vzdialenosť \(|AB| = 70\, \mathrm{cm}\).
(Pozn.: uhol dopadu lúča je uhol medzi kolmicou dopadu a dopadajúcim
lúčom.) Výsledok zaokrúhlite na celé stupne.}
{\(42^{\circ }\)}{\(37^{\circ }\)}{\(38^{\circ }\)}{\(48^{\circ }\)}{}{}}
\LANGes{
\Question{
El punto \(A\) se sitúa a \(20\, \mathrm{cm}\) desde el espejo y el punto \(B\) se sitúa a \(50\, \mathrm{cm}\) desde le mismo espejo. La distancia directa entre \(A\) y
\(B\) (la distancia del segmento \(AB\)) es \(70\, \mathrm{cm}\). Calcula el ángulo de incidencia del rayo a través del punto \(A\) que se refleja en el punto \(B\) y redondea el resultado a los grados más cercanos. (El ángulo de incidencia es el ángulo entre el rayo incidente y el normal al espejo).}
{\(42^{\circ }\)}{\(37^{\circ }\)}{\(38^{\circ }\)}{\(48^{\circ }\)}{}{}}
\End

\Data\ID{9000036110}
\Updated{2021-04-24 05:04:29}
\Level{C}
\SubArea{Triangles}
\LANGen{
\Question{
A tower is observed from two different places \(A\) and \(B\). The direct distance between \(A\) and \(B\) is \(65\, \mathrm{m}\). If we denote the bottom of the tower by \(C\), we get a triangle \(ABC\) in which the measure of \(\measuredangle CAB \) is \(71^{\circ }\) and the measure of \(\measuredangle ABC \) is \( 34^{\circ }\). From the point \(A\) the angle of elevation of the top of the tower is \(40^{\circ }18'\). Find the height of the tower. Suppose that all \(A\), \(B\) and \(C\) are in the same height above sea level and round your answers to nearest meters.}
{\(32\, \mathrm{m}\)}{\(30\, \mathrm{m}\)}{\(35\, \mathrm{m}\)}{\(38\, \mathrm{m}\)}{}{}}
\LANGpl{
\Question{
Wieża jest obserwowana z dwóch różnych miejsc \(A\) i
\(B\). Bezpośrednia odległość między \(A\) i
\(B\) wynosi \(65\, \mathrm{m}\). Jeśli oznaczymy dno wieży przez \(C\), otryzmamy trójkąt \(ABC\), w ktorým miara kąta \( CAB \) wynosi \( 71^{\circ }\) i 
miara kąta \(ABC\) wynosi \( 34^{\circ }\). Z punktu \(A\) kąt wzniesienia szczytu wieży ma miarę \(40^{\circ }18'\).  Wyznacz wysokość wieży. Załóżmy, że wszystkie miejsca \(A\),
\(B\) i \(C\) znajdują się na tej samej wysokości nad poziomem morza. Zaokrąglij swoje odpowiedzi do pełnych metrów.}
{\(32\, \mathrm{m}\)}{\(30\, \mathrm{m}\)}{\(35\, \mathrm{m}\)}{\(38\, \mathrm{m}\)}{}{}}
\LANGcs{
\Question{
Určete výšku rozhledny, kterou pozorujeme ze dvou míst
\(A\) a
\(B\). Pata
rozhledny \(P\) 
a body \(A\) a
\(B\) tvoří vrcholy
trojúhelníku \(ABP\),
\(|AB| = 65\, \mathrm{m}\),
\(|\measuredangle PAB| = 71^{\circ }\),
\(|\measuredangle ABP| = 34^{\circ }\). Vrchol rozhledny je
vidět z místa \(A\) pod
výškovým úhlem \(40^{\circ }18'\).
Body \(A\),
\(B\) a
\(P\) 
leží ve stejné nadmořské výšce. Výsledek zaokrouhlete na celé metry.}
{\(32\, \mathrm{m}\)}{\(30\, \mathrm{m}\)}{\(35\, \mathrm{m}\)}{\(38\, \mathrm{m}\)}{}{}}
\LANGsk{
\Question{
Určte výšku rozhľadne, ktorú pozorujeme z dvoch miest
\(A\) a
\(B\). Päta
rozhľadne \(P\) 
a body \(A\) a
\(B\) tvoria vrcholy
trojuholníka \(ABP\),
\(|AB| = 65\, \mathrm{m}\),
\(|\measuredangle PAB| = 71^{\circ }\),
\(|\measuredangle ABP| = 34^{\circ }\). Vrchol rozhľadne je
vidieť z miesta \(A\) pod
výškovým uhlom \(40^{\circ }18'\).
Body \(A\),
\(B\) a
\(P\) 
ležia v rovnakej nadmorskej výške. Výsledok zaokrúhlite na celé metre.}
{\(32\, \mathrm{m}\)}{\(30\, \mathrm{m}\)}{\(35\, \mathrm{m}\)}{\(38\, \mathrm{m}\)}{}{}}
\LANGes{
\Question{
Se observa una torre desde los lugares \(A\) y \(B\). La distancia directa entre \(A\) y \(B\) es de \(65\, \mathrm{m}\). Si denotamos la parte inferior de la torre por \(C\), obtenemos un triángulo \(ABC\) en el que la medida de \(\measuredangle CAB \) es de\(71^{\circ }\) y al medida de \(\measuredangle ABC \) es de\( 34^{\circ }\). Desde el punto \(A\) el ángulo de elevación de la parte superior de la torre es de \(40^{\circ }18'\). Halla la altura de la torre. Suponemos que \(A\), \(B\) y \(C\) tienen la misma altitud. Redondea el resultado a los metros más cercanos.}
{\(32\, \mathrm{m}\)}{\(30\, \mathrm{m}\)}{\(35\, \mathrm{m}\)}{\(38\, \mathrm{m}\)}{}{}}
\End

\Data\ID{9000038701}
\Updated{2022-01-22 07:01:59}
\Level{C}
\SubArea{Triangles}
\IMG{000387_01-figure0.pdf}
\LANGen{
\Question{
The box is on the slope as in the picture. The angle of the slope is
\(\alpha \).
The forces acting on the box are the force of gravity
\(\vec{F_{G}}\) and the
friction \(\vec{F_{t}}\).
The force of gravity can be replaced by two components
\(\vec{F_{1}}\) and
\(\vec{F_{n}}\). (The force
\(\vec{F_{1}}\) is parallel to the slope
and \(\vec{F_{n}}\) is perpendicular to
the slope.) The friction \(F_{t}\) 
is given by the formula \(F_{t} = fF_{n}\),
where \(f\) 
is the coefficient of the friction.What is the influence of the increasing angle
\(\alpha \) on
the forces acting on the box?}
{\(F_{1}\) becomes
bigger and \(F_{t}\) 
becomes smaller}{both \(F_{1}\) 
and \(F_{t}\) 
become smaller}{\(F_{1}\) becomes
bigger, \(F_{t}\) 
does not change}{\(F_{1}\) becomes
smaller, \(F_{t}\) 
does not change}{both \(F_{1}\) 
and \(F_{t}\) 
become bigger}{\(F_{1}\) becomes
smaller and \(F_{t}\) 
becomes bigger}}
\LANGpl{
\Question{
Pudełko znajduje się na równi pochyłej o kącie nachylenia \( \alpha \) (jak na zdjęciu). Siły działające na pudełko to siła grawitacji \(\vec{F_{G}}\) i siła tarcia 
\(\vec{F_{t}}\). Siłę grawitacji można zastąpić dwoma składowymi  \(\vec{F_{1}}\) i \(\vec{F_{n}}\). (Siła \(\vec{F_{1}}\) jest równoległa do powierzchni równi, a siła \(\vec{F_{n}}\) jest prostopadła do powierzchni równi.) Tarcie \(F_{t}\)  jest określone przez wzór \(F_{t} = fF_{n}\), gdzie \(f\)  jest współczynnikiem tarcia. Jaki jest wpływ narastającego kąta  \(\alpha \)  na siły działające na pudło?}
{\(F_{1}\) zwiększa się, a \(F_{t}\) 
maleje.}{Zarówno \(F_{1}\) 
jak i \(F_{t}\) 
zmniejszają się.}{\(F_{1}\) zwiększa się, a \(F_{t}\) 
nie zmienia się.}{\(F_{1}\) zmniejsza się, a  \(F_{t}\) 
nie zmienia się.}{Zarówno \(F_{1}\) 
jak i \(F_{t}\) 
zwiększają się.}{\(F_{1}\) zmniejsza się, a  \(F_{t}\) 
zwiększa się.}}
\LANGcs{
\Question{
Kvádr položíme na nakloněnou rovinu se sklonem
\(\alpha \).
V tíhovém poli Země na něj bude působit tíhová síla
\(\vec{F_{G}}\) a síla
tření \(\vec{F_{t}}\).
Tíhovou sílu můžeme nahradit jejími složkami
\(\vec{F_{1}}\) a
\(\vec{F_{n}}\), kde
\(\vec{F_{1}}\) 
má směr rovnoběžný s nakloněnou rovinou a
\(\vec{F_{n}}\) 
je na ní kolmá. Pro velikost třecí síly platí
\(F_{t} = fF_{n}\), kde
\(f\) je
součinitel smykového tření. Zvětšíme-li úhel \(\alpha \),
pak:}
{se zvětší \(F_{1}\) 
a \(F_{t}\) se
zmenší.}{se zmenší \(F_{1}\) 
i \(F_{t}\).}{se zvětší \(F_{1}\) 
a \(F_{t}\) se
nezmění.}{se zmenší \(F_{1}\) 
a \(F_{t}\) se
nezmění.}{se zvětší \(F_{1}\) 
i \(F_{t}\).}{se zmenší \(F_{1}\) 
a \(F_{t}\) se
zvětší.}}
\LANGsk{
\Question{
Kváder položíme na naklonenú rovinu so sklonom
\(\alpha \).
V gravitačnom poli Zeme na neho bude pôsobiť gravitačná sila
\(\vec{F_{G}}\) a sila
trenia \(\vec{F_{t}}\).
Gravitačnú silu môžeme nahradiť jej zložkami
\(\vec{F_{1}}\) a
\(\vec{F_{n}}\), kde
\(\vec{F_{1}}\) 
má smer rovnobežný s naklonenou rovinou a
\(\vec{F_{n}}\) 
je na ňu kolmá. Pre veľkosť trecej sily platí
\(F_{t} = fF_{n}\), kde
\(f\) je
súčiniteľ šmykového trenia. Ak zväčšíme uhol \(\alpha \),
potom:}
{sa zväčší \(F_{1}\) 
a \(F_{t}\) sa
zmenší.}{sa zmenší \(F_{1}\) 
i \(F_{t}\).}{sa zväčší \(F_{1}\) 
a \(F_{t}\) sa
nezmení.}{sa zmenší \(F_{1}\) 
a \(F_{t}\) sa
nezmení.}{sa zväčší \(F_{1}\) 
i \(F_{t}\).}{sa zmenší \(F_{1}\) 
a \(F_{t}\) sa
zväčší.}}
\LANGes{
\Question{
Un cuerpo está en un plano inclinado (observa la imagen). El ángulo de la pendiente es \(\alpha \). Las fuerzas que actúan sobre el cuerpo son la fuerza de gravedad \(\vec{F_{G}}\) y la fricción \(\vec{F_{t}}\). La fuerza de la gravedad se puede descomponer en dos componentes \(\vec{F_{1}}\) y \(\vec{F_{n}}\). (La fuerza \(\vec{F_{1}}\) es paralela a la pendiente y \(\vec{F_{n}}\) es perpendicular a la pendiente.) La fricción \(F_{t}\) viene dada por la fórmula \(F_{t} = fF_{n}\), donde \(f\) es el coeficiente de fricción. ¿Cómo influye el crecimiento del ángulo \(\alpha \) en las fuerzas que actúan sobre el cuerpo?}
{\(F_{1}\) aumenta y \(F_{t}\) 
disminuye}{tanto \(F_{1}\) 
como \(F_{t}\) 
disminuyen}{\(F_{1}\) aumenta , \(F_{t}\) 
no cambia}{\(F_{1}\) disminuye, \(F_{t}\) 
no cambia}{tanto \(F_{1}\) 
como \(F_{t}\) 
aumentan}{\(F_{1}\)se hace menor y \(F_{t}\) 
aumenta}}
\End

\Data\ID{9000038702}
\Updated{2022-08-03 10:08:57}
\Level{C}
\SubArea{Triangles}
\IMG{000387_02-figure0.pdf}
\LANGen{
\Question{
The box is on the slope as in the picture. The angle of the slope is
\(\alpha \).
The forces acting on the box are the force of gravity
\(\vec{F_{G}}\) and the
friction \(\vec{F_{t}}\).
The force of gravity can be replaced by two components
\(\vec{F_{1}}\) and
\(\vec{F_{n}}\). (The force
\(\vec{F_{1}}\) is parallel to
the slope and \(\vec{F_{n}}\) 
is perpendicular to the slope.) Find \(F_{1}\).}
{\(F_{1} = F_{G}\sin \alpha \)}{\(F_{1} = \frac{F_{G}} 
 {\sin \alpha } \)}{\(F_{1} = F_{G}\mathop{\mathrm{tg}}\nolimits \alpha \)}{\(F_{1} = \frac{F_{G}} 
 {\mathop{\mathrm{tg}}\nolimits \alpha } \)}{\(F_{1} = F_{G}\cos \alpha \)}{\(F_{1} = \frac{F_{G}} 
 {\cos \alpha } \)}}
\LANGpl{
\Question{
Pudełko znajduje się na równi pochyłej o kącie nachylenia \( \alpha \) (jak na zdjęciu). Siły działające na pudełko to siła grawitacji \(\vec{F_{G}}\) i tarcia \(\vec{F_{t}}\). Siłę grawitacji można zastąpić dwoma składowymi  \(\vec{F_{1}}\) i 
\(\vec{F_{n}}\). (Siła \(\vec{F_{1}}\)  jest równoległa do powierzchni równi, a \(\vec{F_{n}}\)  jest prostopadła do powierzchni równi.) Znajdź \(F_{1}\).}
{\(F_{1} = F_{G}\sin \alpha \)}{\(F_{1} = \frac{F_{G}} 
 {\sin \alpha } \)}{\(F_{1} = F_{G}\mathop{\mathrm{tg}}\nolimits \alpha \)}{\(F_{1} = \frac{F_{G}} 
 {\mathop{\mathrm{tg}}\nolimits \alpha } \)}{\(F_{1} = F_{G}\cos \alpha \)}{\(F_{1} = \frac{F_{G}} 
 {\cos \alpha } \)}}
\LANGcs{
\Question{
Kvádr položíme na nakloněnou rovinu se sklonem
\(\alpha \).
V tíhovém poli Země na něj bude působit tíhová síla
\(\vec{F_{G}}\).
Tuto sílu můžeme nahradit jejími složkami
\(\vec{F_{1}}\) a
\(\vec{F_{n}}\), kde
\(\vec{F_{1}}\) 
má směr rovnoběžný s nakloněnou rovinou a
\(\vec{F_{n}}\) je na
ní kolmá. Pro \(F_{1}\) 
platí:}
{\(F_{1} = F_{G}\sin \alpha \)}{\(F_{1} = \frac{F_{G}} 
 {\sin \alpha } \)}{\(F_{1} = F_{G}\mathop{\mathrm{tg}}\nolimits \alpha \)}{\(F_{1} = \frac{F_{G}} 
 {\mathop{\mathrm{tg}}\nolimits \alpha } \)}{\(F_{1} = F_{G}\cos \alpha \)}{\(F_{1} = \frac{F_{G}} 
 {\cos \alpha } \)}}
\LANGsk{
\Question{
Kváder položíme na naklonenú rovinu so sklonom
\(\alpha \).
V gravitačnom poli Zeme na neho bude pôsobiť gravitačná sila
\(\vec{F_{G}}\).
Túto silu môžeme nahradiť jej zložkami
\(\vec{F_{1}}\) a
\(\vec{F_{n}}\), kde
\(\vec{F_{1}}\) 
má smer rovnobežný s naklonenou rovinou a
\(\vec{F_{n}}\) je na
ňu kolmá. Pre \(F_{1}\) 
platí:}
{\(F_{1} = F_{G}\sin \alpha \)}{\(F_{1} = \frac{F_{G}} 
 {\sin \alpha } \)}{\(F_{1} = F_{G}\mathop{\mathrm{tg}}\nolimits \alpha \)}{\(F_{1} = \frac{F_{G}} 
 {\mathop{\mathrm{tg}}\nolimits \alpha } \)}{\(F_{1} = F_{G}\cos \alpha \)}{\(F_{1} = \frac{F_{G}} 
 {\cos \alpha } \)}}
\LANGes{
\Question{
Un ortoedro se encuentra en un plano inclinado (observa la imagen). El ángulo de la pendiente es \(\alpha \). Las fuerzas que actúan sobre el ortoedro son la fuerza dela  gravedad \(\vec{F_{G}}\) y la de la fricción \(\vec{F_{t}}\). La fuerza de la gravedad se puede reemplazar por dos componentes \(\vec{F_{1}}\) y \(\vec{F_{n}}\). (La fuerza \(\vec{F_{1}}\) es paralela a la pendiente y \(\vec{F_{n}}\) es perpendicular a la pendiente).  Halla \(F_{1}\).}
{\(F_{1} = F_{G}\sin \alpha \)}{\(F_{1} = \frac{F_{G}} 
 {\sin \alpha } \)}{\(F_{1} = F_{G}\mathop{\mathrm{tg}}\nolimits \alpha \)}{\(F_{1} = \frac{F_{G}} 
 {\mathop{\mathrm{tg}}\nolimits \alpha } \)}{\(F_{1} = F_{G}\cos \alpha \)}{\(F_{1} = \frac{F_{G}} 
 {\cos \alpha } \)}}
\End

\Data\ID{9000038703}
\Updated{2022-08-03 10:08:17}
\Level{C}
\SubArea{Triangles}
\IMG{000387_03-figure0.pdf}
\LANGen{
\Question{
The box is on the slope as in the picture. The angle of the slope is
\(\alpha \).
The forces acting on the box are the force of gravity
\(\vec{F_{G}}\) and the
friction \(\vec{F_{t}}\).
The force of gravity can be replaced by two components
\(\vec{F_{1}}\) and
\(\vec{F_{n}}\). (The force
\(\vec{F_{1}}\) is parallel to
the slope and \(\vec{F_{n}}\) 
is perpendicular to the slope.) Find \(F_{p}\).}
{\(F_{p} = F_{G}\cos \alpha \)}{\(F_{p} = \frac{F_{G}} 
 {\cos \alpha } \)}{\(F_{p} = F_{G}\mathop{\mathrm{tg}}\nolimits \alpha \)}{\(F_{p} = \frac{F_{G}} 
 {\mathop{\mathrm{tg}}\nolimits \alpha } \)}{\(F_{p} = F_{G}\sin \alpha \)}{\(F_{p} = \frac{F_{G}} 
 {\sin \alpha } \)}}
\LANGpl{
\Question{
Pudełko znajduje się na równi pochyłej o kącie nachylenia \( \alpha \) (jak na zdjęciu). Siły działające na pudełko to siła grawitacji \(\vec{F_{G}}\) i tarcia \(\vec{F_{t}}\). Siłę grawitacji można zastąpić dwoma składowymi  \(\vec{F_{1}}\) i 
\(\vec{F_{n}}\). (Siła \(\vec{F_{1}}\)  jest równoległa do powierzchni równi, a \(\vec{F_{n}}\)  jest prostopadła do powierzchni równi.) Znajdź siłe reakcji podłoża \(F_{p}\), ktorá równoważy \( F_n \).}
{\(F_{p} = F_{G}\cos \alpha \)}{\(F_{p} = \frac{F_{G}} 
 {\cos \alpha } \)}{\(F_{p} = F_{G}\mathop{\mathrm{tg}}\nolimits \alpha \)}{\(F_{p} = \frac{F_{G}} 
 {\mathop{\mathrm{tg}}\nolimits \alpha } \)}{\(F_{p} = F_{G}\sin \alpha \)}{\(F_{p} = \frac{F_{G}} 
 {\sin \alpha } \)}}
\LANGcs{
\Question{
Kvádr položíme na nakloněnou rovinu se sklonem
\(\alpha \).
V tíhovém poli Země na něj bude působit tíhová síla
\(\vec{F_{G}}\), síla od
podložky \(\vec{F_{p}}\) a
síla tření \(\vec{F_{t}}\).
Tíhovou sílu můžeme nahradit jejími složkami
\(\vec{F_{1}}\) a
\(\vec{F_{n}}\), kde
\(\vec{F_{1}}\) 
má směr rovnoběžný s nakloněnou rovinou a
\(\vec{F_{n}}\) je na
ní kolmá. Pro \(F_{p}\) 
platí:}
{\(F_{p} = F_{G}\cos \alpha \)}{\(F_{p} = \frac{F_{G}} 
 {\cos \alpha } \)}{\(F_{p} = F_{G}\mathop{\mathrm{tg}}\nolimits \alpha \)}{\(F_{p} = \frac{F_{G}} 
 {\mathop{\mathrm{tg}}\nolimits \alpha } \)}{\(F_{p} = F_{G}\sin \alpha \)}{\(F_{p} = \frac{F_{G}} 
 {\sin \alpha } \)}}
\LANGsk{
\Question{
Kváder položíme na naklonenú rovinu so sklonom
\(\alpha \).
V gravitačnom poli Zeme na neho bude pôsobiť gravitačná sila
\(\vec{F_{G}}\), sila od
podložky \(\vec{F_{p}}\) a
sila trenia \(\vec{F_{t}}\).
Gravitačnú silu môžeme nahradiť jej zložkami
\(\vec{F_{1}}\) a
\(\vec{F_{n}}\), kde
\(\vec{F_{1}}\) 
má smer rovnobežný s naklonenou rovinou a
\(\vec{F_{n}}\) je na
ňu kolmá. Pre \(F_{p}\) 
platí:}
{\(F_{p} = F_{G}\cos \alpha \)}{\(F_{p} = \frac{F_{G}} 
 {\cos \alpha } \)}{\(F_{p} = F_{G}\mathop{\mathrm{tg}}\nolimits \alpha \)}{\(F_{p} = \frac{F_{G}} 
 {\mathop{\mathrm{tg}}\nolimits \alpha } \)}{\(F_{p} = F_{G}\sin \alpha \)}{\(F_{p} = \frac{F_{G}} 
 {\sin \alpha } \)}}
\LANGes{
\Question{
Un ortoedro se encuentra en un plano inclinado (observa la imagen).  El ángulo de la pendiente es
\(\alpha \).
Las fuerzas que actúan sobre el ortoedro son la fuerza de la gravedad
\(\vec{F_{G}}\) y la de la fricción \(\vec{F_{t}}\).
La fuerza de la gravedad se puede reemplazar por dos componentes
\(\vec{F_{1}}\) y
\(\vec{F_{n}}\). (La fuerza
\(\vec{F_{1}}\) es paralela a la pendiente y \(\vec{F_{n}}\) 
es perpendicular a la pendiente). Halla \(F_{p}\).}
{\(F_{p} = F_{G}\cos \alpha \)}{\(F_{p} = \frac{F_{G}} 
 {\cos \alpha } \)}{\(F_{p} = F_{G}\mathop{\mathrm{tg}}\nolimits \alpha \)}{\(F_{p} = \frac{F_{G}} 
 {\mathop{\mathrm{tg}}\nolimits \alpha } \)}{\(F_{p} = F_{G}\sin \alpha \)}{\(F_{p} = \frac{F_{G}} 
 {\sin \alpha } \)}}
\End

\Data\ID{9000038704}
\Updated{2022-08-03 10:08:04}
\Level{C}
\SubArea{Triangles}
\IMG{000387_04-figure0.pdf}
\LANGen{
\Question{
The box is on the slope as in the picture. The angle of the slope is
\(\alpha \).
The forces acting on the box are the force of gravity
\(\vec{F_{G}}\) and the
friction \(\vec{F_{t}}\).
The force of gravity can be replaced by two components
\(\vec{F_{1}}\) and
\(\vec{F_{n}}\). (The force
\(\vec{F_{1}}\) is parallel to
the slope and \(\vec{F_{n}}\) 
is perpendicular to the slope.) For \(F_{1} = 20\, \mathrm{N}\) and
\(F_{n} = 55\, \mathrm{N}\) find the
corresponding \(\alpha \).}
{\(\alpha \doteq 20^{\circ }\)}{\(\alpha \doteq 21^{\circ }\)}{\(\alpha \doteq 69^{\circ }\)}{\(\alpha \doteq 70^{\circ }\)}{\(\alpha \doteq 30^{\circ }\)}{\(\alpha \doteq 29^{\circ }\)}}
\LANGpl{
\Question{
Pudełko znajduje się na równi pochyłej o kącie nachylenia \( \alpha \) (jak na zdjęciu). Siły działające na pudełko to siła grawitacji \(\vec{F_{G}}\) i tarcia \(\vec{F_{t}}\). Siłę grawitacji można zastąpić dwoma składowymi  \(\vec{F_{1}}\) i 
\(\vec{F_{n}}\). (Siła \(\vec{F_{1}}\)  jest równoległa do powierzchni równi, a \(\vec{F_{n}}\)  jest prostopadła do powierzchni równi.) Dla \(F_{1} = 20\, \mathrm{N}\) i
\(F_{n} = 55\, \mathrm{N}\)  znajdź odpowiednie \(\alpha \).}
{\(\alpha \doteq 20^{\circ }\)}{\(\alpha \doteq 21^{\circ }\)}{\(\alpha \doteq 69^{\circ }\)}{\(\alpha \doteq 70^{\circ }\)}{\(\alpha \doteq 30^{\circ }\)}{\(\alpha \doteq 29^{\circ }\)}}
\LANGcs{
\Question{
Kvádr položíme na nakloněnou rovinu se sklonem
\(\alpha \).
V tíhovém poli Země na něj bude působit tíhová síla
\(\vec{F_{G}}\).
Tuto sílu můžeme nahradit jejími složkami
\(\vec{F_{1}}\) a
\(\vec{F_{n}}\), kde
\(\vec{F_{1}}\) 
má směr rovnoběžný s nakloněnou rovinou a
\(\vec{F_{n}}\) je na
ní kolmá. Je-li \(F_{1} = 20\, \mathrm{N}\) a
\(F_{n} = 55\, \mathrm{N}\), pak pro
úhel \(\alpha \) 
platí:}
{\(\alpha \doteq 20^{\circ }\)}{\(\alpha \doteq 21^{\circ }\)}{\(\alpha \doteq 69^{\circ }\)}{\(\alpha \doteq 70^{\circ }\)}{\(\alpha \doteq 30^{\circ }\)}{\(\alpha \doteq 29^{\circ }\)}}
\LANGsk{
\Question{
Kváder položíme na naklonenú rovinu so sklonom
\(\alpha \).
V gravitačnom poli Zeme na neho bude pôsobiť gravitačná sila
\(\vec{F_{G}}\).
Túto silu môžeme nahradiť jej zložkami
\(\vec{F_{1}}\) a
\(\vec{F_{n}}\), kde
\(\vec{F_{1}}\) 
má smer rovnobežný s naklonenou rovinou a
\(\vec{F_{n}}\) je na
ňu kolmá. Ak je \(F_{1} = 20\, \mathrm{N}\) a
\(F_{n} = 55\, \mathrm{N}\), potom pre
uhol \(\alpha \) 
platí:}
{\(\alpha \doteq 20^{\circ }\)}{\(\alpha \doteq 21^{\circ }\)}{\(\alpha \doteq 69^{\circ }\)}{\(\alpha \doteq 70^{\circ }\)}{\(\alpha \doteq 30^{\circ }\)}{\(\alpha \doteq 29^{\circ }\)}}
\LANGes{
\Question{
Un ortoedro está en un plano inclinado (observa la imagen). El ángulo de la pendiente es
\(\alpha \).
Las fuerzas que actúan sobre el ortoedro son la fuerza de la gravedad
\(\vec{F_{G}}\) y la de la fricción \(\vec{F_{t}}\).
La fuerza de gravedad se puede reemplazar por dos componentes
\(\vec{F_{1}}\) y
\(\vec{F_{n}}\). (La fuerza
\(\vec{F_{1}}\) es paralela a la pendiente y \(\vec{F_{n}}\) 
es perpendicular a la pendiente).  Suponiendo que \(F_{1} = 20\, \mathrm{N}\) y
\(F_{n} = 55\, \mathrm{N}\) halla \(\alpha \).}
{\(\alpha \doteq 20^{\circ }\)}{\(\alpha \doteq 21^{\circ }\)}{\(\alpha \doteq 69^{\circ }\)}{\(\alpha \doteq 70^{\circ }\)}{\(\alpha \doteq 30^{\circ }\)}{\(\alpha \doteq 29^{\circ }\)}}
\End

\Data\ID{9000038705}
\Updated{2022-08-03 10:08:11}
\Level{C}
\SubArea{Triangles}
\IMG{000387_05-figure0.pdf}
\LANGen{
\Question{
The box is on the slope as in the picture. The angle of the slope is
\(\alpha = 45^{\circ }\).
The forces acting on the box are the force of gravity
\(\vec{F_{G}}\) and the
friction \(\vec{F_{t}}\).
The force of gravity can be replaced by two components
\(\vec{F_{1}}\) and
\(\vec{F_{n}}\). (The force
\(\vec{F_{1}}\) is parallel to the slope
and \(\vec{F_{n}}\) is perpendicular to
the slope.) The friction \(F_{t}\) is
given by the formula \(F_{t} = fF_{n}\). The
coefficient of the friction \(f = 0.5\).
Consider the standard acceleration of gravity
\(g = 10\, \mathrm{m\, s^{-2}}\). Find the acceleration of the box.}
{\(a = \frac{5\sqrt{2}} 
 {2} \, \mathrm{m\, s^{-2}}\)}{\(a = 5\sqrt{2}\, \mathrm{m\, s^{-2}}\)}{\(a = 5\sqrt{3}\, \mathrm{m\, s^{-2}}\)}{\(a = 0\, \mathrm{m\, s^{-2}}\)}{\(a = 5\, \mathrm{m\, s^{-2}}\)}{\(a = \frac{5\sqrt{3}} 
 {2} \, \mathrm{m\, s^{-2}}\)}}
\LANGpl{
\Question{
Pudełko znajduje się na równi pochyłej o kącie nachylenia \( \alpha = 45^{\circ }\). Siły działające na pudełko to siła grawitacji \(\vec{F_{G}}\) i siła tarcia \(\vec{F_{t}}\). Siłę grawitacji można zastąpić dwoma składowymi \(\vec{F_{1}}\) i 
\(\vec{F_{n}}\) (Siła \(\vec{F_{1}}\) jest równoległa do powierzchni równi, a \(\vec{F_{n}}\) jest prostopadła do powierzchni równi.) Tarcie \(F_{t}\)  jest podane za pomocą wzoru \(F_{t} = fF_{n}\). Współczynnik tarcia \(f = 0{,}5\).  Rozważ standardowe przyspieszenie grawitacji \(g = 10\, \mathrm{m\, s^{-2}}\). Oblicz przyspieszenie pudełka.}
{\(a = \frac{5\sqrt{2}} 
 {2} \, \mathrm{m\, s^{-2}}\)}{\(a = 5\sqrt{2}\, \mathrm{m\, s^{-2}}\)}{\(a = 5\sqrt{3}\, \mathrm{m\, s^{-2}}\)}{\(a = 0\, \mathrm{m\, s^{-2}}\)}{\(a = 5\, \mathrm{m\, s^{-2}}\)}{\(a = \frac{5\sqrt{3}} 
 {2} \, \mathrm{m\, s^{-2}}\)}}
\LANGcs{
\Question{
Kvádr položíme na nakloněnou rovinu se sklonem
\(\alpha = 45^{\circ }\).
V tíhovém poli Země na něj bude působit tíhová síla
\(\vec{F_{G}}\), síla od
podložky \(\vec{F_{p}}\) a
síla tření \(\vec{F_{t}}\).
Tíhovou sílu můžeme nahradit jejími složkami
\(\vec{F_{1}}\) a
\(\vec{F_{n}}\), kde
\(\vec{F_{1}}\) 
má směr rovnoběžný s nakloněnou rovinou a
\(\vec{F_{n}}\) 
je na ní kolmá. Pro velikost třecí síly platí
\(F_{t} = fF_{n}\). Součinitel
smykového tření \(f = 0{,}5\).
Tíhové zrychlení \(g\doteq 10\, \mathrm{m\, s^{-2}}\). Kvádr se bude pohybovat po nakloněné rovině se zrychlením o velikosti:}
{\(a = \frac{5\sqrt{2}} 
 {2} \, \mathrm{m\, s^{-2}}\)}{\(a = 5\sqrt{2}\, \mathrm{m\, s^{-2}}\)}{\(a = 5\sqrt{3}\, \mathrm{m\, s^{-2}}\)}{\(a = 0\, \mathrm{m\, s^{-2}}\)}{\(a = 5\, \mathrm{m\, s^{-2}}\)}{\(a = \frac{5\sqrt{3}} 
 {2} \, \mathrm{m\, s^{-2}}\)}}
\LANGsk{
\Question{
Kváder položíme na naklonenú rovinu so sklonom
\(\alpha = 45^{\circ }\).
V gravitačnom poli Zeme na ňu bude pôsobiť gravitačná sila
\(\vec{F_{G}}\), sila od
podložky \(\vec{F_{p}}\) a
sila trenia \(\vec{F_{t}}\).
Gravitačnú silu môžeme nahradiť jej zložkami
\(\vec{F_{1}}\) a
\(\vec{F_{n}}\), kde
\(\vec{F_{1}}\) 
má smer rovnobežný s naklonenou rovinou a
\(\vec{F_{n}}\) 
je na ňu kolmá. Pre veľkosť trecej sily platí
\(F_{t} = fF_{n}\). Súčiniteľ
šmykového trenia \(f = 0{,}5\).
Gravitačné zrýchlenie \(g\doteq 10\, \mathrm{m\, s^{-2}}\). Kváder sa bude pohybovať po naklonenej rovine so zrýchlením o veľkosti:}
{\(a = \frac{5\sqrt{2}} 
 {2} \, \mathrm{m\, s^{-2}}\)}{\(a = 5\sqrt{2}\, \mathrm{m\, s^{-2}}\)}{\(a = 5\sqrt{3}\, \mathrm{m\, s^{-2}}\)}{\(a = 0\, \mathrm{m\, s^{-2}}\)}{\(a = 5\, \mathrm{m\, s^{-2}}\)}{\(a = \frac{5\sqrt{3}} 
 {2} \, \mathrm{m\, s^{-2}}\)}}
\LANGes{
\Question{
Un ortoedro se encuentra en un plano inclinado (observa la imagen). El ángulo de la pendiente es
\(\alpha = 45^{\circ }\).
Las fuerzas que actúan sobre el ortoedro son la fuerza de la gravedad
\(\vec{F_{G}}\) y la de la fricción \(\vec{F_{t}}\).
La fuerza de gravedad se puede reemplazar por dos componentes
\(\vec{F_{1}}\) y
\(\vec{F_{n}}\). (La fuerza
\(\vec{F_{1}}\)es paralela a la pendiente y \(\vec{F_{n}}\) es perpendicular a la pendiente). La fricción \(F_{t}\) viene dada por la fórmula \(F_{t} = fF_{n}\). El coeficiente de la fricción es \(f = 0.5\).
Consideramos la aceleración estándar de la gravedad
\(g = 10\, \mathrm{m\, s^{-2}}\). Halla la aceleración del ortoedro.}
{\(a = \frac{5\sqrt{2}} 
 {2} \, \mathrm{m\, s^{-2}}\)}{\(a = 5\sqrt{2}\, \mathrm{m\, s^{-2}}\)}{\(a = 5\sqrt{3}\, \mathrm{m\, s^{-2}}\)}{\(a = 0\, \mathrm{m\, s^{-2}}\)}{\(a = 5\, \mathrm{m\, s^{-2}}\)}{\(a = \frac{5\sqrt{3}} 
 {2} \, \mathrm{m\, s^{-2}}\)}}
\End

\Data\ID{9000038706}
\Updated{2022-08-03 10:08:33}
\Level{C}
\SubArea{Triangles}
\IMG{000387_06-figure0.pdf}
\LANGen{
\Question{
The box is on the slope as in the picture. The angle of the slope is
\(\alpha \).
The forces acting on the box are the force of gravity
\(\vec{F_{G}}\) and the
friction \(\vec{F_{t}}\).
The force of gravity can be replaced by two components
\(\vec{F_{1}}\) and
\(\vec{F_{n}}\). (The force
\(\vec{F_{1}}\) is parallel to the slope
and \(\vec{F_{n}}\) is perpendicular to
the slope.) The friction \(F_{t}\) is
given by the formula \(F_{t} = fF_{n}\). The
coefficient of the friction is \(f = 0.47\).
Consider the standard acceleration of gravity
\(g = 10\, \mathrm{m\, s^{-2}}\). Find the angle \(\alpha \) 
which ensures that the box moves on the slope with zero acceleration.}
{\(\alpha \doteq 25^{\circ }\)}{\(\alpha \doteq 15^{\circ }\)}{\(\alpha \doteq 20^{\circ }\)}{\(\alpha \doteq 65^{\circ }\)}{\(\alpha \doteq 28^{\circ }\)}{\(\alpha \doteq 62^{\circ }\)}}
\LANGpl{
\Question{
Pudełko znajduje się na równi pochyłej o kącie nachylenia \( \alpha \) (jak na zdjęciu). Siły działające na pudełko to siła grawitacji \(\vec{F_{G}}\) i tarcia \(\vec{F_{t}}\). Siłę grawitacji można zastąpić dwoma składowymi  \(\vec{F_{1}}\) i 
\(\vec{F_{n}}\). (Siła \(\vec{F_{1}}\)  jest równoległa do powierzchni równi, a \(\vec{F_{n}}\)  jest prostopadła do powierzchni równi.) Tarcie \(F_{t}\) jest podane za pomocą wzoru \(F_{t} = fF_{n}\). Współczynnik tarcia \(f = 0{,}47\). Rozważ standardowe przyspieszenie grawitacji \(g = 10\, \mathrm{m s^{-2}}\). Znajdź kąt \(\alpha \), który zapewni, że pudełko przesunie się po równi pochyłej z zerowym przyspieszeniem.}
{\(\alpha \doteq 25^{\circ }\)}{\(\alpha \doteq 15^{\circ }\)}{\(\alpha \doteq 20^{\circ }\)}{\(\alpha \doteq 65^{\circ }\)}{\(\alpha \doteq 28^{\circ }\)}{\(\alpha \doteq 62^{\circ }\)}}
\LANGcs{
\Question{
Kvádr položíme na nakloněnou rovinu se sklonem
\(\alpha \).
V tíhovém poli Země na něj bude působit tíhová síla
\(\vec{F_{G}}\), síla od
podložky \(\vec{F_{p}}\) a
síla tření \(\vec{F_{t}}\).
Tíhovou sílu můžeme nahradit jejími složkami
\(\vec{F_{1}}\) a
\(\vec{F_{n}}\), kde
\(\vec{F_{1}}\) 
má směr rovnoběžný s nakloněnou rovinou a
\(\vec{F_{n}}\) 
je na ní kolmá. Pro velikost třecí síly platí
\(F_{t} = fF_{n}\). Součinitel
smykového tření \(f = 0{,}47\).
Tíhové zrychlení \(g\doteq 10\, \mathrm{m\, s^{-2}}\). Při jakém úhlu \(\alpha \) 
se může kvádr po nakloněné rovině pohybovat rovnoměrně?}
{\(\alpha \doteq 25^{\circ }\)}{\(\alpha \doteq 15^{\circ }\)}{\(\alpha \doteq 20^{\circ }\)}{\(\alpha \doteq 65^{\circ }\)}{\(\alpha \doteq 28^{\circ }\)}{\(\alpha \doteq 62^{\circ }\)}}
\LANGsk{
\Question{
Kváder položíme na naklonenú rovinu so sklonom
\(\alpha \).
V gravitačnom poli Zeme na neho bude pôsobiť gravitačná sila
\(\vec{F_{G}}\), sila od
podložky \(\vec{F_{p}}\) a
sila trenia \(\vec{F_{t}}\).
Gravitačnú silu môžeme nahradiť jej zložkami
\(\vec{F_{1}}\) a
\(\vec{F_{n}}\), kde
\(\vec{F_{1}}\) 
má smer rovnobežný s naklonenou rovinou a
\(\vec{F_{n}}\) 
je na ňu kolmá. Pre veľkosť trecej sily platí
\(F_{t} = fF_{n}\). Súčiniteľ
šmykového trenia \(f = 0{,}47\).
Gravitačné zrýchlenie \(g\doteq 10\, \mathrm{m\, s^{-2}}\). Pri akom uhle \(\alpha \) 
sa môže kváder po naklonenej rovine pohybovať rovnomerne?}
{\(\alpha \doteq 25^{\circ }\)}{\(\alpha \doteq 15^{\circ }\)}{\(\alpha \doteq 20^{\circ }\)}{\(\alpha \doteq 65^{\circ }\)}{\(\alpha \doteq 28^{\circ }\)}{\(\alpha \doteq 62^{\circ }\)}}
\LANGes{
\Question{
Un ortoedro se encuentra en un plano inclinado (observa la imagen). El ángulo de la pendiente es
\(\alpha \).
Las fuerzas que actúan sobre el ortoedro son la fuerza de la gravedad
\(\vec{F_{G}}\) y la fricción \(\vec{F_{t}}\).
La fuerza de gravedad se puede reemplazar por dos componentes
\(\vec{F_{1}}\) y
\(\vec{F_{n}}\). (La fuerza
\(\vec{F_{1}}\) es paralela a la pendiente y \(\vec{F_{n}}\) es perpendicular a la pendiente). La fricción \(F_{t}\) viene dada por la fórmula \(F_{t} = fF_{n}\). El coeficiente de fricción es \(f = 0.47\).
Consideramos la aceleración estándar de la gravedad
\(g = 10\, \mathrm{m\, s^{-2}}\). Halla el ángulo \(\alpha \) para que el ortoedro se mueva en el plano inclinado con aceleración cero.}
{\(\alpha \doteq 25^{\circ }\)}{\(\alpha \doteq 15^{\circ }\)}{\(\alpha \doteq 20^{\circ }\)}{\(\alpha \doteq 65^{\circ }\)}{\(\alpha \doteq 28^{\circ }\)}{\(\alpha \doteq 62^{\circ }\)}}
\End

\Data\ID{9000038707}
\Updated{2022-08-03 10:08:01}
\Level{C}
\SubArea{Triangles}
\IMG{000387_07-figure0.pdf}
\LANGen{
\Question{
The box is on the slope as in the picture. The length of the slope is
\(l = 2\, \mathrm{m}\) and the
height is \(h = 1.2\, \mathrm{m}\).
The forces acting on the box are the force of gravity
\(\vec{F_{G}}\) and the
friction \(\vec{F_{t}}\).
The force of gravity can be replaced by two components
\(\vec{F_{1}}\) and
\(\vec{F_{n}}\). (The force
\(\vec{F_{1}}\) is parallel to the slope
and \(\vec{F_{n}}\) is perpendicular to
the slope.) The friction \(F_{t}\) 
is given by the formula \(F_{t} = fF_{n}\) 
where \(f\) 
is the coefficient of the friction. Consider the standard acceleration of gravity
\(g = 10\, \mathrm{m\, s^{-2}}\). Find the minimal value for the coefficient of the friction
\(f\) to
ensure that the box does not move with an acceleration.}
{\(f = 0.75\)}{\(f = 0.6\)}{\(f = 0.65\)}{\(f = 0.7\)}{\(f = 0.55\)}{\(f = 0.8\)}}
\LANGpl{
\Question{
Pudełko znajduje się na równi pochyłej (jak na zdjęciu). Długość równi pochyłej \(l = 2\, \mathrm{m}\), a wysokość \(h = 1.2\, \mathrm{m}\). Siły działające na skrzynkę to siła grawitacji \(\vec{F_{G}}\) i tarcia
\(\vec{F_{t}}\). Siłę grawitacji można zastąpić dwoma składowymi \(\vec{F_{1}}\) i
\(\vec{F_{n}}\). (Siła \(\vec{F_{1}}\) jest równoległa do powierzchni równi, a \(\vec{F_{n}}\) jest prostopadła do powierzchni równi.) Tarcie \(F_{t}\)  jest podane za pomocą wzoru \(F_{t} = fF_{n}\), gdzie \(f\) jest współczynnikiem tarcia. Rozważ standardowe przyspieszenie grawitacji \(g = 10\, \mathrm{m\, s^{-2}}\). Znajdź minimalną wartość współczynnika tarcia \(f\), aby upewnić się, że pudełko nie porusza się z przyspieszeniem.}
{\(f = 0.75\)}{\(f = 0.6\)}{\(f = 0.65\)}{\(f = 0.7\)}{\(f = 0.55\)}{\(f = 0.8\)}}
\LANGcs{
\Question{
Kvádr položíme na nakloněnou rovinu o délce
\(l = 2\, \mathrm{m}\) a
výšce \(h = 1{,}2\, \mathrm{m}\).
V tíhovém poli Země na něj bude působit tíhová síla
\(\vec{F_{G}}\), síla od
podložky \(\vec{F_{p}}\) a
síla tření \(\vec{F_{t}}\).
Tíhovou sílu můžeme nahradit jejími složkami
\(\vec{F_{1}}\) a
\(\vec{F_{n}}\), kde
\(\vec{F_{1}}\) 
má směr rovnoběžný s nakloněnou rovinou a
\(\vec{F_{n}}\) 
je na ní kolmá. Pro velikost třecí síly platí
\(F_{t} = fF_{n}\), kde
\(f\) 
je součinitel smykového tření. Tíhové zrychlení
\(g\doteq 10\, \mathrm{m\, s^{-2}}\). Jak velký musí být součinitel smykového tření
\(f\), aby
se kvádr nepohyboval zrychleně? Musel by být alespoň:}
{\(f = 0{,}75\)}{\(f = 0{,}6\)}{\(f = 0{,}65\)}{\(f = 0{,}7\)}{\(f = 0{,}55\)}{\(f = 0{,}8\)}}
\LANGsk{
\Question{
Kváder položíme na naklonenú rovinu o dĺžke
\(l = 2\, \mathrm{m}\) a
výške \(h = 1{,}2\, \mathrm{m}\).
V gravitačnom poli Zeme na neho bude pôsobiť gravitačná sila
\(\vec{F_{G}}\), sila od
podložky \(\vec{F_{p}}\) a
sila trenia \(\vec{F_{t}}\).
Gravitačnú silu môžeme nahradiť jej zložkami
\(\vec{F_{1}}\) a
\(\vec{F_{n}}\), kde
\(\vec{F_{1}}\) 
má smer rovnobežný s naklonenou rovinou a
\(\vec{F_{n}}\) 
je na ňu kolmá. Pre veľkosť trecej sily platí
\(F_{t} = fF_{n}\), kde
\(f\) 
je súčiniteľ šmykového trenia. Gravitačné zrýchlenie
\(g\doteq 10\, \mathrm{m\, s^{-2}}\). Aký veľký musí byť súčiniteľ šmykového trenia
\(f\), aby
sa kváder nepohyboval zrýchlene? Musel by byť aspoň:}
{\(f = 0{,}75\)}{\(f = 0{,}6\)}{\(f = 0{,}65\)}{\(f = 0{,}7\)}{\(f = 0{,}55\)}{\(f = 0{,}8\)}}
\LANGes{
\Question{
Un ortoedro se encuentra en un plano inclinado (observa la imagen). La longitud del plano inclinado es
\(l = 2\, \mathrm{m}\) y su altura es \(h = 1.2\, \mathrm{m}\).
Las fuerzas que actúan sobre el ortoedro son la fuerza de gravedad
\(\vec{F_{G}}\) y la fricción \(\vec{F_{t}}\).
La fuerza de gravedad se puede reemplazar por dos componentes
\(\vec{F_{1}}\) y
\(\vec{F_{n}}\). (La fuerza
\(\vec{F_{1}}\) es paralela a la pendiente y \(\vec{F_{n}}\) es perpendicular a la pendiente). La fricción \(F_{t}\) 
viene dada por la fórmula \(F_{t} = fF_{n}\) 
donde \(f\) 
es el coeficiente de fricción. Consideramos la aceleración estándar de la gravedad
\(g = 10\, \mathrm{m\, s^{-2}}\). Halla el valor mínimo de coeficiente de  fricción
\(f\) para que el ortoedro no se mueva con aceleración.}
{\(f = 0.75\)}{\(f = 0.6\)}{\(f = 0.65\)}{\(f = 0.7\)}{\(f = 0.55\)}{\(f = 0.8\)}}
\End

\Data\ID{9000124501}
\Updated{2022-08-03 10:08:39}
\Level{C}
\SubArea{Triangles}
\IMG{001245_01-figure0.pdf}
\LANGen{
\Question{
Similar triangles can be used to estimate the
distance from a distant object of a given width. Consider a door of the width
\(85\, \mathrm{cm}\).
A man stands in an unknown distance
from the door and holds a thin pencil
vertically in his arm in the distance
\(35\, \mathrm{cm}\) 
from his face. If he closes the left eye, the right eye, the
pencil and the left side of the door are aligned in
one line. In a similar way, his left eye, the pencil
and the right hand side of the door are also aligned
in one line, which is apparent when closing the
right eye. Assuming the distance
\(6\, \mathrm{cm}\) 
between his eyes, estimate the distance from the
man to the door. Give your answer in meters and
round to one decimal place.}
{\(5.3\, \mathrm{m}\)}{\(5.0\, \mathrm{m}\)}{\(0.5\, \mathrm{m}\)}{\(4.5\, \mathrm{m}\)}{}{}}
\LANGpl{
\Question{
Podobne trójkąty można wykorzystać do oszacowania odległości od odległego obiektu o danej szerokości. Rozważ drzwi o szerokości \(85\, \mathrm{cm}\). Mężczyzna stoi w nieznanej odległości od drzwi i trzyma cienki ołówek pionowo w ręce w odległości  \(35\, \mathrm{cm}\)  od twarzy. Jeśli zamyka lewe oko, prawe oko, ołówek i lewa strona drzwi są wyrównane w jednej linii. W podobny sposób jego lewe oko, ołówek i prawa strona drzwi są również wyrównane w jednej linii, co jest widoczne przy zamykaniu prawego oka. Zakładając odległość między oczami wynosi \(6\, \mathrm{cm}\),  oszacuj odległość  mężczyzny od drzwi. Podaj odpowiedź w metrach i zaokrąglij do jednego miejsca po przecinku.}
{\(5.3\, \mathrm{m}\)}{\(5.0\, \mathrm{m}\)}{\(0.5\, \mathrm{m}\)}{\(4.5\, \mathrm{m}\)}{}{}}
\LANGcs{
\Question{
Když držíme ve vzdálenosti \(35\, \mathrm{cm}\) 
před obličejem tužku (ve svislé poloze) a díváme se střídavě pravým a
levým okem, zjistíme, že při pohledu pravým okem se tužka kryje s levou
hranou dveří a při pohledu levým okem se kryje s pravou hranou dveří. V jaké
vzdálenosti před dveřmi stojíme, je-li vzdálenost mezi očima (zornicemi)
\(6\, \mathrm{cm}\) a dveře mají
standardizovanou šířku \(85\, \mathrm{cm}\)?
Výsledek zaokrouhlete na desetiny metru.}
{\(5{,}3\, \mathrm{m}\)}{\(5\, \mathrm{m}\)}{\(0{,}5\, \mathrm{m}\)}{\(4{,}5\, \mathrm{m}\)}{}{}}
\LANGsk{
\Question{
Keď držíme vo vzdialenosti \(35\, \mathrm{cm}\) 
pred tvárou ceruzku (vo zvislej polohe) a pozeráme sa striedavo pravým a
ľavým okom, zistíme, že pri pohľade pravým okom sa ceruzka kryje s ľavou
hranou dverí a pri pohľade ľavým okom sa kryje s pravou hranou dverí. V akej
vzdialenosti pred dverami stojíme, ak je vzdialenosť medzi očami (zrenicami)
\(6\, \mathrm{cm}\) a dvere majú
štandardnú šírku \(85\, \mathrm{cm}\)?
Výsledok zaokrúhlite na desatiny metra.}
{\(5{,}3\, \mathrm{m}\)}{\(5\, \mathrm{m}\)}{\(0{,}5\, \mathrm{m}\)}{\(4{,}5\, \mathrm{m}\)}{}{}}
\LANGes{
\Question{
Es posible utilizar triángulos semejantes para calcular la distancia desde un objeto a otro. Considera una puerta con un ancho de \(85\, \mathrm{cm}\). Un hombre se encuentra a una distancia desconocida de la puerta y tiene un lápiz en la mano a una distancia de \(35\, \mathrm{cm}\) de su cara. Cerrando el ojo izquierdo, el ojo derecho, el lápiz y la parte izquierda de la puerta están alineados. Cerrando el ojo derecho, el ojo izquierdo, el lápiz y la parte derecha de la puerta también están alineados. Suponiendo que la distancia entre sus ojos es \(6\, \mathrm{cm}\), calcula la distancia desde el hombre hasta la puerta. Expresa el resultado en metros y redóndealo a un decimal.}
{\(5.3\, \mathrm{m}\)}{\(5.0\, \mathrm{m}\)}{\(0.5\, \mathrm{m}\)}{\(4.5\, \mathrm{m}\)}{}{}}
\End

\Data\ID{9000124503}
\Updated{2021-12-13 07:12:56}
\Level{C}
\SubArea{Triangles}
\IMG{001245_03-figure0.pdf}
\LANGen{
\Question{
A tall radio mast is attached by several cables. The length of each cable is
\(30\, \mathrm{m}\) and all cables
are attached \(2\, \mathrm{m}\) 
under the top of the mast. The second end of the cable is anchored to the ground. The cable is
in the height \(6\, \mathrm{m}\) 
if measured directly above the point which is in the distance
\(8\, \mathrm{m}\) from
the point where the cable is anchored to the ground. Find the height of the
mast.}
{\(20\, \mathrm{m}\)}{\(24\, \mathrm{m}\)}{\(22.5\, \mathrm{m}\)}{\(24.5\, \mathrm{m}\)}{}{}}
\LANGpl{
\Question{
Wysoki maszt radiowy jest zamocowany kilkoma linami. Każda lina ma długość \(30\, \mathrm{m}\), a wszystkie liny są przymocowane \(2\, \mathrm{m}\)  pod szczytem masztu. Drugi koniec liny jest zakotwiczony na ziemi. Lina znajduje się na wysokości \(6\, \mathrm{m}\), jeśli mierzona jest bezpośrednio nad punktem znajdującym się w odległości \(8\, \mathrm{m}\) od punktu, w którym lina jest zakotwiczona na ziemi. Oblicz wysokość masztu.}
{\(20\, \mathrm{m}\)}{\(24\, \mathrm{m}\)}{\(22.5\, \mathrm{m}\)}{\(24.5\, \mathrm{m}\)}{}{}}
\LANGcs{
\Question{
Stožár vysílače je ukotven několika lany. Každé z kotvících lan má délku
\(30\, \mathrm{m}\) a je
upevněno \(2\, \mathrm{m}\) 
pod vrcholem vysílače. Druhý konec lana je upevněn na zemi v neznámé
vzdálenosti od vysílače. Jak vysoký je vysílač, víme-li, že ve vzdálenosti
\(8\, \mathrm{m}\) 
od ukotvení lana na zemi je toto lano ve výšce
\(6\, \mathrm{m}\).}
{\(20\, \mathrm{m}\)}{\(24\, \mathrm{m}\)}{\(22{,}5\, \mathrm{m}\)}{\(24{,}5\, \mathrm{m}\)}{}{}}
\LANGsk{
\Question{
Stožiar vysielača je ukotvený niekoľkými lanami. Každé z kotviacich lán má dĺžku
\(30\, \mathrm{m}\) a je
upevnené \(2\, \mathrm{m}\) 
pod vrcholom vysielača. Druhý koniec lana je upevnený na zemi v neznámej
vzdialenosti od vysielača. Aký vysoký je vysielač, ak vieme, že vo vzdialenosti
\(8\, \mathrm{m}\) 
od ukotvenia lana na zemi je toto lano vo výške
\(6\, \mathrm{m}\).}
{\(20\, \mathrm{m}\)}{\(24\, \mathrm{m}\)}{\(22{,}5\, \mathrm{m}\)}{\(24{,}5\, \mathrm{m}\)}{}{}}
\LANGes{
\Question{
Un mástil de radio está atado por varios cables. Cada cable mide
\(30\, \mathrm{m}\) y todos
están atados \(2\, \mathrm{m}\) 
por bajo del punto superior del mástil. El segundo extremo del cable está anclado al suelo. El cable está a una altura de \(6\, \mathrm{m}\) 
si se mide directamente sobre el punto que está a una distancia de
\(8\, \mathrm{m}\) desde donde el cable está anclado al suelo. Halla la altura del mástil.}
{\(20\, \mathrm{m}\)}{\(24\, \mathrm{m}\)}{\(22.5\, \mathrm{m}\)}{\(24.5\, \mathrm{m}\)}{}{}}
\End

\Data\ID{9000124504}
\Updated{2022-08-03 10:08:16}
\Level{C}
\SubArea{Triangles}
\IMG{001245_04-figure0.pdf}
\LANGen{
\Question{
A force due to gravity on a body is
\(1\: 800\, \mathrm{N}\). This body has to be
lifted to the height \(50\, \mathrm{cm}\) 
using a slope. The maximal force which can be used to lift the body is
\(600\, \mathrm{N}\).
Neglect the friction and find the minimal length of the slope required to accomplish
this task.Hint: The force due to gravity can be decomposed into two directions. The normal force
\(F_{1}\) is compensated by the reaction
of the slope. The force \(F_{2}\) 
parallel to the slope is required to overcome if we wish to lift the body (see the picture).}
{\(\frac{3}
{2}\, \mathrm{m}\)}{\(\frac{2}
{3}\, \mathrm{m}\)}{\(\frac{1}
{6}\, \mathrm{m}\)}{\(\frac{20}
 {9} \, \mathrm{m}\)}{}{}}
\LANGpl{
\Question{
Siła spowodowana grawitacją ciała wynosi \(1\: 800\, \mathrm{N}\). Ciało należy podnieść do wysokości \(50\, \mathrm{cm}\) 
za pomocą równi pochyłej. Maksymalna siła, która może zostać użyta do podniesienia ciała, wynosi \(600\, \mathrm{N}\). Pomijając tarcie  znajdź minimalną długość zbocza wymaganą do wykonania tego zadania. Wskazówka: Siłę grawitacji można rozłożyć na dwa kierunki. Siła nacisku \(F_{1}\) jest kompensowana przez siłę reakcji podloża. Siła  \(F_{2}\)  równoległa do równi powierzchni jest wymagana do pokonania jeśli chcemy podnieść ciało (zobacz rysunek).}
{\(\frac{3}
{2}\, \mathrm{m}\)}{\(\frac{2}
{3}\, \mathrm{m}\)}{\(\frac{1}
{6}\, \mathrm{m}\)}{\(\frac{20}
 {9} \, \mathrm{m}\)}{}{}}
\LANGcs{
\Question{
Maximální síla, kterou jsem schopen vyvinout je
\(600\, \mathrm{N}\). Jakou
nejmenší délku musí mít nakloněná rovina, abych pomocí ní dokázal těleso o
tíze \(1\: 800\, \mathrm{N}\) zvednout
do výšky \(50\, \mathrm{cm}\)?
Tření mezi posouvaným tělesem a nakloněnou rovinou zanedbáváme.
(Nápověda: Tíhová síla tělesa na nakloněné rovině se rozloží na dvě
navzájem kolmé složky. Při posunu tělesa po nakloněné rovině musíme překonat
složku \(F_{2}\) 
(viz obrázek).}
{\(\frac{3}
{2}\, \mathrm{m}\)}{\(\frac{2}
{3}\, \mathrm{m}\)}{\(\frac{1}
{6}\, \mathrm{m}\)}{\(\frac{20}
 {9} \, \mathrm{m}\)}{}{}}
\LANGsk{
\Question{
Maximálna sila, ktorú som schopný vyvinúť je
\(600\, \mathrm{N}\). Akú
najmenšiu dĺžku musí mať naklonená rovina, aby pomocou nej dokázalo teleso o
hmotnosti \(1\: 800\, \mathrm{N}\) zdvihnúť
do výšky \(50\, \mathrm{cm}\)?
Trenie medzi posúvaným telesom a naklonenou rovinou zanedbávame.
(Nápoveda: Tiažová sila telesa na naklonenej rovine sa rozloží na dve
navzájom kolmé zložky. Pri posune telesa po naklonenej rovine musíme prekonať
zložku \(F_{2}\) 
(viď obrázok).}
{\(\frac{3}
{2}\, \mathrm{m}\)}{\(\frac{2}
{3}\, \mathrm{m}\)}{\(\frac{1}
{6}\, \mathrm{m}\)}{\(\frac{20}
 {9} \, \mathrm{m}\)}{}{}}
\LANGes{
\Question{
La fuerza de gravedad que actúa sobre un cuerpo  es
\(1\: 800\, \mathrm{N}\). Este cuerpo lo tenemos que elevar hasta una altura de \(50\, \mathrm{cm}\) 
usando un plano inclinado. La fuerza máxima que se puede usar para levantar el cuerpo es
\(600\, \mathrm{N}\).
Omitiendo la fricción, halla la longitud mínima de la pendiente requerida para lograr esta tarea.}
{\(\frac{3}
{2}\, \mathrm{m}\)}{\(\frac{2}
{3}\, \mathrm{m}\)}{\(\frac{1}
{6}\, \mathrm{m}\)}{\(\frac{20}
 {9} \, \mathrm{m}\)}{}{}}
\End

\Data\ID{9000124505}
\Updated{2022-08-03 10:08:20}
\Level{C}
\SubArea{Triangles}
\IMG{001245_05-figure0.pdf}
\LANGen{
\Question{
The picture shows the virtual image \(y'\) 
of the object \(y\) as created by
a concave lens. The points \(F\) 
and \(F'\) 
are focal points of the lens. The distance from the lens to each of the focal points is
\(20\, \mathrm{cm}\). The
object \(y\) is
\(25\, \, \mathrm{cm}\) height and it is
in the distance \(50\, \mathrm{cm}\) 
from the lens. Find the height of the virtual image
\(y'\).}
{\(\frac{50}
 {7} \, \mathrm{cm}\)}{\(10\, \mathrm{cm}\)}{\(\frac{50}
 {3} \, \mathrm{cm}\)}{\(\frac{175}
 {2} \, \mathrm{cm}\)}{}{}}
\LANGpl{
\Question{
Zdjęcie przedstawia wirtualny obraz \(y'\) obiektu \(y\) utworzonego przy użycie soczewki rozpraszającej. Punkty \(F\) 
i \(F'\)  są punktami ogniskowymi soczewki. Odległość środka soczewki od każdego ogniska wynosi \(20\, \mathrm{cm}\). Obiekt \(y\) ma \(25\, \, \mathrm{cm}\) wysokości i znajduje się w odległości  \(50\, \mathrm{cm}\)  od soczewki. Znajdź wysokość wirtualnego obrazu \(y'\).}
{\(\frac{50}
 {7} \, \mathrm{cm}\)}{\(10\, \mathrm{cm}\)}{\(\frac{50}
 {3} \, \mathrm{cm}\)}{\(\frac{175}
 {2} \, \mathrm{cm}\)}{}{}}
\LANGcs{
\Question{
Na obrázku je zakresleno zobrazení předmětu pomocí tenké rozptylné čočky.
Body \(F\) 
a \(F'\) jsou
tzv. ohniska čočky. Vzdálenost ohniska od čočky je tzv. ohnisková vzdálenost
\(f\). Předmět o velikosti
\(25\, \mathrm{cm}\) (\(y\))
a vzdálený \(50\, \mathrm{cm}\) 
(\(a\))
od čočky zobrazíme čočkou, jejíž ohnisková vzdálenost
\(f\) je
\(20\, \mathrm{cm}\). Jaká bude
velikost \(y'\) 
vytvořeného obrazu? (Poznámka: Ve fyzice označujeme ohniskové vzdálenosti
rozptylných čoček záporným číslem.)}
{\(\frac{50}
 {7} \, \mathrm{cm}\)}{\(10\, \mathrm{cm}\)}{\(\frac{50}
 {3} \, \mathrm{cm}\)}{\(\frac{175}
 {2} \, \mathrm{cm}\)}{}{}}
\LANGsk{
\Question{
Na obrázku je zakreslené zobrazenie predmetu pomocou tenkej rozptylnej šošovky.
Body \(F\) 
a \(F'\) sú
tzv. ohniská šošovky. Vzdialenosť ohniska od šošovky je tzv. ohnisková vzdialenosť
\(f\). Predmet o veľkosti
\(25\, \mathrm{cm}\) (\(y\))
a vzdialený \(50\, \mathrm{cm}\) 
(\(a\))
od šošovky zobrazíme šošovkou, ktorej ohnisková vzdialenosť
\(f\) je
\(20\, \mathrm{cm}\). Aká bude
veľkosť \(y'\) 
vytvoreného obrazu? (Poznámka: Vo fyzike označujeme ohniskové vzdialenosti
rozptylných šošoviek záporným číslom.)}
{\(\frac{50}
 {7} \, \mathrm{cm}\)}{\(10\, \mathrm{cm}\)}{\(\frac{50}
 {3} \, \mathrm{cm}\)}{\(\frac{175}
 {2} \, \mathrm{cm}\)}{}{}}
\LANGes{
\Question{
La imagen representa una imagen virtual \(y'\) 
del objeto \(y\) creado por una lente cóncava. Los puntos \(F\) 
y \(F'\) son focos de la lente. La distancia desde la lente a cada uno de los focos es de
\(20\, \mathrm{cm}\). La altura del objeto \(y\) es
\(25\, \, \mathrm{cm}\) y está a una distancia de \(50\, \mathrm{cm}\) 
desde la lente. Halla la altura de la imagen virtual \(y'\).}
{\(\frac{50}
 {7} \, \mathrm{cm}\)}{\(10\, \mathrm{cm}\)}{\(\frac{50}
 {3} \, \mathrm{cm}\)}{\(\frac{175}
 {2} \, \mathrm{cm}\)}{}{}}
\End

\Data\ID{9000150501}
\Updated{2022-04-23 05:04:05}
\Level{C}
\SubArea{Triangles}
\LANGen{
\Question{
A man of height \(180\, \mathrm{cm}\) casts a \(200\, \mathrm{cm}\) shadow. At the same moment, a tree of an unknown height casts a \(35\, \mathrm{m}\) shadow. Find the height of the tree.}
{\(\frac{63}
 {2} \, \mathrm{m}\)}{\(\frac{350}
 {9} \, \mathrm{m}\)}{\(\frac{72}
 {7} \, \mathrm{m}\)}{\(\frac{36}
{35}\, \mathrm{m}\)}{}{}}
\LANGpl{
\Question{
Mężczyzna o wzroście  \(180\, \mathrm{cm}\) rzuca \(200\, \mathrm{cm}\) cień. Drzewo o niewiadomej wysokości rzuca  \(35\, \mathrm{m}\) 
cień w tym samym czasie. Oblicz wysokość drzewa.}
{\(\frac{63}
 {2} \, \mathrm{m}\)}{\(\frac{350}
 {9} \, \mathrm{m}\)}{\(\frac{72}
 {7} \, \mathrm{m}\)}{\(\frac{36}
{35}\, \mathrm{m}\)}{}{}}
\LANGcs{
\Question{
Jak vysoký je strom, jestliže vrhá stín dlouhý
\(35\, \mathrm{m}\)? Ve stejnou dobu
vrhá \(180\, \mathrm{cm}\) vysoká postava
stín o délce \(200\, \mathrm{cm}\).}
{\(\frac{63}
 {2} \, \mathrm{m}\)}{\(\frac{350}
 {9} \, \mathrm{m}\)}{\(\frac{72}
 {7} \, \mathrm{m}\)}{\(\frac{36}
{35}\, \mathrm{m}\)}{}{}}
\LANGsk{
\Question{
Aký vysoký je strom, ak vrhá tieň dlhý
\(35\, \mathrm{m}\)? V rovnakej dobe
vrhá \(180\, \mathrm{cm}\) vysoká postava
tieň o dĺžke \(200\, \mathrm{cm}\).}
{\(\frac{63}
 {2} \, \mathrm{m}\)}{\(\frac{350}
 {9} \, \mathrm{m}\)}{\(\frac{72}
 {7} \, \mathrm{m}\)}{\(\frac{36}
{35}\, \mathrm{m}\)}{}{}}
\LANGes{
\Question{
Un hombre que mide \(180\, \mathrm{cm}\) proyecta una sombra de \(200\, \mathrm{cm}\). Un árbol de una altura desconocida proyecta una sombra de \(35\, \mathrm{m}\). Calcula la altura del árbol.}
{\(\frac{63}
 {2} \, \mathrm{m}\)}{\(\frac{350}
 {9} \, \mathrm{m}\)}{\(\frac{72}
 {7} \, \mathrm{m}\)}{\(\frac{36}
{35}\, \mathrm{m}\)}{}{}}
\End

\Data\ID{9000150503}
\Updated{2022-08-03 10:08:21}
\Level{C}
\SubArea{Triangles}
\IMG{001505_03-figure0.pdf}
\LANGen{
\Question{
A pendulum constituted of a rope of the length
\(l\) and a
body is displaced from it's equilibrium. The force due to gravity on the body
\(F_{g} = 20\, \mathrm{N}\). The body
is higher by \(h = 10\, \mathrm{cm}\) 
in the displaced position (comparing to the equilibrium
position). The tension in the rope in the displaced position is
\(F_{1} = 12\, \mathrm{N}\). Find the length
of the rope \(l\). Hint: Using a parallelogram, the force of gravity on the body can be decomposed into a force
\(F_{1}\) in the direction
of the rope and \(F_{2}\) 
in the perpendicular direction.}
{\(25\, \mathrm{cm}\)}{\(25\, \mathrm{m}\)}{\(6\, \mathrm{cm}\)}{\(16\frac{2}
{3}\, \mathrm{cm}\)}{}{}}
\LANGpl{
\Question{
Wahadło składające się ciała zawieszanego na nici o długości \(l\) jest wychylone z położenia równowagi tak, że ciało jest na wysokości \(h = 10\, \mathrm{cm}\) (w porównaniu do pozycją równowagi). Siła grawitacji \(F_{g} = 20\, \mathrm{N}\) i składowa radialna (skierowana przeciwnie do siły naprężenia nici) \(F_{1}=10\,\mathrm{N}\). Wyznacz długość nici \(l\). Wskazówka: Używając równoległoboku siła grawitacji na ciało może zostać rozłożona na siłę  \(F_{1}\) w kierunku nici  i \(F_{2}\)   w kierunku prostopadłym.}
{\(25\, \mathrm{cm}\)}{\(25\, \mathrm{m}\)}{\(6\, \mathrm{cm}\)}{\(16\frac{2}
{3}\, \mathrm{cm}\)}{}{}}
\LANGcs{
\Question{
Na vlákno zavěsíme těleso o tíze
\(20\, \mathrm{N}\) 
(\(F_{g}\)) a takto
vzniklé kyvadlo vychýlíme. Vychýlením kyvadla se zvýší poloha tělesa nad
podložkou o \(10\, \mathrm{cm}\) 
(\(h\)). V této poloze je
vlákno napínáno silou \(12\, \mathrm{N}\) 
(\(F_{1}\)). Určete délku
vlákna (\(l\)).
(Nápověda: Tíha zavěšeného tělesa se rozloží na síly
\(F_{1}\) a
\(F_{2}\) (složky tíhové
síly). Síla \(F_{1}\) způsobuje
napínání vlákna a \(F_{2}\) 
vrací kyvadlo do svislé polohy. Rozklad sil se provádí pomocí tzv. rovnoběžníku
sil.)}
{\(25\, \mathrm{cm}\)}{\(25\, \mathrm{m}\)}{\(6\, \mathrm{cm}\)}{\(16\frac{2}
{3}\, \mathrm{cm}\)}{}{}}
\LANGsk{
\Question{
Na vlákno zavesíme teleso o hmotnosti
\(20\, \mathrm{N}\) 
(\(F_{g}\)) a takto
vzniknuté kyvadlo vychýlime. Vychýlením kyvadla sa zvýši poloha telesa nad
podložkou o \(10\, \mathrm{cm}\) 
(\(h\)). V tejto polohe je
vlákno napínané silou \(12\, \mathrm{N}\) 
(\(F_{1}\)). Určte dĺžku
vlákna (\(l\)).
(Nápoveda: Hmotnosť zaveseného telesa sa rozloží na sily
\(F_{1}\) a
\(F_{2}\) (zložky tiažovej
sily). Sila \(F_{1}\) spôsobuje
napínanie vlákna a \(F_{2}\) 
vracia kyvadlo do zvislej polohy. Rozklad síl sa prevádza pomocou tzv. rovnobežníka
síl.)}
{\(25\, \mathrm{cm}\)}{\(25\, \mathrm{m}\)}{\(6\, \mathrm{cm}\)}{\(16\frac{2}
{3}\, \mathrm{cm}\)}{}{}}
\LANGes{
\Question{
Tenemos un péndulo constituido por una cuerda de longitud \(l\) y un cuerpo que se desplaza desde su posición de equilibrio. La fuerza de gravedad que actúa sobre el cuerpo es \(F_{g} = 20\, \mathrm{N}\). El cuerpo está un \(h = 10\, \mathrm{cm}\) más alto en la posición desplazada (comparando con la posición de equilibrio). La tensión de la cuerda en la posición desplazada es \(F_{1} = 12\, \mathrm{N}\). Halla la longitud de la cuerda \(l\). Sugerencia: Usando un paralelogramo, la fuerza de gravedad sobre el cuerpo puede descomponerse en una fuerza \(F_{1}\) en la dirección de la cuerda y \(F_{2}\) en la dirección perpendicular.}
{\(25\, \mathrm{cm}\)}{\(25\, \mathrm{m}\)}{\(6\, \mathrm{cm}\)}{\(16\frac{2}
{3}\, \mathrm{cm}\)}{}{}}
\End

\Data\ID{9000150504}
\Updated{2020-07-15 03:07:37}
\Level{C}
\SubArea{Triangles}
\IMG{001505_04-figure0.pdf}
\LANGen{
\Question{
The object \(y\) is projected
using a lens with foci at \(F\) 
and \(F'\).
The focal length of the lens (the distance from the focus to the lens) 
\(f = 20\, \mathrm{cm}\). The distance from the
object \(y\) to the lens 
\(a = 60\, \mathrm{cm}\). Find the distance from
the lens to the image \(y'\).}
{\(30\, \mathrm{cm}\)}{\(600\, \mathrm{cm}\)}{\(\frac{20}
 {3} \, \mathrm{cm}\)}{\(25\, \mathrm{cm}\)}{}{}}
\LANGpl{
\Question{
Obiekt \(y\) jest rzutowany za pomocą soczewki skupiającej z ogniskami w \(F\) 
i \(F'\).  Ogniskowa soczewki (odległość środka soczewki od ogniska)  \(f = 20\, \mathrm{cm}\). Odległość od obiektu \(y\)  do soczewki \(a = 60\, \mathrm{cm}\). Znajdź odległość od soczewki do obrazu \(y'\).}
{\(30\, \mathrm{cm}\)}{\(600\, \mathrm{cm}\)}{\(\frac{20}
 {3} \, \mathrm{cm}\)}{\(25\, \mathrm{cm}\)}{}{}}
\LANGcs{
\Question{
Na obrázku je zakresleno zobrazení předmětu
\(y\) pomocí tenké
spojné čočky. Body \(F\) 
a \(F'\) jsou
tzv. ohniska čočky. Vzdálenost ohniska od čočky je tzv. ohnisková vzdálenost
\(f\). Předmět umístíme
ve vzdálenosti \(a = 60\, \mathrm{cm}\) od čočky
s ohniskovou vzdálenosti \(f = 20\, \mathrm{cm}\).
Určete v jaké vzdálenosti \(a'\) 
od čočky se vytvoří obraz \(y'\).}
{\(30\, \mathrm{cm}\)}{\(600\, \mathrm{cm}\)}{\(\frac{20}
 {3} \, \mathrm{cm}\)}{\(25\, \mathrm{cm}\)}{}{}}
\LANGsk{
\Question{
Na obrázku je zakreslené zobrazenie predmetu
\(y\) pomocou tenkej
spojenej šošovky. Body \(F\) 
a \(F'\) sú
tzv. ohniská šošovky. Vzdialenosť ohniska od šošovky je tzv. ohnisková vzdialenosť
\(f\). Predmet umiestníme
vo vzdialenosti \(a = 60\, \mathrm{cm}\) od šošovky
s ohniskovou vzdialenosťou \(f = 20\, \mathrm{cm}\).
Určte v akej vzdialenosti \(a'\) 
od šošovky sa vytvorí obraz \(y'\).}
{\(30\, \mathrm{cm}\)}{\(600\, \mathrm{cm}\)}{\(\frac{20}
 {3} \, \mathrm{cm}\)}{\(25\, \mathrm{cm}\)}{}{}}
\LANGes{
\Question{
El objeto \(y\) se proyecta
usando una lente con focos \(F\) 
y \(F'\).
La distancia focal de la lente es (la distancia desde los puntos focales hacia la lente) 
\(f = 20\, \mathrm{cm}\). La distancia desde el objeto \(y\) a la lente es 
\(a = 60\, \mathrm{cm}\). Halla la distancia desde la lente hacia la imagen virtual \(y'\).}
{\(30\, \mathrm{cm}\)}{\(600\, \mathrm{cm}\)}{\(\frac{20}
 {3} \, \mathrm{cm}\)}{\(25\, \mathrm{cm}\)}{}{}}
\End

\Data\ID{9000150505}
\Updated{2022-04-13 09:04:11}
\Level{C}
\SubArea{Triangles}
\IMG{001505_05-figure0.pdf}
\LANGen{
\Question{
The iron support has the shape of the right triangle
\(ABC\) with the side
\(AB\) of the length
\(30\, \mathrm{cm}\) and the hypotenuse
\(AC\) of the length
\(50\, \mathrm{cm}\) (see the picture). The
maximal allowed force \(F_{1}\) 
on \(AB\) is
\(270\, \mathrm{N}\). Find the maximal
force \(G\) allowed
at the point \(A\). Hint: The load \(G\) 
at the point \(A\) 
can be decomposed to the direction of the hypotenuse and the other side of the triangle as shown in
the picture.}
{\(360\, \mathrm{N}\)}{\(450\, \mathrm{N}\)}{\(540\, \mathrm{N}\)}{\(162\, \mathrm{N}\)}{}{}}
\LANGpl{
\Question{
Żelazne podparcie ma kształt trójkąta prostokątnego \(ABC\) o boku \(AB\) długości 
\(30\, \mathrm{cm}\) i przeciwprostokątnej \(AC\) o długości \(50\, \mathrm{cm}\) (spójrz rysunek). Maksymalna dozwolona siła \(F_{1}\) 
działająca na  \(AB\) wynosi 
\(270\, \mathrm{N}\). Znajdź maksymalne dopuszczalne obciążenie \(G\) dozwolonej 
w punkcie  \(A\). Wskazówka: Obciążenie \(G\) 
w punkcie \(A\) może zostać rozłożone na kierunek przeciwprostokątnej i na drugi bok trójkąta, jak pokazano na rysunku.}
{\(360\, \mathrm{N}\)}{\(450\, \mathrm{N}\)}{\(540\, \mathrm{N}\)}{\(162\, \mathrm{N}\)}{}{}}
\LANGcs{
\Question{
Nosník má tvar pravoúhlého trojúhelníku (viz obrázek) s odvěsnou
\(AB\) o délce
\(30\, \mathrm{cm}\) a přeponou
\(AC\) o délce
\(50\, \mathrm{cm}\). Jakou
maximální tíhu \(G\) 
může mít břemeno zavěšené v bodě
\(A\), jestliže maximální
povolená tahová síla \(F_{1}\) 
na trám \(AB\) 
je \(270\, \mathrm{N}\)?
(Nápověda: Tíha zavěšeného tělesa se rozloží na dvě složky. Síla
\(F_{1}\) má charakter tahové
síly na část nosníku \(AB\),
složka \(F_{2}\) 
má charakter tlakové síly na část nosníku
\(AC\) - viz
obrázek.)}
{\(360\, \mathrm{N}\)}{\(450\, \mathrm{N}\)}{\(540\, \mathrm{N}\)}{\(162\, \mathrm{N}\)}{}{}}
\LANGsk{
\Question{
Nosník má tvar pravouhlého trojuholníka (viď obrázok) s odvesnou
\(AB\) o dĺžke
\(30\, \mathrm{cm}\) a preponou
\(AC\) o dĺžke
\(50\, \mathrm{cm}\). Akú
maximálnu tiaž \(G\) 
môže mať bremeno zavesené v bode
\(A\), ak maximálna
povolená ťahová sila \(F_{1}\) 
na trám \(AB\) 
je \(270\, \mathrm{N}\)?
(Nápoveda: Tiaž zaveseného telesa sa rozloží na dve zložky. Sila
\(F_{1}\) má charakter ťahovej
sily na časť nosníka \(AB\),
zložka \(F_{2}\) 
má charakter tlakovej sily na časť nosníka
\(AC\) - viď
obrázok.)}
{\(360\, \mathrm{N}\)}{\(450\, \mathrm{N}\)}{\(540\, \mathrm{N}\)}{\(162\, \mathrm{N}\)}{}{}}
\LANGes{
\Question{
Un soporte de hierro tiene una forma de triángulo rectángulo
\(ABC\) con el cateto
\(AB\) que mide
\(30\, \mathrm{cm}\) y la hipotenusa
\(AC\) que mide
\(50\, \mathrm{cm}\) (mira la imagen). La fuerza máxima permitida \(F_{1}\) 
sobre \(AB\) es
\(270\, \mathrm{N}\). Halla la fuerza máxima permitida \(G\) sobre el punto \(A\). Sugerencia: La carga \(G\) 
sobre el punto \(A\) puede descomponerse en la dirección de la hipotenusa y en el otro lado del triángulo como se muestra en la imagen.}
{\(360\, \mathrm{N}\)}{\(450\, \mathrm{N}\)}{\(540\, \mathrm{N}\)}{\(162\, \mathrm{N}\)}{}{}}
\End
